\documentclass[11pt,a4paper]{article}
\usepackage[a4paper,margin=0.95in]{geometry}
\usepackage[T1]{fontenc}
\usepackage{lmodern,microtype}
\usepackage{amsmath,amssymb,amsthm,mathtools}
\usepackage{booktabs,array,tabularx,longtable,multirow,float}
\usepackage{enumitem,xcolor,url,listings,needspace,multicol}
\usepackage[hidelinks]{hyperref}
\usepackage[nameinlink,noabbrev,capitalise]{cleveref}
\usepackage{fancyhdr}

\definecolor{deepblue}{RGB}{28,69,125}
\hypersetup{colorlinks=true,linkcolor=deepblue,citecolor=deepblue,urlcolor=deepblue,
 pdftitle={Structural Cryptanalysis of QR-UOV: Projective Direct Forgery, Equivalent-Key Recovery, and Limits of Vinegar Retuning},
 pdfauthor={Justin Thaler, Kai Zhang, Scott Kominers, Quang Dao},
 pdfsubject={Projective direct forgery, verified equivalent-key recovery, and parameter implications for QR-UOV},
 pdfkeywords={QR-UOV,UOV,projective direct forgery,equivalent-key recovery,Chevalley-Warning,nonsingular roots,Karatsuba,symbolic homotopy,exact verification}}

\theoremstyle{definition}
\newtheorem{theorem}{Theorem}[section]
\newtheorem{lemma}[theorem]{Lemma}
\newtheorem{proposition}[theorem]{Proposition}
\newtheorem{corollary}[theorem]{Corollary}
\newtheorem{definition}[theorem]{Definition}
\theoremstyle{remark}
\newtheorem{remark}[theorem]{Remark}

\newcommand{\F}{\mathbb F}
\newcommand{\PP}{\mathbb P}
\newcommand{\cO}{\mathcal O}
\newcommand{\cB}{\mathcal B}
\newcommand{\cX}{\mathcal X}
\newcommand{\cR}{\mathcal R}
\newcommand{\cU}{\mathcal U}
\newcommand{\cV}{\mathcal V}
\newcommand{\cW}{\mathcal W}
\newcommand{\cG}{\mathcal G}
\newcommand{\cE}{\mathcal E}
\newcommand{\Span}{\operatorname{span}}
\newcommand{\rank}{\operatorname{rank}}
\newcommand{\codim}{\operatorname{codim}}
\newcommand{\cdeg}{\operatorname{cdeg}}
\newcommand{\Sym}{\operatorname{Sym}}
\newcommand{\disc}{\mathrm{disc}}
\newcommand{\bad}{\mathrm{bad}}
\newcommand{\oil}{\mathrm{oil}}
\newcommand{\off}{\mathrm{off}}
\newcommand{\spanbad}{\mathrm{span}}
\newcommand{\QRUOV}{\textsf{QR-UOV}}
\newcommand{\Solve}{\operatorname{Solve}}
\newcommand{\Adv}{\operatorname{Adv}}
\newcommand{\trans}{\mathsf T}
\newcommand{\softO}{\widetilde O}
\newcommand{\Mmul}{\mathsf M}
\newcommand{\Tr}{\operatorname{Tr}}
\newcommand{\Gmul}{\mathsf G_{\rm mul}}
\newcommand{\proofstep}[1]{\par\smallskip\noindent\textbf{#1}\enspace}
\newcommand{\ord}{\operatorname{ord}}
\newcommand{\M}{\mathsf M}

\setlist[itemize]{leftmargin=*,itemsep=0.10em,topsep=0.20em}
\setlist[enumerate]{leftmargin=*,itemsep=0.12em,topsep=0.20em}
\setlength{\parindent}{1.2em}
\setlength{\parskip}{0.20em}
\linespread{1.04}
\setlength{\emergencystretch}{2em}
\raggedbottom
\pagestyle{fancy}
\fancyhf{}
\fancyfoot[C]{\thepage}
\renewcommand{\headrulewidth}{0pt}

\title{\bfseries Structural Cryptanalysis of \QRUOV{}:\ Projective Direct Forgery, Local-Separator Descent, and Equivalent-Key Recovery}
\author{Justin Thaler\thanks{Georgetown University and a16z crypto research.}
\and Kai Zhang\thanks{MIT.}
\and Scott Kominers\thanks{Harvard University and a16z crypto research.}
\and Quang Dao\thanks{Carnegie Mellon University and LayerZero Labs.}}
\date{31 July 2026}

\begin{document}
\maketitle

\begin{abstract}
\QRUOV{} compresses Unbalanced Oil and Vinegar (UOV) by representing variables
in a degree-three quotient ring.  We give two complementary structural attacks
and a local-separator refinement that, under an explicit symbolic-homotopy gate
model, yields a candidate Level-I break.

The direct attack sends a chosen target to $(0,\ldots,0,1)$, constructs an
$(m-3)$-dimensional common isotropic space for three zero-target quadrics, and
quotients the remaining homogeneous system by scaling.  The resulting square
projective cores have dimensions $50$, $74$, and $101$.  The equivalent-key
attack recovers the hidden oil graph from one regular directional root.

The main new refinement splits the homotopy before lifting and replaces a
global primitive-element requirement by local separation of one eligible
simple target branch.  A generic characteristic polynomial and its coefficient
polars remain valid even when the chosen linear form collides on other branches.
On a simple target factor, a single Frobenius sketch is exact.  In the inherited
Boolean-gate model, the resulting costs are $140.383995$, $191.027133$, and
$247.125439$ bits, respectively $2.616$, $15.973$, and $24.875$ bits below the
NIST category targets.  Using the parent paper's weaker root-probability bound
still gives $140.531461$ bits at Level I.

We also implement the complete local-separator pipeline on reduced QR-UOV
parameters.  From only the public key, message, and salt, the program constructs
the residual system, lifts and recombines all branches, performs rational
specialization and exact Frobenius descent, reconstructs one point, and outputs
a forgery accepted by the unchanged public verifier.  Five deterministic
instances and a supplementary larger-field instance all succeed and replay
from public transcripts.  This closes the end-to-end composability objection at
toy scale, but it does not validate the full-parameter time or memory
extrapolation.  The Level-I statement remains conditional on the stated
characteristic-polynomial degree transfer, the fast online-product schedule,
the public-random expansion model, and astronomical symbolic memory.
\end{abstract}

\section{Introduction}\label{sec:intro}

\subsection*{Background and attack goals}
A multivariate signature scheme publishes a quadratic map
\[
 P:\F_q^n\longrightarrow\F_q^m
\]
and signs a target $y$ by finding $x$ with $P(x)=y$.  UOV makes this inversion
easy with a secret split between \emph{vinegar} and \emph{oil} variables.  No
central equation contains an oil--oil product, so fixing the vinegar variables
leaves a linear system in the oil variables~\cite{KPG1999}.

An attacker does not need the submitters' seed.  Two endpoints are enough:
\begin{itemize}
\item a \emph{direct forgery} finds one preimage of the target derived from a
message and salt; and
\item an \emph{equivalent-key attack} finds any hidden decomposition that can
invert every target accepted by the unchanged verifier.
\end{itemize}
This paper gives one attack of each kind.  They work in different coordinate
systems and expose different parts of QR-UOV's structure.  Their combination
therefore says more about parameter repair than either attack alone.

QR-UOV uses a quotient ring of degree $\ell=3$.  Its recommended
$q=127$ profiles have base-field dimensions
\[
 (n,m)=(210,54),\qquad(306,78),\qquad(411,105),
\]
and extension-field dimensions
\[
 (V,O,m)=(52,18,54),\qquad(76,26,78),\qquad(102,35,105)
\]
over $K=\F_{127^3}$.  NIST advanced QR-UOV to the third round of its
Additional Digital Signatures process in May 2026; third-round teams may submit
updated specifications and implementations~\cite{NISTRound3,NISTIR8610}.

\subsection*{Attack I: remove three equations, then quotient by scaling}
Let $y\ne0$ be a verification target.  A public output transformation sends it
to
\[
 (0,\ldots,0,1).
\]
Call the first three transformed quadrics $Q_1,Q_2,Q_3$.  Choose a random
seven-dimensional seed subspace $S_0$.  Chevalley--Warning guarantees a
nonzero common zero $w\in S_0$, found by enumerating $\PP^6(\F_q)$.  For every
fixed hidden oil space, $S_0$ intersects it nontrivially with probability at
most
\[
 \frac{q^m-1}{q-1}\frac{q^7-1}{q^n-1},
\]
so $w$ is off oil with overwhelming probability over public attack coins.
This fixed-key statement avoids converting an average density calculation into
an unjustified per-key rejection-sampling claim.

Starting from $w$, the rest of the subspace construction is deterministic.  If
$W$ is a $j$-dimensional subspace on which all three quadrics vanish, intersect
their three polar-orthogonal spaces.  Modulo $W$, the result has dimension at
least $n-4j$.  Whenever it has dimension at least seven, restrict to any
seven-dimensional subspace and use Chevalley--Warning again to add one vector.
The construction reaches dimension $r$ whenever
\[
 n\ge4r+3.
\]
For $r=m-3$, the official sets have slacks $3,3,0$.

Full target normalization makes all transformed target coordinates except the
last one zero.  On the constructed $r$-space, the next $r-1=m-4$ homogeneous
quadrics therefore define a square system in $\PP^{r-1}$.  A random affine
chart turns it into $r-1$ quadrics in $r-1$ variables.  The last transformed
quadric is not included in the solver core: it only tests whether the
projective point can be scaled to the nonzero target, and then supplies that
scale.  Thus the actual geometric solver dimensions are
\[
 50,\qquad74,\qquad101.
\]
The affine dimensions $51,75,102$ are the $\alpha=3$ Hashimoto line already
used in QR-UOV's direct-attack analysis~\cite{Hashimoto2023,QRUOVSpec}; the
one-variable projective quotient is an additional exact reduction.

For a point outside the oil space, one public-random pull-back form has uniform
evaluation and a full affine family of first derivatives.  Evaluation at two
off-oil points is independent unless the points are base-field scalar
multiples.  A second-moment calculation proves that an eligible projective
root with the required nonsingularity exists with probability at least
\[
 0.326336998767509\ldots.
\]
A random nonzero chart contains any fixed projective root with probability
strictly greater than $1-1/q$.  The root-probability theorem uses an idealized
model in which the secret graph seed is independent of the public central-form
seed.  The official construction has this property in the ideal-XOF and
ideal-cipher models except on a seed collision; \cref{app:seeded-transfer}
charges that event explicitly.  Every reconstructed vector is evaluated under
the original public key, so failure can omit a signature but cannot create a
false one.

\subsection*{Attack II: one directional root gives an equivalent key}
In normalized extension-field coordinates, the hidden oil space is the graph
\[
 \cO=\{(-Gc,c):c\in K^O\}.
\]
Fixing a public direction $c\in K^O$ yields $m$ public quadrics in only $V$
variables, with planted solution $Gc$.  If the planted derivative has rank
$V$, one recovered root determines every column of $G$ through public linear
systems.  The graph is checked against every public identity, pulled back to
the official base-field coordinates, and accompanied by an invertible signing
matrix.  An accepted output is therefore an exact universal-forgery key.

The solver adapts the nonsingular-solutions homotopy of Safey El Din and Schost
to characteristic $127$~\cite{SafeySchost2018}.  An explicit Vandermonde
product start and a full random separator replace the two short
integer-indexed families responsible for the published characteristic
restriction.  A flattened coefficient algebra and sparse irreducible working
fields improve the finite-size ledger.  The audit also corrects a previous
undercount: rational reconstruction requires precision $2C_0$, not $C_0$.

\subsection*{The fixed-output security scissors}
The attacks react oppositely to the vinegar count $V$.  For quotient degree
$\ell$, the base-field dimension is $n=\ell V+m$.  Target-orthogonal slicing
of $s$ equations is possible whenever
\[
 \ell V\ge s(m-s).
\]
After projectivization, that reduction has square solver core $m-s-1$ rather
than $m-s$.  Taking $s=\ell$ gives a direct core $m-\ell-1$ whenever
$V\ge m-\ell$.  If $V\le m-\ell-1$, the directional core is already at most
$V\le m-\ell-1$.  Since there is no integer between these two cases,
\[
 r_{\rm attack}(V)\le m-\ell-1
\]
for every vinegar count and every nonzero target.  For QR-UOV, the ceiling is
$m-4$.  This is an algebraic-core statement, not a claim that time depends
only on dimension.  Its design implication is exact: increasing $V$ while
holding $m$ fixed cannot raise the combined bottleneck beyond the crossing,
while decreasing $V$ makes the directional attack cheaper.  A repair must
also increase $m$, or alter the construction so that one reduction no longer
applies.

\subsection*{Concrete conclusions}
Using the QR-UOV specification's broad Boolean-gate unit, explicit recursive
Karatsuba circuits for the auxiliary fields, the corrected lifting precision,
local separation, and complete success normalization, the best rows are
\[
 140.384,\qquad191.027,\qquad247.125
\]
bits at Levels I, III, and V.  Therefore all three parameter sets fall below
their category targets in the stated model, by $2.616$, $15.973$, and $24.875$
bits respectively.  The Level-I conclusion is a candidate model-based break,
not a practical computation: its finite margin depends on the explicit
characteristic-polynomial degree transfer and fast online-product envelope.
The reduced-parameter implementation nevertheless executes the entire attack
pipeline and produces publicly accepted forgeries, so the construction is no
longer supported only by disconnected component tests.
The symbolic states are astronomical.  We report time, memory, exact
reductions, and distributional assumptions separately rather than presenting
the rows as implementation benchmarks.

\subsection*{What is exact, and what is modeled}
\begin{table}[t]
\centering
\small
\caption{Status of the main claims.}
\label{tab:claim-status-final}
\begin{tabularx}{\textwidth}{@{}p{0.31\textwidth}X@{}}
\toprule
Claim & Status \\
\midrule
Projective direct core $m-4$ & Exact for every nonzero target once a seed root is supplied; Chevalley--Warning guarantees such a root in every seven-dimensional seed, and homogeneity quotients the residual by scaling. \\
Seed root is off oil & Fixed-key probability bound over public attack coins; no random-key assumption. \\
Common-isotropic extension & Deterministic and finite after the seed, using linear algebra and exhaustive search in $\PP^6(\F_q)$. \\
Eligible projective-root probability & Exact lower bound in the independent-graph public-random degree-three model. \\
One regular direction determines an equivalent key & Exact for every expanded key satisfying the public rank and signing checks. \\
Returned signature or key is valid & Exact; every candidate is verified against the unchanged public map. \\
Finite cost rows & Reproducible symbolic-algorithm diagnostics, not measured circuits. \\
Full-parameter memory & Prohibitive for the direct representations analyzed here. \\
Fixed SHAKE/AES expansion & Key-by-key verification is exact; density statements transfer only in the stated ideal-primitive models, with explicit buffer and seed-collision terms. \\
\bottomrule
\end{tabularx}
\end{table}

\subsection*{Organization}
\Cref{sec:related} positions the two attacks.  \Cref{sec:model} explains the
official pull-back, the public-random models, and exact signing certificates.
\Cref{sec:direct} proves target-orthogonal slicing and the security scissors.
\Cref{sec:recovery,sec:local-geometry} give directional graph recovery.
\Cref{sec:finite-char-solver,sec:flat-lifting,sec:sparse-fields} establish the
characteristic-$127$ solver and cost refinements.  \Cref{sec:security} combines
the attacks and gives parameter consequences.  Independent solver routes and
seeded-expansion arguments are retained in the appendices.

\section{Related work}\label{sec:related}

\paragraph{Direct attacks on underdetermined MQ.}
Thomae--Wolf, Furue et al., and Hashimoto reduce underdetermined quadratic
systems before applying an MQ solver~\cite{ThomaeWolf2012,FurueUnderdetermined2021,Hashimoto2023}.
Hashimoto's line is the benchmark used in the QR-UOV specification.  The
QR-UOV team's January 2026 response reports classical Hashimoto estimates of
$158$, $217$, and $285$ bits for the recommended sets~\cite{QRUOVResponse2026}.
The affine residual dimensions $51,75,102$ are therefore prior art.  Our
contribution is a self-contained homogeneous construction, a fixed-key
seed-disjointness bound, a QR-UOV-specific nonsingular-root theorem, and the
observation that full target normalization lets the solver quotient by radial
scaling, reducing the square geometric cores to $50,74,101$.  We then cost
those projective cores with a separate finite-characteristic solver.  The recent ``Just Guess''
algorithm improves the quantum analysis of underdetermined MQ but does not
improve the team's reported classical QR-UOV estimates~\cite{JustGuess,QRUOVResponse2026}.

\paragraph{Equivalent-key recovery for UOV.}
P\'ebereau showed that one suitable oil vector can suffice for polynomial-time
UOV equivalent-key recovery~\cite{Pebereau}.  We do not claim that endpoint as
new.  The QR-UOV-specific contribution is the public mechanism for finding the
needed information: a direction exposes a planted root $Gc$, full derivative
rank makes it isolated, and cross-direction graph identities recover all of
$G$.  Exact pull-back and signing checks then certify a key-only universal
forgery.

Beullens' intersection attack and later work recover hidden UOV geometry
through structured linear algebra~\cite{Beullens}.  Rectangular-MinRank
estimates were sharpened for QR-UOV and related schemes~\cite{FurueIkematsu,Suzuki}.
Truncated-polynomial-ring methods improve some UOV parameters but do not report
an improvement for the QR-UOV sets studied here~\cite{FurueIkematsu2026}.
Odd-characteristic symmetric-algebra attacks apply to QR-UOV but, at
characteristic $127$, do not improve the existing QR-UOV estimates~\cite{JinEtAl2026}.

\paragraph{Polynomial-system solving.}
The arbitrary-core engine is the nonsingular-solutions homotopy of Safey El
Din and Schost~\cite{SafeySchost2018}, building on symbolic homotopy and
geometric-resolution methods~\cite{Schost2003,GiustiLecerfSalvy2001}.  We
replace the start-system and separator choices responsible for its published
characteristic restriction.  The appendices retain independent checks using
the near-quadratic solver of van der Hoeven--Lecerf and a locally closed
Kronecker algorithm~\cite{vdHL,Gimenez2026}.

\paragraph{Chevalley--Warning.}
The direct reduction uses the classical fact that, when the sum of the degrees
is smaller than the number of variables, the number of common zeros over a
finite field is divisible by the characteristic~\cite{Chevalley1935,Warning1935}.
Applied to three quadrics in seven variables, it guarantees a nonzero extension
vector.  Exhaustive projective search makes the theorem constructive, so no
unpriced existence oracle is hidden in the preprocessing.

\section{Expanded-key model and exact signing certificates}\label{sec:model}

This section connects the algebra used by the attack to the official verification map.  There are two coordinate systems:
\begin{itemize}
\item the official base-field system over $\F_{127}$, in which signatures are verified; and
\item a smaller extension-field system over $K=\F_{127^3}$, in which the hidden oil space and the attack are easier to describe.
\end{itemize}
The specification gives an exact, invertible pull-back between them.  We first recover and verify the hidden structure over $K$, then pull it back and check the resulting base-field factorization against the unchanged public key.  One invertible oil system then certifies that the recovered decomposition can sign every target.

The Round-2 $q=127$ parameter sets use quotient-ring degree $\ell=3$.  The specification has $v=\ell V$ base-field vinegar variables and $m=\ell O$ base-field oil variables and output coordinates.  After quotient-ring expansion, the attacker works with the same $m$ outputs as quadratic forms in only $N=V+O$ variables over
\[
K=\F_{q^\ell}=\F_{127^3},\qquad Q=|K|=127^3=2{,}048{,}383,
\qquad \log_2 Q=20.966054\ldots.
\]
Thus $m$ denotes both the number of public output coordinates and the base-field oil dimension, whereas $O=m/\ell$ is the oil dimension over $K$.

Let $\Phi_f:K\to\F_q^{\ell\times\ell}$ be the regular representation from the specification, and let $W_f\in\F_q^{\ell\times\ell}$ be its invertible symmetrizer.  Write $\phi$ for the blockwise inverse of $\Phi_f$, and write $W_f^{(a)}$ for the block-diagonal matrix with $a$ copies of $W_f$.

The attack uses the expanded extension-field tuple
\[
P=(p_1,\ldots,p_m),\qquad p_i(x)=x^{\mathsf T}P_ix,
\qquad x\in K^{N},
\]
where the symmetric matrices $P_i$ share a hidden totally isotropic $O$-space $\cO$:
\begin{equation}\label{eq:isotropy}
u^{\mathsf T}P_iw=0\qquad(u,w\in\cO,\ 1\le i\le m).
\end{equation}
The Round-2 sets analyzed in this paper are
\begin{center}
\begin{tabular}{@{}lrrrrrr@{}}
\toprule
Level & $q$ & $\ell$ & $v$ & $m$ & $V=v/\ell$ & $O=m/\ell$\\
\midrule
I&127&3&156&54&52&18\\
III&127&3&228&78&76&26\\
V&127&3&306&105&102&35\\
\bottomrule
\end{tabular}
\end{center}

\begin{center}
\small
\setlength{\tabcolsep}{5pt}
\begin{tabularx}{0.94\textwidth}{@{}>{\bfseries}p{0.16\textwidth}X@{}}
\toprule
Symbol & Meaning \\
\midrule
$N=V+O$ & Number of extension-field variables. \\
$e=m-V$ & Number of public equations not used in the square $V$-equation core. \\
$D=2^V$ & B\'ezout degree of a smooth projective core. \\
$L_s=\F_{Q^s}$ & Solver working field; the same field supplies slicing, recombination, and rational-root extraction. \\
\bottomrule
\end{tabularx}
\end{center}

All concrete theorems below concern these three sets, for which $m<q$.  The directional and projective equivalent-key analyses use \cref{def:ideal} and allow arbitrary correlation between the secret graph and the symmetric public blocks.  The direct-forgery root-counting theorem uses the stronger independent-graph specialization of \cref{def:direct-ideal}.  The fastest projective route has cost $\softO(2^V2^{(1+\varepsilon)(V-1)}3\log_2 127)$ for every fixed rational $\varepsilon>0$, while the graph-direction route uses the explicit locally closed finite-field theorem of Gim\'enez et al.  \Cref{app:seeded-transfer} gives separate ideal-XOF and ideal-cipher transfers for the two models and isolates the remaining fixed-primitive gap.

\subsection{Public-random expansion model and relation to seeded expansion}

\begin{definition}[Public-random expansion with an arbitrarily correlated secret graph]\label{def:ideal}
First sample
\[
(A_1,\ldots,A_m)\leftarrow\Sym_V(K)^m
\]
from the product-uniform distribution.  Conditional on the resulting tuple $A=(A_i)_i$, sample the secret graph $G\in K^{V\times O}$ according to an arbitrary conditional distribution.  Conditional on $(A,G)$, sample
\[
(E_1,\ldots,E_m)\leftarrow(K^{V\times O})^m
\]
from the product-uniform distribution, set each lower-right public block by the key-generation identity, and expose the resulting expanded public tuple.  Values are sampled once and reused consistently.  No independence between $G$ and $A$ is assumed.
\end{definition}

\begin{definition}[Independent-graph specialization]\label{def:direct-ideal}
The \emph{independent-graph public-random model} is \cref{def:ideal} with the
additional requirement that the marginal law of $G$ be independent of $A$.
Equivalently, after centralization the product-uniform coefficient pairs
$(A_i,B_i)_{1\le i\le m}$ are independent of $G$.  The marginal law of $G$
may otherwise be arbitrary.
\end{definition}

\begin{proposition}[Exact central-block factorization under graph correlation]\label{prop:ideal-factorization}
In the block notation of \eqref{eq:central}--\eqref{eq:graphid}, define
\[
B_i:=E_i-A_iG.
\]
Conditional on every fixed pair $(A,G)$, the tuple $(B_1,\ldots,B_m)$ is product-uniform and its conditional law does not depend on $(A,G)$.  Consequently, $B$ is independent of $(A,G)$, and the pairs $(A_i,B_i)_{1\le i\le m}$ are independent and uniform on $\Sym_V(K)\times K^{V\times O}$.  The graph $G$ may nevertheless remain arbitrarily correlated with $A$.  The lower-right public blocks $C_i$ are determined by \eqref{eq:graphid}.
\end{proposition}

\begin{proof}
Fix $(A,G)=(a,g)$.  For each $i$, the map
\[
E_i\longmapsto B_i=E_i-a_ig
\]
is a translation of the additive space $K^{V\times O}$ and hence a bijection.  It sends the product-uniform law of $E$ to the product-uniform law of $B$, independently of $(a,g)$.  Thus $B$ is independent of $(A,G)$.  Since $A$ itself is product-uniform, $(A_i,B_i)_i$ has the claimed product law.  The formula for $C_i$ is the key-generation block identity \eqref{eq:graphid}.
\end{proof}

The directional and projective bad-key calculations depend on the central pairs $(A_i,B_i)$, uniform public slices, and fresh attack coins.  A secret coordinate change depending on $G$ maps a uniform public slice to a uniform central slice for every fixed key, so those calculations tolerate arbitrary $G$--$A$ correlation.  The direct-forgery reduction itself is also key by key, but its average-case residual-root theorem conditions on an off-oil point and then uses independence of $G$ from the central coefficient pairs; it is therefore stated under \cref{def:direct-ideal}.  \Cref{app:seeded-transfer} proves separate ideal-primitive transfers, including the public/secret seed-collision term needed for the independent-graph model.  All returned keys and signatures are verified against the unchanged public map, independently of any distributional model.

\subsection{Central and public graph coordinates}

In secret central coordinates $K^V\oplus K^O$,
\begin{equation}\label{eq:central}
F_i=\begin{pmatrix}A_i&B_i\\B_i^{\trans}&0\end{pmatrix},
\qquad F_i(v,o)=v^{\trans}A_iv+2v^{\trans}B_io,
\end{equation}
with $A_i\in\Sym_V(K)$ and $B_i\in K^{V\times O}$ independent and uniform.  The planted oil space is $\cO_0=\{0\}\oplus K^O$.

In normalized public coordinates the hidden space is a graph
\begin{equation}\label{eq:graph}
\cO=\left\{\binom{-Gc}{c}:c\in K^O\right\}.
\end{equation}
Writing
\[
P_i=\begin{pmatrix}A_i&E_i\\E_i^{\trans}&C_i\end{pmatrix},
\]
total isotropy is equivalent to
\begin{equation}\label{eq:graphid}
C_i=G^{\trans}E_i+E_i^{\trans}G-G^{\trans}A_iG.
\end{equation}

\begin{proposition}[Verified graph gives an exact expanded UOV decomposition]\label{prop:verified-key}
Let $G'\in K^{V\times O}$ satisfy every identity \eqref{eq:graphid}, and put
\[
B_i'=E_i-A_iG',\qquad
F_i'=\begin{pmatrix}A_i&B_i'\\ {B_i'}^{\trans}&0\end{pmatrix},\qquad
S_{G'}=\begin{pmatrix}I_V&G'\\0&I_O\end{pmatrix}.
\]
Then
\[
P_i=S_{G'}^{\trans}F_i'S_{G'}
\qquad(1\le i\le m).
\]
Consequently, the specification's inverse pull-back gives an expanded base-field UOV decomposition of the unchanged verification map.
\end{proposition}

\begin{proof}
Since
\[
S_{G'}^{-1}=\begin{pmatrix}I_V&-G'\\0&I_O\end{pmatrix},
\]
a direct block multiplication gives
\[
S_{G'}^{-\trans}P_iS_{G'}^{-1}
=
\begin{pmatrix}
A_i&E_i-A_iG'\\
E_i^{\trans}-{G'}^{\trans}A_i&
C_i-{G'}^{\trans}E_i-E_i^{\trans}G'+{G'}^{\trans}A_iG'
\end{pmatrix}.
\]
The lower-right block is zero by \eqref{eq:graphid}; the remaining blocks are exactly $F_i'$.  Rearranging proves the displayed factorization.

To make the interface with the specification explicit, let hats denote the corresponding base-field matrices.  The forward pull-back of Section~4.1 of the specification has the form
\[
P_i=\phi\!\left((W_f^{(N)})^{-1}\widehat P_i\right),\qquad
F_i'=\phi\!\left((W_f^{(N)})^{-1}\widehat F_i'\right),\qquad
S_{G'}=\phi(\widehat S_{G'}).
\]
Hence the inverse pull-back sets
\[
\widehat P_i=W_f^{(N)}\phi^{-1}(P_i),\qquad
\widehat F_i'=W_f^{(N)}\phi^{-1}(F_i'),\qquad
\widehat S_{G'}=\phi^{-1}(S_{G'}).
\]
The pull-back identity in the specification converts the verified extension-field factorization into
$\widehat P_i=\widehat S_{G'}^{\trans}\widehat F_i'\widehat S_{G'}$ for the unchanged base-field verification map~\cite[Sec.~4.1]{QRUOVSpec}.  These matrix identities publicly certify an exact expanded UOV decomposition.  Signing uses this base-field pull-back, not the extension-field tuple with only $O=m/3$ oil variables.  The separate test in \cref{prop:signing-witness} certifies that the decomposition can invert every target.
\end{proof}

For an expanded base-field decomposition $\widehat P_i=\widehat S^{\mathsf T}\widehat F_i\widehat S$, write the central variables as $(a,o)\in\F_q^v\oplus\F_q^m$.  Because the oil--oil block of every $\widehat F_i$ is zero, fixing $a$ leaves a linear map in the oil variables.  Denote its matrix by
\begin{equation}\label{eq:signing-matrix}
\mathsf L_{\widehat F}(a)\in\F_q^{m\times m},
\qquad
\bigl(\mathsf L_{\widehat F}(a)o\bigr)_i
:=\widehat F_i(a,o)-\widehat F_i(a,0).
\end{equation}

\begin{proposition}[Public signing witness]\label{prop:signing-witness}
If $\det\mathsf L_{\widehat F}(a)\ne0$, then $(\widehat F,\widehat S,a)$ is a publicly checkable signing certificate.  For every target $y\in\F_q^m$, one computes
\[
o=\mathsf L_{\widehat F}(a)^{-1}\bigl(y-\widehat F(a,0)\bigr),
\qquad
x=\widehat S^{-1}(a,o),
\]
and obtains $\widehat P(x)=y$.
\end{proposition}

\begin{proof}
Direct substitution gives $\widehat F(a,o)=y$, followed by $\widehat P(x)=\widehat F(\widehat Sx)=y$.  These signatures verify under the unchanged public map; matching the official signing distribution is neither needed nor claimed.
\end{proof}

\begin{corollary}[Key-only universal chosen-message forgery]\label{cor:universal-forgery}
Given a publicly verified signing certificate from \cref{prop:signing-witness}, an attacker can forge an accepted \QRUOV{} signature on every chosen message without querying a signing oracle.  Concretely, for a message $M$, choose any allowed salt $r$, compute the official target
\[
 \mu=\operatorname{SHAKE256}(\mathsf{seed}_{\rm pk}\mathbin\|\operatorname{BytesToBits}(M),512),
 \qquad y=\operatorname{Hash}(\mu,r),
\]
invert $y$ using \cref{prop:signing-witness}, and output $(r,x)$.  The official verifier accepts because $\widehat P(x)=\operatorname{Hash}(\mu,r)$.
\end{corollary}

\begin{proof}
The specification's verifier recomputes $\mu$ and $\operatorname{Hash}(\mu,r)$ and accepts exactly when the public map evaluated at the signature vector equals that target~\cite[Secs.~4.2--4.3]{QRUOVSpec}.  \Cref{prop:signing-witness} inverts every target under the unchanged public map.  No distributional claim about the forged signature is needed for verification, and the secret seed is never recovered.
\end{proof}

Fix a nonzero $\xi\in\F_q^\ell$.  For each quotient-ring vinegar block $1\le j\le V$, let $a_j$ place $\xi$ in block $j$ and zero elsewhere, and set
\[
M_j:=\mathsf L_{\widehat F}(a_j)\in\F_q^{m\times m}.
\]
For $0\le r\le m$, write
\begin{equation}\label{eq:rank-probability}
\varrho_r:=\Pr_{M\leftarrow\F_q^{m\times m}}[\rank M=r]
=q^{-m^2}\prod_{i=0}^{r-1}\frac{(q^m-q^i)^2}{q^r-q^i}.
\end{equation}
In particular,
\begin{equation}\label{eq:square-invertible}
 \pi_m:=\varrho_m=\prod_{k=1}^m(1-q^{-k}),
 \qquad \sigma_m:=1-\pi_m.
\end{equation}
For the stronger fallback, define
\begin{equation}\label{eq:pair-failure}
\chi_m:=q^{-(m-1)}+(1-q^{-(m-1)})q^{-1},
\qquad
\zeta_{\rm pair}:=
\varrho_{m-1}q^{-1}\chi_m
+(1-\varrho_m-\varrho_{m-1})(1-\varrho_m).
\end{equation}

\begin{proposition}[Basis-first deterministic signing search]\label{prop:signing-search}
For the planted decomposition in the public-random expansion model, the matrices
$M_1,\ldots,M_V$ are independent and uniform in $\F_q^{m\times m}$.
Consequently:
\begin{enumerate}
\item Testing the basis assignments $a_1,a_2,\ldots$ in order uses fewer than
\[
 \pi_m^{-1}<1.008
\]
basis tests on average, stopping at the first success or after all $V$ tests.  It fails to find an invertible matrix exactly when all $V$ basis matrices are singular, an event of probability
\begin{equation}\label{eq:basis-sign-failure}
 \eta_{\rm basis}:=\sigma_m^V,
\end{equation}
whose base-two logarithms at Levels I, III, and V are
\[
-362.823,\qquad-530.280,\qquad-711.692.
\]

\item Much shorter fixed lists already suffice at the category scale.  Testing the first
\[
 b_{\rm sign}=22,\qquad31,\qquad40
\]
basis assignments at Levels I, III, and V fails with probabilities $\sigma_m^{b_{\rm sign}}$, whose base-two logarithms are
\begin{equation}\label{eq:compact-sign-logs}
 -153.502,\qquad-216.298,\qquad-279.095.
\end{equation}

\item Pair the basis matrices as $(M_1,M_2),(M_3,M_4),\ldots$.  Let $\cE_{\rm sign}$ be the event that every designated pair span contains only singular matrices.  Then
\begin{equation}\label{eq:eta-sign}
 \Pr[\cE_{\rm sign}]
 \le \eta_{\rm sign}:=\zeta_{\rm pair}^{\lfloor V/2\rfloor},
\end{equation}
with base-two logarithms below
\[
-544.820,\qquad-796.276,\qquad-1068.687.
\]
Outside $\cE_{\rm sign}$, a deterministic fallback succeeds after at most
\begin{equation}\label{eq:signing-scan-count}
 V+\lfloor V/2\rfloor(q-1)
\end{equation}
oil-matrix tests: first test every $M_j$; if necessary, for each pair $(A,B)$ test $A+tB$ for all $t\in\F_q^\times$.  The basis tests already cover the two remaining projective points $A$ and $B$.
\end{enumerate}

For any fixed exact candidate, define its oil-matrix determinant polynomial
\[
\Delta(c_1,\ldots,c_V):=\det\!\left(\sum_{j=1}^V c_jM_j\right).
\]
Whenever $\Delta$ is nonzero, a uniform scalar-block vinegar assignment is invertible with probability at least $1-m/q$.  The corresponding expected numbers of assignments are at most
\[
\frac{q}{q-m}=1.740,\qquad2.592,\qquad5.773
\]
at Levels I, III, and V.  This last statement is key by key.  For a recovered candidate not known to be planted, every successful determinant test is still an exact public signing certificate.
\end{proposition}

\begin{proof}
Write $b_{i,j,k}\in K=\F_{q^\ell}$ for entry $(j,k)$ of the planted extension-field block $B_i$.  Under the public-random experiment these elements are independent and uniform.  The $(i,k)$ block of the oil coefficient row produced by $a_j$ is
\[
2\xi^{\mathsf T}W_f\Phi_f(b_{i,j,k})\in\F_q^{1\times\ell}.
\]
The factor $2$ is invertible.  The $\F_q$-linear map
\[
K\longrightarrow\F_q^{1\times\ell},
\qquad b\longmapsto \xi^{\mathsf T}W_f\Phi_f(b)
\]
is injective: if $b\ne0$, then $\Phi_f(b)$ is invertible, while $\xi^{\mathsf T}W_f\ne0$.  Equal dimensions make it an isomorphism.  Thus each $M_j$ is uniform, and distinct $j$ use disjoint entries of the independent blocks $B_i$; hence the $M_j$ are independent.

A basis test succeeds with probability $\pi_m$.  Truncating a geometric random variable at $V$ gives expected test count
\[
 \sum_{j=0}^{V-1}\sigma_m^j
 =\frac{1-\sigma_m^V}{\pi_m}<\pi_m^{-1}.
\]
Independence gives $\sigma_m^b$ for every fixed prefix of length $b$, yielding all displayed basis-list probabilities.

It remains to bound one pair.  Let $A,B$ be independent uniform $m\times m$ matrices, and let $E_{A,B}$ be the event that their span contains no invertible matrix.  If $A$ is invertible, this event is impossible.  If $\rank A=m-1$, invertible left and right changes of basis reduce $A$ to $\operatorname{diag}(I_{m-1},0)$ without changing the distribution of $B$.  Write
\[
B=\begin{pmatrix}C&b\\ c^{\mathsf T}&d\end{pmatrix}.
\]
On $E_{A,B}$, the polynomial $\det(A+tB)$ vanishes for every $t\in\F_q$.  Its degree is at most $m<q$, so it is the zero polynomial.  Its linear coefficient is $d$, and, after $d=0$, its quadratic coefficient is $-c^{\mathsf T}b$.  These conditions have probability at most $q^{-1}\chi_m$.  If $\rank A\le m-2$, then $E_{A,B}$ implies that $B$ is singular, an event of probability $1-\varrho_m$.  Conditioning on the rank of $A$ gives $\Pr[E_{A,B}]\le\zeta_{\rm pair}$.  The designated pairs are independent, proving \eqref{eq:eta-sign}.

The basis matrices and $A+tB$ for $t\in\F_q^\times$ enumerate every projective point on each designated pair line, so the deterministic scan succeeds whenever one pair span contains an invertible matrix.  Finally, $\Delta$ has total degree $m$.  If it is nonzero, Schwartz--Zippel gives success probability at least $1-m/q$ for a uniform coefficient vector.
\end{proof}

\section{Projective target-orthogonal direct forgery}\label{sec:direct}

The directional attack uses QR-UOV's extension-field graph.  This section gives
a second route that works directly with the official base-field verification
map
\[
 \widehat P:\F_q^n\longrightarrow\F_q^m,
 \qquad n=v+m,
 \qquad q=127.
\]
It does not recover a key.  Given a nonzero target $y$, it constructs a public
subspace on which three target equations disappear.  Full output normalization
then leaves $m-4$ homogeneous zero-target equations in projective dimension
$m-4$; one final equation only recovers the affine scale.  The square solver
cores are
\[
 50,\qquad74,\qquad101
\]
for Levels I, III, and V.

The preceding affine dimensions $51,75,102$ coincide with the $\alpha=3$
Hashimoto line already used in the QR-UOV specification's direct-attack
analysis~\cite{Hashimoto2023,QRUOVSpec}.  We do not claim those dimensions as
new.  The further one-variable reduction follows from two facts specific to
the homogeneous verification problem: every target but one can be made zero,
and the remaining zero equations are invariant under nonzero scalar
multiplication.  The other new ingredients are a fixed-key seed bound, a
finite Chevalley--Warning construction, a second-moment lower bound for an
eligible nonsingular projective root, and a characteristic-$127$ solver ledger
independent of WXL semi-regularity estimates.

The algebraic reduction below is key by key.  Its probability calculation uses
the independent-graph public-random model of \cref{def:direct-ideal}: the
hidden graph $G$ is independent of the product-uniform central coefficient
pairs $(A_i,B_i)$.  This independence is needed only for the direct attack's
average-case root analysis; the directional results continue to allow
arbitrary $G$--$A$ correlation.  In the ideal-XOF and ideal-cipher models,
official public and secret expansion have this independence whenever their
independently sampled seeds differ.  \Cref{app:seeded-transfer} adds the
seed-collision and buffer terms explicitly.  No such transfer is claimed for
the fixed SHAKE or AES functions.

\subsection{The pull-back still supplies full first-order randomness}

The official degree-three pull-back produces a restricted family of base-field quadrics, not the set of all UOV quadrics over $\F_q$.  We therefore record the exact local property needed below.

Let $E=K^{V+O}$, viewed as an $n=3(V+O)$ dimensional vector space over $\F_q$, and let $\widehat{\cO}\subset E$ be the hidden $m=3O$ dimensional base-field oil space.  In central coordinates, one public-random expanded form is obtained from
\[
 p(v,o)=v^{\mathsf T}Av+2v^{\mathsf T}Bo
\]
with $A\in\Sym_V(K)$ and $B\in K^{V\times O}$ uniform.  The specification's pull-back evaluates a fixed nonzero $\F_q$-linear functional of $p(v,o)$.  Equivalently, after choosing a nonzero $\gamma\in K$, it has the form
\begin{equation}\label{eq:trace-pullback}
 Q(v,o)=\Tr_{K/\F_q}\!\left(\gamma p(v,o)\right).
\end{equation}
The precise value of $\gamma$ depends on the chosen regular-representation basis and is irrelevant; nondegeneracy of the trace pairing is what matters.

\begin{lemma}[Evaluation and one-jet surjectivity]\label{lem:pullback-one-jet}
Let $x\in E\setminus\widehat{\cO}$, and let $U\subseteq E$ be an $\F_q$-linear subspace containing $x$.  For a uniform pull-back form $Q$ as in \eqref{eq:trace-pullback}:
\begin{enumerate}
\item $Q(x)$ is uniform in $\F_q$;
\item the map
\[
 Q\longmapsto\bigl(Q(x),dQ_x|_U\bigr)
\]
is onto
\[
 \{(a,\lambda)\in\F_q\times U^*: \lambda(x)=2a\};
\]
consequently, conditioned on $Q(x)=a$, the derivative $dQ_x|_U$ is uniform in the displayed affine hyperplane;
\item for $x,z\in E\setminus\widehat{\cO}$, the two evaluation functionals $Q\mapsto Q(x)$ and $Q\mapsto Q(z)$ are proportional over $\F_q$ only if $z=\lambda x$ for some $\lambda\in K$ with $\lambda^2\in\F_q$.  Since $[K:\F_q]=3$ is odd, this forces $\lambda\in\F_q$.
\end{enumerate}
\end{lemma}

\begin{proof}
Write $x=(v_x,o_x)$ in central $K$-coordinates.  Since $x$ is off oil, $v_x\ne0$.  Choose a $K$-linear functional $\varphi:E\to K$ with
$\varphi|_{\widehat{\cO}}=0$ and $\varphi(x)=1$.  The map
\[
 h\longmapsto\bigl(u\mapsto\Tr_{K/\F_q}(\gamma h(u))\bigr)
\]
from $K$-linear maps $E\to K$ to $\F_q$-linear maps $E\to\F_q$ is an isomorphism, by nondegeneracy of the trace pairing and equality of dimensions.

Fix a desired pair $(c,\lambda)$ with $\lambda(x)=2c$.  Choose a $K$-linear $h$ whose trace functional restricts to $\lambda/2$ on $U$, and define
\[
 b_h(u,w)=\varphi(u)h(w)+h(u)\varphi(w)-h(x)\varphi(u)\varphi(w).
\]
This symmetric $K$-bilinear form vanishes on
$\widehat{\cO}\times\widehat{\cO}$ and satisfies $b_h(x,u)=h(u)$.  Hence
\[
 p_h(u):=b_h(u,u),
 \qquad
 Q_h(u):=\Tr_{K/\F_q}(\gamma p_h(u))
\]
is a pull-back UOV form with
\[
 Q_h(x)=\Tr(\gamma h(x))=c,
 \qquad
 d(Q_h)_x(u)=2\Tr(\gamma h(u))=\lambda(u)
 \quad(u\in U).
\]
This proves the first two claims.

For the last claim, compare coefficients in central coordinates.  If the two evaluation functionals are proportional by $\rho\in\F_q^*$, nondegeneracy of the trace pairing gives
\[
 v_zv_z^{\mathsf T}=\rho v_xv_x^{\mathsf T},
 \qquad
 v_zo_z^{\mathsf T}=\rho v_xo_x^{\mathsf T}.
\]
The first rank-one identity gives $v_z=\lambda v_x$ and $\lambda^2=\rho$ for some $\lambda\in K^*$.  The second then gives $o_z=\lambda o_x$, so $z=\lambda x$ and $\lambda^2\in\F_q$.  Finally,
\[
 \gcd\bigl(2(q-1),q^3-1\bigr)=q-1
\]
because $q^2+q+1$ is odd.  Thus $\lambda^{q-1}=1$ and $\lambda\in\F_q$.
\end{proof}

\subsection{Normalize every target coordinate except one}

If $y=0$, then $x=0$ is already a valid preimage, so only the nonzero case
needs analysis.  Let $y\in\F_q^m$ be nonzero.  Choose a public matrix
$T_y\in\operatorname{GL}_m(\F_q)$ whose first $m-1$ rows span $y^\perp$ and
whose last row has inner product one with $y$.  Then
\begin{equation}\label{eq:target-normalization}
 T_yy=(0,\ldots,0,1)^{\mathsf T}.
\end{equation}
Write $Q_1,\ldots,Q_m$ for the corresponding output recombinations of the
public forms.  Since $T_y$ is fixed by $y$ and invertible, product-uniform
public forms remain product-uniform in the public-random expanded-key model.

Only $Q_1,Q_2,Q_3$ are used to construct the attack subspace.  Choose a
uniformly random seven-dimensional subspace $S_0\subset E$, restrict these
three forms to $S_0$, and exhaustively test the points of
$\PP(S_0)\simeq\PP^6(\F_q)$.  Chevalley--Warning guarantees a nonzero common
zero because three quadrics have total degree six in seven variables.  Denote
the first such point by $w$.

The hidden oil space is not known to the attacker, so the algorithm cannot
test whether $w$ is off oil.  The random seed subspace makes this unnecessary:
with overwhelming probability it contains no nonzero oil vector at all.

\begin{proposition}[A random seed misses every fixed oil space]\label{prop:off-oil-seed}
Fix an arbitrary $m$-dimensional subspace $\widehat{\cO}\subset\F_q^n$ and
choose a uniform seven-dimensional subspace $S_0\subset\F_q^n$.  Then
\begin{equation}\label{eq:seed-oil-bound}
 \Pr[S_0\cap\widehat{\cO}\ne\{0\}]
 \le
 \epsilon_{\rm seed}
 :=\frac{q^m-1}{q-1}\,
   \frac{q^7-1}{q^n-1}.
\end{equation}
Consequently, except with probability at most $\epsilon_{\rm seed}$ over
public attack coins, every nonzero common zero found inside $S_0$ is off oil.
For the three official profiles,
\begin{equation}\label{eq:seed-oil-numbers}
 \log_2\epsilon_{\rm seed}
 \le -1048.291,\qquad-1551.477,\qquad-2096.594.
\end{equation}
The bound holds for every fixed hidden oil space; it does not average over key
generation.
\end{proposition}

\begin{proof}
There are $(q^m-1)/(q-1)$ projective oil points.  For any fixed projective
point $[u]$, a uniform seven-dimensional subspace contains $[u]$ with
probability $(q^7-1)/(q^n-1)$.  A union bound proves
\eqref{eq:seed-oil-bound}.  The numerical values substitute the three official
pairs $(n,m)$.
\end{proof}

\subsection{A deterministic Chevalley--Warning extension}

For a quadratic form $Q$ over a field of odd characteristic, write
\[
 B_Q(u,z)=\frac{Q(u+z)-Q(u)-Q(z)}2
\]
for its polar bilinear form.  A subspace is \emph{totally isotropic} for $Q$ when $Q$ vanishes identically on it.

\begin{theorem}[Common-isotropic extension from one seed point]\label{thm:cw-extension}
Let $Q_1,\ldots,Q_s$ be homogeneous quadrics on $\F_q^n$, and let $0\ne w$ satisfy $Q_i(w)=0$ for every $i$.  If
\begin{equation}\label{eq:cw-feasibility}
 n\ge(s+1)r+s,
\end{equation}
then a deterministic finite algorithm constructs an $r$-dimensional subspace $U$ containing $w$ on which every $Q_i$ vanishes identically.

At each extension step the algorithm uses linear algebra and examines at most
\begin{equation}\label{eq:cw-projective-count}
 \left|\PP^{2s}(\F_q)\right|
 =\frac{q^{2s+1}-1}{q-1}
\end{equation}
projective points.
\end{theorem}

\begin{proof}
Suppose a $j$-dimensional common totally isotropic subspace $W$ containing $w$ has already been constructed.  Put
\[
 L_W=\bigcap_{i=1}^s W^{\perp_{B_{Q_i}}}.
\]
The $sj$ orthogonality conditions give
\[
 \dim(L_W/W)\ge n-(s+1)j.
\]
For $j\le r-1$, \eqref{eq:cw-feasibility} makes this dimension at least $2s+1$.  Choose any $(2s+1)$-dimensional subspace $S$ of the quotient.  The $s$ induced quadrics on $S$ have total degree $2s<2s+1$.  The Chevalley--Warning theorem says that their number of common zeros is divisible by the characteristic~\cite{Chevalley1935,Warning1935}.  Since the zero vector is one zero, there is a nonzero zero.  Enumerating \eqref{eq:cw-projective-count} finds one deterministically.

Lift that zero to $z\in L_W$.  Then $Q_i(z)=0$ and $B_{Q_i}(z,W)=0$, so $W\oplus\langle z\rangle$ is again common totally isotropic.  Iterating reaches dimension $r$.
\end{proof}

\begin{corollary}[Target-orthogonal affine reduction]\label{cor:generic-target-slicing}
Let $P:\F_q^n\to\F_q^m$ be any homogeneous quadratic map, let $y\ne0$, and
choose $1\le s<m$.  After an invertible output recombination sends $y$ to the
last coordinate, a nonzero common zero of any $s$ of the first $m-1$
transformed quadrics can be extended to an $(m-s)$-dimensional common totally
isotropic subspace whenever
\begin{equation}\label{eq:generic-square-condition}
 n\ge(s+1)m-s^2.
\end{equation}
On that subspace, the affine problem has $m-s$ equations in $m-s$ variables.
\end{corollary}

\begin{proof}
Apply \cref{thm:cw-extension} with $r=m-s$.  Its condition becomes
$n\ge(s+1)(m-s)+s=(s+1)m-s^2$.  The selected $s$ forms vanish identically on
the constructed subspace and have zero target, leaving the stated affine
residual.
\end{proof}

The affine formulation leaves one variable that homogeneity makes redundant.
The next corollary removes it exactly.

\begin{corollary}[Projective target-orthogonal reduction]\label{cor:projective-target-slicing}
Assume that $q$ is odd.  Under the hypotheses of \cref{cor:generic-target-slicing}, put $r=m-s$ and let
$U$ be the constructed $r$-dimensional subspace.  Normalize the target to
$(0,\ldots,0,1)$ and write
\[
 H_1,\ldots,H_{r-1},H_0
\]
for the restrictions to $U$ of the remaining transformed forms, with $H_0$
corresponding to the nonzero target coordinate.  Then:
\begin{enumerate}
\item affine solutions of
\begin{equation}\label{eq:projective-affine-pair}
 H_1(z)=\cdots=H_{r-1}(z)=0,
 \qquad H_0(z)=1
\end{equation}
modulo the pair $z\sim-z$ are in bijection with projective points
$[z]\in\PP(U)$ satisfying
\begin{equation}\label{eq:projective-zero-system}
 H_1(z)=\cdots=H_{r-1}(z)=0
\end{equation}
and $H_0(z)$ a nonzero square in $\F_q$;
\item a solution of \eqref{eq:projective-affine-pair} has invertible affine
Jacobian if and only if its projective point is nonsingular for
\eqref{eq:projective-zero-system}; and
\item after choosing a uniform nonzero linear form $L\in U^*$, the chart
$L(z)=1$ turns \eqref{eq:projective-zero-system} into $r-1$ quadrics in
$r-1$ affine variables.  Any fixed projective root lies in this chart with
probability
\[
 1-\frac{q^{r-1}-1}{q^r-1}>1-\frac1q.
\]
\end{enumerate}
Thus the square solver core has dimension $m-s-1$ and geometric degree at most
$2^{m-s-1}$.
\end{corollary}

\begin{proof}
For a nonzero common zero $z$ of $H_1,\ldots,H_{r-1}$, replacing $z$ by
$\lambda z$ multiplies every quadratic value by $\lambda^2$.  Hence some
representative satisfies $H_0(\lambda z)=1$ exactly when $H_0(z)$ is a nonzero
square, and then there are exactly two choices $\lambda$ and $-\lambda$.

Euler's identity gives $dH_i(z)(z)=2H_i(z)=0$ for $1\le i<r$, whereas
$dH_0(z)(z)=2$.  The first $r-1$ Jacobian rows therefore have rank $r-1$
exactly when their common kernel is the radial line $\langle z\rangle$; in
that case the last row is not in their span and the affine Jacobian has rank
$r$.  Conversely, an invertible affine Jacobian forces the first $r-1$ rows
to have rank $r-1$.  Finally, among the $q^r-1$ nonzero linear forms on $U$,
exactly $q^{r-1}-1$ vanish at a fixed projective point.  Dehomogenizing in a
chart that contains the point preserves nonsingularity and gives the claimed
square affine system.
\end{proof}

For the QR-UOV parameters, set $s=3$ and $r=m-3$.  The subspace condition is
\begin{equation}\label{eq:official-cw-condition}
 n\ge4m-9.
\end{equation}
The three official profiles give
\[
 (n,4m-9)=(210,207),\quad(306,303),\quad(411,411).
\]
The resulting projective solver dimensions are
\begin{equation}\label{eq:official-projective-cores}
 k=r-1=m-4\in\{50,74,101\}.
\end{equation}
The seed root and the extension are both obtained by finite exhaustive
searches, so this algebraic reduction is unconditional.  One seed search and
the $r-1$ extension steps each evaluate three seven-variable quadrics on at
most $|\PP^6(\F_{127})|$ points.  The companion audit additionally overcharges
uniform seed generation, recomputation of every polar-orthogonal space by
dense elimination, the random chart, and the complete success normalization.
The resulting preprocessing bounds are
\[
 2^{63.545},\qquad2^{64.077},\qquad2^{64.501}
\]
Boolean gates, negligible beside every terminal solver estimate below.

\subsection{An eligible nonsingular projective root with constant probability}

Let $U$ be the subspace from \cref{thm:cw-extension}, and choose an $\F_q$-basis
to identify $U$ with $\F_q^r$.  Under the normalization
\eqref{eq:target-normalization}, write
\[
 H_j=Q_{3+j}|_U\quad(1\le j\le r-1),
 \qquad H_0=Q_m|_U.
\]
The affine residual is
\begin{equation}\label{eq:direct-residual}
 H_1(z)=\cdots=H_{r-1}(z)=0,
 \qquad H_0(z)=1,
 \qquad r=m-3.
\end{equation}
The projective solver uses only the first $r-1$ equations; $H_0$ filters and
scales their projective roots.

The next statement gives an explicit success probability rather than assuming
that a random-looking square system has a root.  Let
\begin{equation}\label{eq:pi-r-direct}
 \pi_{r-1}(q)=\prod_{j=1}^{r-1}(1-q^{-j}).
\end{equation}

\begin{proposition}[Affine nonsingular-root probability]\label{prop:direct-root-probability}
Condition on the selected forms, the attack coins used to construct $U$, and
an off-oil seed point.  Put
\[
 d=\dim_{\F_q}(U\cap\widehat{\cO}),
 \qquad
 \mu=1-q^{d-r}.
\]
Then $d\le r-1$.  Over the independent remaining public-random QR-UOV forms,
the probability that \eqref{eq:direct-residual} has an $\F_q$-rational root
with invertible $r\times r$ Jacobian is at least
\begin{equation}\label{eq:direct-root-lower-bound}
 \frac{\mu\,\pi_{r-1}(q)^2}
 {2+\mu-(q-1)q^{-r}}.
\end{equation}
For $q=127$ and each of $r=51,75,102$, this is at least
\begin{equation}\label{eq:direct-root-number}
 p_*=0.326336998767509\ldots.
\end{equation}
\end{proposition}

\begin{proof}
The off-oil seed point belongs to $U$, so $U$ is not contained in
$\widehat{\cO}$ and $d\le r-1$.  Let $S=U\setminus\widehat{\cO}$ and let $M$
count all roots of \eqref{eq:direct-residual} in $S$.  For a fixed $x\in S$,
\cref{lem:pullback-one-jet} gives
\[
 \Pr[x\text{ is a root}]=q^{-r}.
\]
Hence $\mathbb E[M]=|S|q^{-r}=\mu$.

For two off-oil points, the evaluation functionals are independent unless the
points are scalar multiples.  If $z=\lambda x$ and both solve
\eqref{eq:direct-residual}, homogeneity and the last equation give
$\lambda^2=1$.  By \cref{lem:pullback-one-jet}, the only dependent ordered
pairs that can both solve the target are therefore $z=x$ and $z=-x$; the
other $q-3$ nontrivial base-field multiples have joint probability zero.  Every
nonproportional ordered pair has joint probability $q^{-2r}$.  Consequently,
\begin{equation}\label{eq:root-second-moment}
 \mathbb E[M^2]
 =2\mu+\mu^2-(q-1)\mu q^{-r}.
\end{equation}

Let $N$ count only roots with invertible Jacobian.  At a fixed root $x$, the
conditional one-jet statement makes the Jacobian rows independent and uniform
subject only to Euler's relation $dQ_x(x)=2Q(x)$.  The target vector is
nonzero, so one column is fixed and nonzero while the other $r-1$ columns are
uniform.  Thus
\[
 \Pr[J(x)\text{ invertible}\mid x\text{ is a root}]
 =\pi_{r-1}(q),
\]
and $\mathbb E[N]=\mu\pi_{r-1}(q)$.  Since $N^2\le M^2$, the second-moment
inequality gives
\[
 \Pr[N>0]
 \ge\frac{\mathbb E[N]^2}{\mathbb E[M^2]},
\]
which is \eqref{eq:direct-root-lower-bound}.  The right-hand side increases
with $\mu$, so use $\mu\ge1-q^{-1}$ and substitute the three values of $r$.
\end{proof}

\begin{corollary}[Eligible projective-root probability]\label{cor:projective-root-probability}
Under the same conditioning, the probability that the projective system
\[
 H_1=\cdots=H_{r-1}=0\quad\text{in }\PP^{r-1}(\F_q)
\]
has a nonsingular point $[z]$ for which $H_0(z)$ is a nonzero square is at
least the bound in \eqref{eq:direct-root-lower-bound}, and hence at least
$p_*$ for every official profile.
\end{corollary}

\begin{proof}
By \cref{cor:projective-target-slicing}, every affine nonsingular root of
\eqref{eq:direct-residual} maps to an eligible nonsingular projective root;
the two affine roots $z$ and $-z$ map to the same projective point.  Existence
of an affine root therefore implies the displayed projective event.
\end{proof}

The factor $1/p_*=3.0643169600\ldots$ costs $1.615566$ bits in the
work-to-success ratio.  A random chart adds less than
$\log_2(127/126)=0.011405$ bits.  Unlike a Poisson heuristic, the root bound
follows from the QR-UOV pull-back distribution and a second-moment argument.
The seed-disjointness statement is fixed-key; the residual-root statement
averages over the independent-graph idealized expansion model.  We do not
claim that repeated invocations on every fixed seeded key are independent.

\subsection{Verified projective direct forgery}

\begin{theorem}[Projective target-orthogonal QR-UOV forgery]\label{thm:verified-direct-forgery}
Work in the independent-graph public-random model and the official degree-three
pull-back.  For any nonzero target $y$, the following public algorithm has
one-sided error.
\begin{enumerate}
\item Normalize the target to $(0,\ldots,0,1)$ as in
\eqref{eq:target-normalization}.
\item Choose a uniform seven-dimensional seed subspace, enumerate its
projective points, and take a nonzero common zero of the first three forms.
\item Use \cref{thm:cw-extension} to construct an $(m-3)$-dimensional common
totally isotropic subspace for those forms.
\item Choose a uniform nonzero chart form $L$ on that subspace.  Dehomogenize
the $m-4$ remaining zero-target quadrics in the chart $L=1$ and apply
\cref{thm:finite-char-nonsingular} to the resulting square system in
\[
 k=m-4
\]
variables.
\item For each represented $\F_q$-rational root, recover its projective point,
test whether the last quadratic value is a nonzero square, scale to make that
value one, reverse the public transformations, and evaluate the unchanged
public map.
\end{enumerate}
Every returned vector is a valid signature preimage.  Over the
independent-graph public-random key and the public attack coins, one complete
invocation reaches an off-oil residual with a represented eligible
$\F_q$-rational nonsingular projective root with probability at least
\begin{equation}\label{eq:direct-total-success}
 (1-\epsilon_{\rm seed})p_*
 \left(1-\frac1q\right)
 \left(1-\frac{2^{2k}}{|k'|}\right),
 \qquad k=m-4,
\end{equation}
where $k'/\F_q$ is the solver working extension.  Taking
\[
 [k':\F_q]=15,\qquad22,\qquad30
\]
at Levels I, III, and V gives the finite costs reported in
\cref{sec:security}.  The inverse of \eqref{eq:direct-total-success} is used as
a conservative work-to-success normalization; it is not asserted to be a
fixed-key Las Vegas restart factor.
\end{theorem}

\begin{proof}
The first $m-1$ target coordinates are zero.  By
\cref{prop:off-oil-seed}, the seed subspace is disjoint from the hidden oil
space with probability at least $1-\epsilon_{\rm seed}$ for every fixed key.
Chevalley--Warning then supplies an off-oil seed root, and
\cref{thm:cw-extension} constructs the advertised subspace.  Conditional on
the selected first three forms, the graph, the attack coins, and this off-oil
seed point, the remaining transformed forms are still independent uniform
pull-back forms.  \Cref{cor:projective-root-probability} supplies an eligible
nonsingular projective root with probability at least $p_*$.  A random nonzero
chart contains a fixed such root with probability greater than $1-1/q$.
The finite-characteristic solver includes every root in the chart with the
separator probability in \eqref{eq:direct-total-success}.  Square testing and
scaling recover the associated affine solutions exactly.  Re-evaluation of
all $m$ original public forms rejects every malformed or extraneous candidate,
so the algorithm cannot output an invalid signature.
\end{proof}

\subsection{The fixed-output security scissors}

The direct reduction becomes especially informative when combined with the
directional equivalent-key attack.  Let the quotient degree be $\ell$, so that
\[
 n=\ell V+m.
\]
Applying \cref{cor:generic-target-slicing} with a general value of $s$ gives
the required $(m-s)$-dimensional common-isotropic subspace whenever
\begin{equation}\label{eq:scissors-feasibility}
 \ell V\ge s(m-s).
\end{equation}
By \cref{cor:projective-target-slicing}, the corresponding nonzero-target
solver core has dimension $m-s-1$.

\begin{theorem}[Projective security scissors at fixed output count]\label{thm:security-scissors}
Work over a field of odd characteristic.  Fix $m$ and quotient degree $\ell$.  For every nonzero target, within the two
attack families analyzed in this paper an algebraic core of dimension at most
$m-\ell-1$ is available for every integer vinegar count $V$:
\begin{equation}\label{eq:scissors-bound}
 r_{\rm attack}(V)\le m-\ell-1.
\end{equation}
More precisely:
\begin{enumerate}
\item if $V\le m-\ell-1$, directional equivalent-key recovery uses a square
core of dimension $V$;
\item if $V\ge m-\ell$, target-orthogonal slicing with $s=\ell$, followed by
projectivization, uses a square core of dimension $m-\ell-1$.
\end{enumerate}
These cases cover every integer $V$.  For QR-UOV, $\ell=3$, so the ceiling is
$m-4$.
\end{theorem}

\begin{proof}
In the first case, $V<m$, so the $m$ directional equations admit a
$V$-equation square recombination and the planted-root attack of
\cref{sec:recovery} has core $V\le m-\ell-1$.  In the second case,
\[
 \ell V\ge\ell(m-\ell),
\]
which is \eqref{eq:scissors-feasibility} with $s=\ell$.  The affine direct
residual has dimension $m-\ell$ and
\cref{cor:projective-target-slicing} removes its radial coordinate, leaving
core $m-\ell-1$.  Since consecutive integers $m-\ell-1$ and $m-\ell$ meet the
two inequalities, no case is omitted.
\end{proof}

The theorem is an algebraic-core statement, not a claim that running time
depends only on the number of variables.  For the official degree-three sets,
\cref{prop:off-oil-seed,cor:projective-root-probability} supply the success
factors charged in the concrete ledger.  Its design implication is exact:
once $V$ reaches $m-\ell$, increasing $V$ while holding $m$ fixed no longer
enlarges the projective direct-forgery core.  Conversely, decreasing $V$ into
the other case makes directional recovery no larger than the same ceiling.
Moving both attack lines therefore requires changing $m$ as well as $V$, or
modifying the construction so that one reduction no longer applies.

\subsection{Exact small-parameter execution}
The companion program \texttt{QRUOV\_direct\_qr\_toy.py} instantiates the
degree-three model over $\F_{3^3}$, generates independent extension-field UOV
forms, pulls each back by an explicit trace, hides the oil space with a random
extension-field change of variables, and performs the public attack entirely
over $\F_3$.  It now executes the projective reduction itself: after building
$U$, it enumerates the projective roots of the remaining zero-target forms and
recovers the affine scale from the final equation.  In the frozen $200$-key run
with
\[
 (V,O,n,m,r,k)=(3,2,15,6,3,2),
\]
$199$ keys yielded a verified forgery within twelve outer trials.  Every
returned vector passed both the base-field public map and an independent
extension-field evaluator; the remaining key produced no candidate within the
fixed trial budget.  $185$ of the $200$ final seed subspaces were disjoint from
the hidden oil space, while all $200$ selected seed points were off oil.  The
distinction is expected at this tiny field size: seed-space disjointness is a
sufficient production-scale guarantee, not a condition that the public
algorithm tests.  This is a correctness regression test, not a production
benchmark.

\begin{remark}[Relation to Hashimoto and the specification]\label{rem:direct-prior-art}
The affine reduction to $m-3$ variables is prior art.  Hashimoto's
underdetermined-MQ method and the QR-UOV specification already analyze those
residual sizes~\cite{Hashimoto2023,QRUOVSpec}.  The exact projective quotient
to $m-4$ variables is additional: it uses full target normalization, solves
only the homogeneous zero-target equations modulo scaling, and lets the final
nonzero-target equation reconstruct the affine representative.  The other
upgrades are the finite Chevalley--Warning construction, fixed-key seed bound,
explicit degree-three nonsingular-root probability, and finite-characteristic
solver cost.
\end{remark}

\section{Graph-direction and derivative-kernel recovery}\label{sec:recovery}

The attack has two local endpoints.  The graph-direction endpoint is the stronger structural statement: one affine root determines the complete hidden graph by public linear algebra.  The derivative-kernel endpoint is retained for the faster projective route.  Both endpoints are followed by the same exact graph, pull-back, factorization, and signing checks.

\subsection{A public graph direction produces a planted root}

Fix a direction $c\in K^O$.  Substitute the public oil coordinate $c$ and the negated vinegar coordinate $-x$ into the public form $P_i$:
\begin{equation}\label{eq:direction-system}
 q_{i,c}(x):=P_i(-x,c)
 =x^{\mathsf T}A_ix-2x^{\mathsf T}E_ic+c^{\mathsf T}C_ic,
 \qquad x\in K^V.
\end{equation}
The system $q_{1,c}=\cdots=q_{m,c}=0$ is public and has only $V$ variables.  Define the public $m\times V$ matrix
\begin{equation}\label{eq:direction-matrix}
 \mathcal D_c(x):=
 \begin{pmatrix}
 (E_1c-A_1x)^{\mathsf T}\\[-1mm]
 \vdots\\[-1mm]
 (E_mc-A_mx)^{\mathsf T}
 \end{pmatrix}.
\end{equation}
Recall $B_i=E_i-A_iG$, and write $\mathcal B(c)$ for the matrix whose $i$th row is $(B_ic)^{\mathsf T}$.

\begin{proposition}[Graph-direction identity]\label{prop:graph-direction-identity}
For every $c\in K^O$, the point
\[
 x_c:=Gc
\]
is a common zero of the public system \eqref{eq:direction-system}.  More precisely, for every $z\in K^V$,
\begin{equation}\label{eq:direction-translate}
 q_{i,c}(Gc+z)=z^{\mathsf T}A_iz-2(B_ic)^{\mathsf T}z.
\end{equation}
Consequently,
\[
 \mathcal D_c(x_c)=\mathcal B(c),
 \qquad
 J_q(x_c)=-2\mathcal B(c).
\]
\end{proposition}

\begin{proof}
Expand \eqref{eq:direction-system} at $x=Gc+z$.  Its constant term is
\[
 c^{\mathsf T}(G^{\mathsf T}A_iG-G^{\mathsf T}E_i-E_i^{\mathsf T}G+C_i)c=0
\]
by \eqref{eq:graphid}.  Its linear term is
$2z^{\mathsf T}(A_iG-E_i)c=-2(B_ic)^{\mathsf T}z$, proving \eqref{eq:direction-translate}.  The two matrix identities follow by differentiation and by $E_i-A_iG=B_i$.
\end{proof}

\begin{proposition}[One regular direction determines the entire graph]\label{prop:graph-direction-completion}
Let $c\in K^O$, let $x$ be a common zero of all the equations $q_{i,c}$, and suppose $\rank\mathcal D_c(x)=V$.  For $d\in K^O$, define
\begin{equation}\label{eq:direction-rhs}
 b_{c,d}(x)_i:=c^{\mathsf T}C_id-x^{\mathsf T}E_id
 \qquad(1\le i\le m).
\end{equation}
Then the overdetermined public system
\begin{equation}\label{eq:direction-completion-system}
 \mathcal D_c(x)y=b_{c,d}(x)
\end{equation}
has at most one solution $y\in K^V$.  If $x=Gc$, its unique solution is $y=Gd$.  Thus solving \eqref{eq:direction-completion-system} for the $O$ coordinate vectors $d=e_j$ recovers $G$ exactly.  Requiring every identity \eqref{eq:graphid} then gives a public polynomial-time certificate of a valid expanded UOV decomposition.
\end{proposition}

\begin{proof}
Full column rank gives uniqueness.  For $x=Gc$, multiply \eqref{eq:graphid} on the left by $c^{\mathsf T}$ and on the right by $d$.  Rearrangement yields
\[
 (E_ic-A_iGc)^{\mathsf T}Gd
 =c^{\mathsf T}C_id-(Gc)^{\mathsf T}E_id,
\]
which is exactly row $i$ of \eqref{eq:direction-completion-system}.  Hence $Gd$ is the unique solution.  The final identities are precisely the hypotheses of \cref{prop:verified-key}.
\end{proof}

\begin{proposition}[Several partial-rank directions can be combined]\label{prop:multi-direction-completion}
Let $c_1,\ldots,c_r\in K^O$ and suppose that roots $x_j=Gc_j$ of the corresponding public direction systems have been found.  Stack their public completion matrices and right-hand sides:
\[
 \mathcal D_{\boldsymbol c}(\boldsymbol x)
 :=\begin{pmatrix}\mathcal D_{c_1}(x_1)\\ \vdots\\ \mathcal D_{c_r}(x_r)\end{pmatrix},
 \qquad
 b_{\boldsymbol c,d}(\boldsymbol x)
 :=\begin{pmatrix}b_{c_1,d}(x_1)\\ \vdots\\ b_{c_r,d}(x_r)\end{pmatrix}.
\]
If the stacked matrix has column rank $V$, then for every $d\in K^O$ the public system
\[
 \mathcal D_{\boldsymbol c}(\boldsymbol x)y=b_{\boldsymbol c,d}(\boldsymbol x)
\]
has the unique solution $y=Gd$.  Thus no individual direction needs full rank: enough complementary derivative information across several recovered direction roots also determines the whole graph.
\end{proposition}

\begin{proof}
For each $j$, the proof of \cref{prop:graph-direction-completion} shows that $Gd$ satisfies the $j$th block of equations.  It therefore satisfies the stacked system.  Full column rank gives uniqueness.
\end{proof}

\begin{remark}[A compact attack certificate]\label{rem:graph-certificate}
A successful graph-direction branch can be certified by the tuple $(c,x,G)$ together with one full-rank minor of $\mathcal D_c(x)$.  A verifier checks the $m$ equations $q_{i,c}(x)=0$, the $O$ linear systems \eqref{eq:direction-completion-system}, all $m$ graph identities, the inverse pull-back, and the signing witness.  No claim about the other roots or about completeness of the solver output enters this certificate.
\end{remark}

\subsection{How often a direction is regular, and how to choose a square core}

Put
\begin{equation}\label{eq:rho-mv}
 \rho_{m,V}(Q):=\prod_{a=0}^{V-1}(1-Q^{a-m}),
 \qquad
 \tau_{m,V}(Q):=1-\rho_{m,V}(Q).
\end{equation}
Thus $\rho_{m,V}(Q)$ is the full-column-rank probability of a uniform $m\times V$ matrix over $K$.

The direction system has $m$ equations but only $V$ variables.  The obvious deterministic reduction to a square system tries all $\binom mV$ row subsets.  The next lemma replaces that list by an optimal-size family of public linear combinations.

Index the $m$ equations by $r=0,\ldots,m-1$.  For $0\le a<V$, write
$(r)_a=r(r-1)\cdots(r-a+1)$, with $(r)_0=1$, and define
\begin{equation}\label{eq:wronskian-condenser}
 \mathsf C(t)_{a+1,r+1}:=(r)_a t^{r-a}
 \qquad(t\in K),
\end{equation}
where the entry is zero when $r<a$.  The first $V$ columns of $\mathsf C(t)$ are triangular with diagonal $0!,1!,\ldots,(V-1)!$.  They are therefore nonsingular because $V-1<127=\operatorname{char}K$.

\begin{lemma}[A one-parameter rank condenser]\label{lem:wronskian-condenser}
Let $B\in K^{m\times V}$ have column rank $V$, where $m\le127$.  Then
\begin{equation}\label{eq:wronskian-polynomial}
 W_B(T):=\det\bigl(\mathsf C(T)B\bigr)
\end{equation}
is a nonzero polynomial of degree at most $V(m-V)$.
\end{lemma}

\begin{proof}
For column $j$ of $B$, form the polynomial
$f_j(T)=\sum_{r=0}^{m-1}B_{r+1,j}T^r$.  The columns of $B$ are independent exactly when the polynomials $f_1,\ldots,f_V$ are independent, and $\mathsf C(T)B$ is their ordinary Wronskian matrix.

Choose another basis $g_1,\ldots,g_V$ of their span with distinct increasing degrees
$0\le d_1<\cdots<d_V<m$.  The leading coefficient of the Wronskian of the $g_j$ is, up to the nonzero product of their leading coefficients,
\[
 \det\bigl((d_j)_a\bigr)_{0\le a<V,\,1\le j\le V}
 =\prod_{i<j}(d_j-d_i).
\]
This is nonzero in $K$ because all $d_j$ are distinct and below $127$.  A change of polynomial basis only multiplies the Wronskian by a nonzero scalar, so $W_B$ is nonzero.  Finally,
\[
 \deg W_B
 \le \sum_{j=1}^V d_j-\frac{V(V-1)}2
 \le V(m-V),
\]
where the last bound uses the $V$ largest possible distinct degrees below $m$.
\end{proof}

\begin{proposition}[Key-by-key sampling and public-random density]\label{prop:graph-direction-density}
For a fixed expanded key, suppose some $V\times V$ minor of $\mathcal B(c)$ is not the zero polynomial in $c$.  A fresh uniform $c\leftarrow K^O$ then satisfies $\rank\mathcal B(c)=V$ with probability at least
\begin{equation}\label{eq:random-direction-success}
 1-\frac{V}{Q}.
\end{equation}
Conditional on this event, a fresh uniform $t\leftarrow K$ makes
$\mathsf C(t)\mathcal B(c)$ invertible with probability at least
\begin{equation}\label{eq:random-recombination-success}
 1-\frac{V(m-V)}{Q}.
\end{equation}

In the public-random model, $\mathcal B(c)$ is a uniform $m\times V$ matrix for every fixed nonzero $c\in K^O$.  In particular, one preselected coordinate direction fails with probability $\tau_{m,V}(Q)$, whose base-two logarithm is
\[
 -62.898161,\qquad -62.898161,\qquad -83.864216
\]
at Levels I, III, and V.  More strongly, the $O$ coordinate-direction matrices
$\mathcal B(e_1),\ldots,\mathcal B(e_O)$ are independent uniform matrices.  Therefore
\begin{equation}\label{eq:all-coordinate-directions-bad}
 \Pr\bigl[\rank\mathcal B(e_j)<V\text{ for every }j\bigr]
 =\tau_{m,V}(Q)^O,
\end{equation}
with base-two logarithms
\[
 -1132.166907,\qquad-1635.352198,\qquad-2935.247544.
\]
\end{proposition}

\begin{proof}
Every maximal minor of $\mathcal B(c)$ is homogeneous of degree $V$ in $c$, so Schwartz--Zippel gives \eqref{eq:random-direction-success}.  Conditional on full rank, \cref{lem:wronskian-condenser} and the univariate root bound give \eqref{eq:random-recombination-success}.

Under \cref{def:ideal}, each $B_i$ is an independent uniform linear map $K^O\to K^V$.  For fixed nonzero $c$, the vector $B_ic$ is uniform in $K^V$, independently across $i$.  For the coordinate directions, the vectors $B_ie_j$ use independent columns of the $B_i$, proving both the fixed-direction law and \eqref{eq:all-coordinate-directions-bad}.  The displayed values are exact evaluations of the rank formulas.
\end{proof}

\begin{corollary}[A short deterministic public list]\label{cor:deterministic-outer-list}
Fix any public list of $V(m-V)+1$ distinct elements of $K$.  Whenever $\mathcal B(c)$ has column rank $V$, at least one $t$ on the list makes the square system
\begin{equation}\label{eq:condensed-direction-system}
 \mathsf C(t)\,(q_{1,c},\ldots,q_{m,c})^{\mathsf T}=0
\end{equation}
have the planted root as a nonsingular solution.

For one fixed coordinate direction, the list sizes at Levels I, III, and V are
\[
 105,\qquad153,\qquad307.
\]
Scanning all coordinate directions uses respectively
\[
 1{,}890,\qquad3{,}978,\qquad10{,}745
\]
square systems, and fails in the public-random model only with probability
$\tau_{m,V}(Q)^O$.  Thus all outer choices can be deterministic; only the polynomial-system solver needs random coins.
\end{corollary}

\begin{proof}
The determinant in \eqref{eq:wronskian-polynomial} has degree at most $V(m-V)$ and cannot vanish at every point of a longer distinct list.  At the planted root, the Jacobian of \eqref{eq:condensed-direction-system} is
$-2\mathsf C(t)\mathcal B(c)$ by \cref{prop:graph-direction-identity}.  The numerical counts use $m-V\in\{2,2,3\}$.
\end{proof}

\begin{corollary}[Category-scale deterministic direction lists]\label{cor:compact-direction-list}
Use the first
\[
 r_{\rm dir}=3,\qquad4,\qquad4
\]
coordinate directions at Levels I, III, and V, and for each direction use the condenser list from \cref{cor:deterministic-outer-list}.  The resulting fully public lists contain
\[
315,\qquad612,\qquad1228
\]
square systems.  In the public-random model, the probability that none has the planted point as a nonsingular root is exactly $\tau_{m,V}(Q)^{r_{\rm dir}}$, with base-two logarithms
\begin{equation}\label{eq:compact-direction-logs}
 -188.694,\qquad-251.593,\qquad-335.457.
\end{equation}
These bounds are already below the Level-I, Level-III, and Level-V category exponents $143,207,272$.  Scanning all $O$ coordinate directions gives the stronger probabilities in \eqref{eq:all-coordinate-directions-bad}.
\end{corollary}

\begin{proof}
The coordinate-direction derivative matrices are independent by \cref{prop:graph-direction-density}.  If one of the first $r_{\rm dir}$ matrices has full column rank, \cref{cor:deterministic-outer-list} supplies a condenser value with nonsingular planted Jacobian.  Multiplying the one-direction failure logarithm by $r_{\rm dir}$ gives \eqref{eq:compact-direction-logs}; multiplying the per-direction list lengths $105,153,307$ gives the displayed system counts.
\end{proof}

\begin{remark}[The list length is geometrically optimal]\label{rem:condenser-optimal}
The Grassmannian of $V$-subspaces of an $m$-dimensional space has dimension $V(m-V)$.  For any fixed full-rank linear map $R:K^m\to K^V$, the subspaces on which $R$ is singular form a hyperplane section in Pl\"ucker coordinates.  Over an algebraic closure, at most $V(m-V)$ such hyperplane sections always have a common point.  Hence no family of at most $V(m-V)$ linear recombinations can work for every full-rank $m\times V$ matrix after all field extensions.  The family in \cref{cor:deterministic-outer-list} meets this natural lower bound exactly.
\end{remark}

\paragraph{Randomized versus deterministic use.}
Sampling one field element $t$ is convenient for the expected-cost theorem.  The short list in \cref{cor:deterministic-outer-list} removes that choice entirely and is far smaller than enumerating all $V$-row subsets.

\subsection{Derivative-kernel recovery for the projective route}

For a public common zero $x$, form
\[
C_x=(P_1x\ \cdots\ P_mx)\in L^{(V+O)\times m}.
\]

\begin{lemma}[Derivative-kernel recovery]\label{lem:kernel}
If $x\in\PP(\cO_L)$, then $\cO_L\subseteq\ker C_x^{\trans}$.  If $\rank C_x=V$, equality holds.
\end{lemma}

\begin{proof}
For $u\in\cO_L$, \eqref{eq:isotropy} gives $u^{\trans}P_ix=0$ for every $i$.  If $\rank C_x=V$, its transpose kernel has dimension $(V+O)-V=O$.
\end{proof}

At the oil point, a smooth restricted core has derivative rank $V$.  Total isotropy bounds the ambient derivative rank by the same value, so \cref{lem:kernel} recovers the complete oil space.  Row reduction into the public graph chart, Frobenius descent to $K$, verification of \eqref{eq:graphid}, and \cref{prop:verified-key} then recover the expanded decomposition.  Exact factorization certifies that decomposition, while \cref{prop:signing-search} supplies an invertible oil matrix whose inverse certifies target inversion.

\paragraph{Optional identification with the planted oil space.}
Neither endpoint needs uniqueness: the algorithm keeps every exactly verified candidate and tests all of them for a signing witness.  \Cref{app:uniqueness} separately shows that a second $K$-rational oil space is extraordinarily unlikely.  Outside that additional event, every accepted oil space is the planted one.

\section{Jacobian localization removes global regularity}\label{sec:local-geometry}

Fix $c\in K^O$ and a public recombination $R\in K^{V\times m}$.  Let
\[
 f=(f_1,\ldots,f_V)^{\mathsf T}:=R(q_{1,c},\ldots,q_{m,c})^{\mathsf T}
\]
and put
\begin{equation}\label{eq:local-jacobian-polynomial}
 h(x):=\det J_f(x).
\end{equation}
The graph-direction condition is local: if $\mathcal B(c)$ has full column rank and $R\mathcal B(c)$ is invertible, then \cref{prop:graph-direction-identity} gives
\begin{equation}\label{eq:h-at-planted}
 h(Gc)=(-2)^V\det(R\mathcal B(c))\ne0.
\end{equation}
No assertion is needed about the system on the complementary hypersurface $h=0$.

\begin{lemma}[Localization gives a reduced regular sequence]\label{lem:jacobian-localization}
Let $k$ be a perfect field, let $f_1,\ldots,f_V\in k[x_1,\ldots,x_V]$, and let $h=\det J_f$.  In the localization
\[
 S_h:=k[x_1,\ldots,x_V]_h,
\]
the sequence $f_1,\ldots,f_V$ is regular.  For every $1\le s\le V$, the ideal $(f_1,\ldots,f_s)S_h$ is radical.  Equivalently, the sequence is reduced and regular on the open set $h\ne0$.
\end{lemma}

\begin{proof}
Let $\mathfrak p$ be a prime of $S_h$ containing the first $s$ equations.  Because $h$ is a unit in $S_h$, the full Jacobian matrix is invertible modulo $\mathfrak p$.  Its first $s$ rows are therefore linearly independent.  The Jacobian criterion over a perfect field shows that the local zero set of $f_1,\ldots,f_s$ is smooth of codimension $s$ at every such prime.  Hence the localized quotient is reduced and every minimal prime has height $s$.  The ring $S_h$ is regular, hence Cohen--Macaulay; an ideal generated by $s$ elements and of height $s$ is a complete intersection.  Thus the generators form a regular sequence, and smoothness gives radicality for every prefix.
\end{proof}

\begin{corollary}[Only the planted simple root is needed]\label{cor:local-hypotheses}
If $\mathcal B(c)$ has full column rank and $R\mathcal B(c)$ is invertible, then $Gc$ belongs to the locally closed set
\[
 \{f_1=\cdots=f_V=0,\ h\ne0\}.
\]
The input tuple $(f_1,\ldots,f_V,h)$ satisfies hypotheses (H1)--(H2) of the locally closed Kronecker algorithm of Gim\'enez et al.  These conclusions are independent of every root and irreducible component contained in $h=0$.
\end{corollary}

\begin{proof}
Equation \eqref{eq:h-at-planted} puts the planted root in the open set.  \Cref{lem:jacobian-localization} is exactly the reduced-regular-sequence condition required after localization.
\end{proof}

\begin{remark}[Why this is a separate geometric route]\label{rem:local-independent}
The local theorem does not use the projective slice, the off-oil incidence, the cumulative-degree bound, the joint discriminant, smoothness of the full projective core, or radicality of components on $h=0$.  A failure in any of those arguments can affect the faster projective line without affecting graph-direction recovery.
\end{remark}


\section{Finite-characteristic nonsingular-root recovery}\label{sec:finite-char-solver}

The directional reduction does not need a description of every component of the square core.  It needs one isolated root at which the Jacobian is invertible.  This section adapts the symbolic-homotopy algorithm of Safey El Din and Schost to that exact task in characteristic $127$.

\subsection{The solver target and its degree parameters}
Let $k$ be a perfect field and let
\[
 f=(f_1,\ldots,f_V)\in k[X_1,\ldots,X_V]^V
\]
be a square system of affine quadrics.  Write
\begin{equation}\label{eq:ns-locus}
 Z_{\rm ns}(f)
 :=\{x\in\overline k^V:f(x)=0,\ \det J_f(x)\ne0\}.
\end{equation}
Every point in $Z_{\rm ns}(f)$ is isolated.  B\'ezout's theorem gives
\begin{equation}\label{eq:C-C0}
 C:=2^V,
 \qquad
 C_0:=2^V+V2^{V-1}=2^{V-1}(V+2).
\end{equation}
Here $C$ bounds the number of isolated roots, while $C_0$ bounds the degree of the one-parameter homotopy curve.  Rational reconstruction of a degree-$C_0$ curve uses coefficients through precision $2C_0$; the lifting ledger below therefore runs to that precision.  The second formula is the coefficient sum of
$(\vartheta_0+2\vartheta)^V$ after discarding $\vartheta_0^2$ and $\vartheta^{V+1}$, exactly as in the multihomogeneous degree definition of Safey El Din and Schost~\cite{SafeySchost2018}.

The published algorithm starts from a product system with known simple roots, follows its $C$ branches symbolically, specializes at the target fiber, and removes singular outputs by a Jacobian-based \textsc{Clean} operation.  Its large-characteristic hypothesis supports the particular finite families used for the start and the separator.  We replace both families.

\subsection{A simple product start over any sufficiently large finite field}

\begin{lemma}[Vandermonde product start]\label{lem:ff-vandermonde-start}
Let $k'$ contain $2V$ distinct elements $a_1,\ldots,a_{2V}$.  For $a\in k'$, put
\begin{equation}\label{eq:La}
 L_a(X):=X_1+aX_2+\cdots+a^{V-1}X_V+a^V,
\end{equation}
and, for $1\le i\le V$, set
\begin{equation}\label{eq:vandermonde-start}
 g_i(X):=L_{a_{2i-1}}(X)L_{a_{2i}}(X).
\end{equation}
Then $g=(g_1,\ldots,g_V)$ has exactly $C=2^V$ affine roots over $k'$.  Every root is simple.  A zero-dimensional parametrization of all roots can be constructed in $O(V^2C)$ operations in $k'$.
\end{lemma}

\begin{proof}
Choose one node $b_i\in\{a_{2i-1},a_{2i}\}$ from each pair.  The equations
$L_{b_1}=\cdots=L_{b_V}=0$ have coefficient matrix
\[
 \begin{pmatrix}
 1&b_1&\cdots&b_1^{V-1}\\
 \vdots&\vdots&&\vdots\\
 1&b_V&\cdots&b_V^{V-1}
 \end{pmatrix},
\]
which is Vandermonde and invertible.  Hence the choice gives one solution.

At that solution, the monic polynomial
\[
 T^V+X_VT^{V-1}+\cdots+X_2T+X_1
\]
has $b_1,\ldots,b_V$ as its roots.  It therefore vanishes at none of the $V$ unchosen nodes.  Consequently the solution satisfies exactly one factor in each product $g_i$, and different choices give different solutions.  Conversely, any common zero of the products must choose at least one factor from each pair.  It cannot satisfy both factors in one pair, because that would give at least $V+1$ distinct roots of the displayed monic polynomial.  Thus there are exactly $2^V$ roots.

At a root, row $i$ of the Jacobian of $g$ is the Vandermonde row belonging to the chosen node, multiplied by the nonzero value of the unchosen factor.  The Jacobian is therefore an invertible Vandermonde matrix times an invertible diagonal matrix.  Enumerating the $2^V$ choices and forming the corresponding degree-$V$ polynomials gives the stated $O(V^2C)$ bound.
\end{proof}

This construction uses no division by an integer and no condition on the characteristic beyond the existence of $2V$ distinct field elements.  For the QR-UOV working extensions below, this condition is trivial.

\subsection{A full random separator}
A primitive linear form must separate finitely many branch values and avoid finitely many leading-coefficient vectors.  The published proof draws from a short integer-indexed family; that is the source of its second characteristic bound.  Sampling the entire dual space is simpler over a finite extension.

\begin{lemma}[Separator over a finite extension]\label{lem:full-random-separator}
Let $A$ be a domain containing the finite field $k'$, and let
$z_1,\ldots,z_s\in A^V$ be nonzero.  For a uniform
$\lambda=(\lambda_1,\ldots,\lambda_V)\leftarrow(k')^V$,
\begin{equation}\label{eq:separator-bound}
 \Pr\bigl[\lambda\cdot z_j=0\text{ for some }j\bigr]
 \le \frac{s}{|k'|}.
\end{equation}
In particular, if $s\le C^2$, the separator succeeds with probability at least
$1-C^2/|k'|$.
\end{lemma}

\begin{proof}
Fix a nonzero vector $z$.  Choose an index $h$ with $z_h\ne0$.  After the other $V-1$ coefficients of $\lambda$ are fixed, at most one value of $\lambda_h\in k'$ can make $\lambda\cdot z=0$: two such values would give a nonzero scalar in $k'$ annihilating the nonzero domain element $z_h$.  Thus one vector is killed with probability at most $1/|k'|$.  A union bound proves the claim.
\end{proof}

Safey El Din and Schost identify at most $C^2$ vectors that a well-separating form must avoid: differences between start roots, differences between relevant target specializations, and leading-coefficient vectors of branches~\cite[Lem.~14 and Rem.~15]{SafeySchost2018}.  \Cref{lem:full-random-separator} therefore replaces their integer family without changing the success criterion.

\subsection{The lifting identities do not require large characteristic}
For completeness, we isolate the two local identities used by precision doubling.  They work over a commutative ring and do not divide by factorials or by the characteristic.

\begin{lemma}[Newton--Hensel precision doubling]\label{lem:ring-newton-hensel}
Let $R$ be a commutative ring, let $I\subset R$ be an ideal, and let
$F\in R[X_1,\ldots,X_V]^V$.  Suppose
\[
 F(x)\equiv0\pmod{I^m}
\]
and that $J_F(x)$ is invertible modulo $I^m$.  Put
\begin{equation}\label{eq:ring-newton-update}
 x^+:=x-J_F(x)^{-1}F(x)\pmod{I^{2m}}.
\end{equation}
Then $F(x^+)\equiv0\pmod{I^{2m}}$.
\end{lemma}

\begin{proof}
The correction $h=x^+-x$ lies in $I^m$.  Polynomial Taylor expansion gives
\[
 F(x+h)=F(x)+J_F(x)h+R_2(x,h),
\]
where every monomial in $R_2$ contains at least two coordinates of $h$.  Hence
$R_2(x,h)\in I^{2m}$.  The first two terms cancel by \eqref{eq:ring-newton-update}.  No denominator is used.
\end{proof}

\begin{lemma}[First-order change of primitive parameter]\label{lem:primitive-renormalization}
Work modulo $T^{2m}$.  Let $Q(T,U)$ be monic in $U$, let
$X(T,U)\in(R[T,U]/(T^{2m},Q))^V$, and let $\lambda$ be linear.  Suppose
\[
 e(T,U):=\lambda(X(T,U))-U\in T^mR[T,U]/(T^{2m},Q).
\]
Define
\begin{equation}\label{eq:primitive-update}
 X^*:=X-(\partial_U X)e,
 \qquad
 Q^*:=Q-(\partial_U Q)e.
\end{equation}
Then the change of parameter $U_{\rm old}=U-e(T,U)$ represents the same lifted branches modulo $T^{2m}$, and
\begin{equation}\label{eq:primitive-restored}
 \lambda(X^*)\equiv U\pmod{(T^{2m},Q^*)}.
\end{equation}
\end{lemma}

\begin{proof}
Because $e\in T^m$, substitution of $U-e$ into any polynomial is, modulo $T^{2m}$, its first-order Taylor expansion.  This gives exactly \eqref{eq:primitive-update}.  Moreover
\[
 \lambda(X^*)-U
 =e-(\partial_U\lambda(X))e
 =-(\partial_U e)e.
\]
Differentiation in $U$ does not change $T$-adic order, so $\partial_Ue\in T^m$ and the product lies in $T^{2m}$.  Again, the argument uses no integer denominator.
\end{proof}

\subsection{Proof crosswalk for the finite-characteristic solver}

\begin{proposition}[Where the published characteristic bound is used]\label{prop:characteristic-crosswalk}
For a square system of quadrics over a finite perfect field, replace the start construction in Safey El Din--Schost Lemma~7 by \cref{lem:ff-vandermonde-start}, and replace the separator family in their Lemma~14 by \cref{lem:full-random-separator}.  The remaining correctness proof of their \textsc{NonsingularSolutions} algorithm applies over the extension field without a further lower bound on the characteristic.
\end{proposition}

\begin{proof}
We follow the proof in the order in which its ingredients are used.

\proofstep{Start fiber.}
The original lower bound
$\operatorname{char}k\ge\sum_i d_i$ ensures that its integer-indexed affine forms are distinct and that the product start has the full multihomogeneous root count with invertible Jacobian.  \Cref{lem:ff-vandermonde-start} supplies the same three facts directly: degree two, exactly $C$ explicit roots, and a simple start fiber.

\proofstep{Homotopy curve.}
The homotopy is
\[
 H(T,X)=(1-T)g(X)+Tf(X).
\]
The published construction saturates by the $X$-Jacobian determinant.  Lazard's lemma and the multihomogeneous degree bound then give a one-dimensional curve of degree at most $C_0$.  These arguments require the simple start just established, not an ordering of integers inside the field.

\proofstep{Target specialization.}
The target system can have fewer than $C$ isolated nonsingular roots and may have singular or positive-dimensional components.  The published proof handles this by generalized power series, explicitly because ordinary Puiseux series are not sufficient in arbitrary characteristic~\cite[Sec.~3.4]{SafeySchost2018}.  Its valuation and specialization lemmas therefore already cover the finite-field setting.

\proofstep{Precision doubling and primitive normalization.}
The global lifting step is Newton--Hensel lifting of a geometric resolution.  Its local identities are \cref{lem:ring-newton-hensel,lem:primitive-renormalization}; interpolation and polynomial arithmetic use inverses guaranteed by the simple fiber and separating form.  No factorial or characteristic-dependent scalar is introduced.

\proofstep{Separator and cleaning.}
The published well-separating condition is avoidance of at most $C^2$ nonzero vectors.  \Cref{lem:full-random-separator} supplies it with probability at least $1-C^2/|k'|$.  Finally, the published \textsc{Clean} call removes every point where the target Jacobian vanishes.  As noted in its Remark~15, every output is therefore a subset of the nonsingular target roots; a bad separator may omit roots or cause failure, but cannot introduce an invalid point.

These replacements cover the two uses of the integer-indexed finite families.  All other steps are field algebra, generalized-power-series specialization, or Jacobian cleaning over the perfect extension field.
\end{proof}

\begin{theorem}[Nonsingular roots in characteristic $127$]\label{thm:finite-char-nonsingular}
Let $k$ be a finite field and let $f=(f_1,\ldots,f_V)$ be affine quadrics over $k$, represented by a straight-line program of size $L$.  Let $k'/k$ be a finite extension containing at least $2V$ elements.  There is a randomized algorithm with the following properties.

\begin{enumerate}
\item Every returned point belongs to $Z_{\rm ns}(f)$.
\item Every point of $Z_{\rm ns}(f)$ is represented in the output with probability at least
\begin{equation}\label{eq:finite-char-success}
 1-\frac{C^2}{|k'|}.
\end{equation}
\item The arithmetic cost is
\begin{equation}\label{eq:published-slp-specialization}
 \softO\!\left(
 C C_0\bigl(L+2V+V^2\bigr)V
 \right)
\end{equation}
operations in $k'$.
\end{enumerate}
In particular, $|k'|\ge8C^2$ gives one-invocation inclusion probability at least $7/8$.
\end{theorem}

\begin{proof}
Run the symbolic-homotopy algorithm of Safey El Din and Schost with the start from \cref{lem:ff-vandermonde-start} and a uniform full separator from \cref{lem:full-random-separator}.  \Cref{prop:characteristic-crosswalk} gives correctness over $k'$.  The well-separating event has probability at least \eqref{eq:finite-char-success}; on that event the published specialization recovers every isolated nonsingular target branch.  The final \textsc{Clean} step proves item~1 on every branch, including branches with a bad separator.

For $V$ quadrics in one block, the multihomogeneous quantities in the published Proposition~5 are exactly $C$ and $C_0$ from \eqref{eq:C-C0}; the sum of degrees is $2V$.  The new start costs $O(V^2C)$ and does not change the dominant lifting term.  Substitution into that proposition gives \eqref{eq:published-slp-specialization}.
\end{proof}

\begin{remark}[One-sided failure is enough for cryptanalysis]\label{rem:one-sided-enough}
The original solver paper notes that it cannot cheaply certify completeness of a generic output.  QR-UOV does not need such a certificate.  We impose every unrecombined directional equation, require a $K$-rational root, complete the graph, and verify the unchanged public map.  An omitted planted root causes a restart.  An extraneous or malformed candidate is rejected.
\end{remark}

\subsection{End-to-end verified directional recovery}

For a random $V\times m$ matrix $R$ over $K$, put
\[
 f_{c,R}:=R(q_{1,c},\ldots,q_{m,c})^{\mathsf T}.
\]
If $\mathcal B(c)$ has rank $V$, then $R\mathcal B(c)$ is a uniform $V\times V$ matrix.  Write
\begin{equation}\label{eq:mu-square}
 \mu_V(Q):=\prod_{j=1}^V(1-Q^{-j}),
 \qquad Q=|K|=127^3.
\end{equation}

Choose $K'=\F_{Q^r}$ with
\begin{equation}\label{eq:r-choice}
 r:=\min\{s\ge1:Q^s\ge8\cdot2^{2V}\}.
\end{equation}
The values are
\begin{equation}\label{eq:r-values}
 r=6,\qquad8,\qquad10
\end{equation}
at Levels I, III, and V.

\begin{theorem}[Verified directional equivalent-key recovery]\label{thm:verified-directional-recovery}
Let a fixed expanded \QRUOV{} public key have a valid graph decomposition $G$.  Suppose
\begin{enumerate}
\item some maximal minor of $\mathcal B(c)$ is a nonzero polynomial in $c$; and
\item the recovered planted decomposition has a nonzero signing determinant polynomial $\Delta$ from \cref{prop:signing-search}.
\end{enumerate}
Then the following Las Vegas algorithm returns a publicly verified functional equivalent key.

\begin{enumerate}
\item Sample $c\leftarrow K^O$ and $R\leftarrow K^{V\times m}$.
\item Run \cref{thm:finite-char-nonsingular} on $f_{c,R}$ over $K'$ from \eqref{eq:r-choice}.
\item Intersect the returned parametrization with all $m$ original directional equations and with the $K$-rational locus.
\item For every retained root $x$, run the graph completion of \cref{prop:graph-direction-completion}; require all graph identities, the inverse pull-back, exact public-map factorization, and a signing witness.
\item Return the first accepted key; otherwise restart with fresh public coins.
\end{enumerate}

One invocation includes the planted branch with probability at least
\begin{equation}\label{eq:directional-invocation-success}
 p_{\rm dir}
 \ge
 \left(1-\frac VQ\right)
 \mu_V(Q)
 \left(1-\frac{2^{2V}}{Q^r}\right).
\end{equation}
The expected numbers of full solver invocations are at most
\begin{equation}\label{eq:directional-restart-values}
 1.0000262,
 \qquad1.0000561,
 \qquad1.0202201.
\end{equation}
Every returned key is correct, regardless of whether failed solver invocations were complete.
\end{theorem}

\begin{proof}
By \cref{prop:graph-direction-density}, a random direction has full planted derivative with probability at least $1-V/Q$.  Conditional on that event, $R\mathcal B(c)$ is uniform square, so it is invertible with probability $\mu_V(Q)$.  The planted point is then an isolated nonsingular root of $f_{c,R}$ by \cref{prop:graph-direction-identity}.  \Cref{thm:finite-char-nonsingular} includes it with probability at least $1-2^{2V}/Q^r$.

The all-equation and rational-root filters retain the planted point.  \Cref{prop:graph-direction-completion} reconstructs $G$, and \cref{prop:verified-key,prop:signing-search} supply the exact factorization and signing certificate.  Conversely, any candidate that fails an original equation, completion equation, graph identity, pull-back identity, factorization check, or oil-matrix test is rejected.  This proves zero-error acceptance.  Fresh restarts and \eqref{eq:directional-invocation-success} give the geometric tail and the numerical expectations.
\end{proof}

\begin{corollary}[Public-random key density]\label{cor:directional-public-random}
In the public-random expanded-key model, the probability that every coordinate direction has deficient planted derivative is
$\tau_{m,V}(Q)^O$, with base-two logarithms
\[
 -1132.166907,\qquad-1635.352198,\qquad-2935.247544.
\]
The probability that the deterministic pair-line signing search fails has base-two logarithm below
\[
 -544.820,\qquad-796.276,\qquad-1068.687.
\]
Outside their union, a deterministic scan of coordinate directions and signing witnesses contains a successful outer branch; only the internal nonsingular-root solver requires random coins.
\end{corollary}

\begin{proof}
Combine \cref{prop:graph-direction-density,prop:signing-search}.  The two exceptional events need not be independent; a union bound suffices.
\end{proof}

\section{Dense-quadratic batching and flattened lifting}\label{sec:flat-lifting}

\subsection{Why the direct straight-line-program bound is not the last word}
For dense quadrics, a direct straight-line program has size $L=O(V^3)$.  Substituting it into \eqref{eq:published-slp-specialization} gives
\begin{equation}\label{eq:slp-asymptotic}
 \softO(V^5 4^V)
\end{equation}
operations in $K'$.  This is a valid generic specialization, but it repeatedly charges contractions that all use the same dense quadratic coefficient tensor.

Write
\begin{equation}\label{eq:dense-quad-form}
 f_i(X)=X^{\mathsf T}A_iX+b_i^{\mathsf T}X+c_i.
\end{equation}
At precision $k$, the lifted coordinates live in a coefficient algebra
\begin{equation}\label{eq:coefficient-algebra}
 \mathcal A_k
 :=K'[T,U]/(T^k,Q_k(T,U)),
\end{equation}
where $Q_k$ is monic of $U$-degree at most $C$.  The $V$ coordinate vectors form a $V\times(Ck)$ coefficient matrix.  Flattening the matrices $A_i$ turns all Jacobian contractions into one rectangular matrix product over $K'$.  Quadratic values follow from the Jacobian and $X$, while the Newton solve uses matrix algebra over $\mathcal A_k$.

\begin{proposition}[Dense-quadratic lifting ledger]\label{prop:dense-batched-ledger}
Let $\omega$ be an admissible square-matrix multiplication exponent, and let $\mu(C,k)$ be the cost of one multiplication in $\mathcal A_k$.  The precision-$k$ work of the global Newton lift can be organized using
\begin{equation}\label{eq:precision-k-ledger}
 \softO\bigl(V^3Ck+(V^\omega+V^2)\mu(C,k)\bigr)
\end{equation}
operations in $K'$.  Summed over precision doubling through $2C_0$, this gives
\begin{equation}\label{eq:batched-general}
 \softO\bigl(V^3CC_0+(V^\omega+V^2)\mu(C,2C_0)\bigr).
\end{equation}
\end{proposition}

\begin{proof}
Stack the $V$ matrices $A_i$ into a $V^2\times V$ scalar matrix.  Multiplication by the $V\times(Ck)$ lifted-coordinate matrix forms every linear Jacobian contraction simultaneously in $O(V^3Ck)$ scalar operations.  This term is written explicitly rather than hidden inside the coefficient-algebra multiplication count.

The identity
\[
 f_i(X)=\frac12X^{\mathsf T}\bigl(2A_iX+b_i\bigr)
       +\frac12b_i^{\mathsf T}X+c_i
\]
computes all values with $O(V^2)$ products in $\mathcal A_k$.  Newton matrix inversion or solving costs $O(V^\omega)$ such products, and primitive-element renormalization lies within the same envelope.  The terminal precision is $2C_0$; quasi-linear multiplication makes the last doubling stage dominate up to logarithmic factors.
\end{proof}

\subsection{One multiplication, not a product of two multiplication costs}\label{sec:flat-multiplication}

The earlier stress analysis bounded multiplication in \eqref{eq:coefficient-algebra} by applying a quasi-linear polynomial-multiplication bound first in $U$ and then in $T$.  This gives a product $\Mmul(C)\Mmul(k)$.  A bivariate array can instead be flattened once.

\begin{lemma}[Collision-free Kronecker flattening]\label{lem:kronecker-flat}
Let
$R_k=K'[T]/(T^k)$ and let $Q_k(U)\in R_k[U]$ be monic of degree $C$.  After precomputing the reciprocal of the reversed $Q_k$, one multiplication in
$R_k[U]/(Q_k)$ can be performed using at most three ordinary polynomial multiplications whose input length is below $2Ck$, plus $O(Ck)$ additions.  The reciprocal precomputation costs
$O(\Mmul(2Ck)\log C)$ operations in $K'$.
\end{lemma}

\begin{proof}
For a polynomial
\[
 A(U,T)=\sum_{0\le i<C,\ 0\le j<k}a_{i,j}U^iT^j,
\]
choose the radix $B=2C-1$ and encode
\begin{equation}\label{eq:kronecker-map}
 U^iT^j\longmapsto Z^{i+Bj}.
\end{equation}
In a product, the $U$-exponent is at most $2C-2=B-1$.  It therefore never carries into the $T$ digit.  The largest input exponent is
$(C-1)+B(k-1)<Bk<2Ck$, so one ordinary univariate multiplication recovers the complete bivariate product modulo $T^k$.

To reduce the $U$-degree modulo the monic polynomial $Q_k$, use the standard reciprocal division algorithm over the coefficient ring $R_k$.  With the reciprocal precomputed, one product obtains the quotient estimate and one product forms the quotient times $Q_k$.  Together with the unreduced product, this is three products of the same padded length.  Newton inversion of the reversed monic polynomial gives the stated one-time reciprocal cost over any commutative coefficient ring.
\end{proof}

\begin{corollary}[Flattened lifting complexity]\label{cor:flattened-lifting}
Let $\Mmul(n)$ denote a quasi-linear polynomial-multiplication cost.  The batched lift can be implemented in
\begin{equation}\label{eq:flattened-asymptotic}
 \softO\bigl(V^3CC_0+(V^\omega+V^2)\Mmul(CC_0)\bigr)
 =\softO(V^{\omega+1}4^V)
\end{equation}
operations in $K'$.

For the explicit stress function
\begin{equation}\label{eq:M-explicit}
 \Mmul(n):=n\log_2n\log_2\log_2n,
\end{equation}
the entire precision-doubling schedule for one coefficient-algebra product is bounded by
\begin{equation}\label{eq:flat-schedule-bound}
 3\Mmul(8CC_0).
\end{equation}
\end{corollary}

\begin{proof}
Apply \cref{lem:kronecker-flat} at precisions
$2C_0,C_0,C_0/2,\ldots$.  At the terminal stage, each encoded input has length below $4CC_0$.  Monotonicity of the logarithmic factors gives
\[
 \sum_{j\ge0}\Mmul\!\left(\frac{4CC_0}{2^j}\right)
 <2\Mmul(4CC_0)
 <\Mmul(8CC_0).
\]
There are three ordinary products per modular multiplication.  Reciprocal precomputation contributes an additional logarithmic factor but is dominated by the explicit $V^2$ outer allowance.  Finally, $CC_0=2^{2V-1}(V+2)$.
\end{proof}

\begin{remark}[What was improved and what was corrected]\label{rem:no-double-log-product}
Flattening does not alter the exponential factor $4^V$; it replaces the product of two independent quasi-linear overheads by the overhead at the total coefficient-array length.  The factor $8$ in \eqref{eq:flat-schedule-bound} is important: an earlier draft stopped at precision $C_0$, whereas rational reconstruction requires $2C_0$.  All finite values in \cref{sec:security} use the corrected precision.
\end{remark}

\section{Sparse working-field models}\label{sec:sparse-fields}
Both solvers may choose any convenient representation of their finite working
field.  The directional solver starts over $K=\F_{127^3}$, so its auxiliary
fields must contain $K$; the direct solver starts over $\F_{127}$ and has no
such divisibility requirement.  Projectivization lowers the Level-I direct
working degree from $16$ to $15$.

\begin{lemma}[Explicit sparse auxiliary fields]\label{lem:sparse-fields}
The following polynomials are irreducible over $\F_{127}$:
\begin{align*}
 \phi_{15}(X)&=X^{15}+X^{12}-X+1,
 &\phi_{18}(X)&=X^{18}+X^4+1,\\
 \phi_{22}(X)&=X^{22}+X+1,
 &\phi_{24}(X)&=X^{24}+X^5-X+1,\\
 \phi_{30}(X)&=X^{30}+X^7+1.
\end{align*}
The projective direct-forgery solver uses degrees $15,22,30$ at Levels I,
III, and V.  The directional solver uses degrees $18,24,30$.
\end{lemma}

\begin{proof}
The companion audit applies Rabin's irreducibility criterion~\cite{Rabin1980}.
For every displayed degree $d$, it verifies
\[
 X^{127^d}=X\pmod{\phi_d}
\]
and, for each prime divisor $e$ of $d$,
\[
 \gcd\bigl(\phi_d,X^{127^{d/e}}-X\bigr)=1.
\]
All computations are exact over $\F_{127}$, and the machine-readable output
records every gcd degree.  The same audit certifies the additional sparse
fields used in the retuning table of \cref{tab:combined-retuning}, including
$X^{26}+X^3-X+1$.
\end{proof}

Schoolbook convolution is unnecessary at these modest extension degrees.  The
next statement freezes an explicit recursive circuit rather than invoking an
asymptotic multiplication theorem.

\begin{proposition}[Audited Karatsuba extension-multiplication envelope]\label{prop:sparse-extension-envelope}
Let
\[
 g_{\rm mul}=\Gmul(\log_2 127),
 \qquad g_{\rm add}=4\log_2 127.
\]
For length-$d$ coefficient vectors, let $(M_d,A_d)$ be chosen recursively as
follows.  The schoolbook candidate is
\[
 (d^2,(d-1)^2).
\]
For $d\ge2$, put $a=\lceil d/2\rceil$ and $b=\lfloor d/2\rfloor$.  The
Karatsuba candidate is
\begin{equation}\label{eq:karatsuba-recurrence}
 \bigl(2M_a+M_b,\;2A_a+A_b+4d-4\bigr).
\end{equation}
Choose whichever candidate has smaller value $Mg_{\rm mul}+Ag_{\rm add}$.
If $s_d$ is the number of nonleading terms of $\phi_d$, one multiplication in
$\F_{127^d}$ followed by sparse reduction is then charged by
\begin{equation}\label{eq:sparse-extension-envelope}
 \mathsf G^{\rm sparse}_{127^d}
 :=M_dg_{\rm mul}+\bigl(A_d+s_d(d-1)\bigr)g_{\rm add}.
\end{equation}
For the official working fields, the selected plans are:
\begin{center}
\begin{tabular}{@{}rrrr@{}}
\toprule
$d$ & $M_d$ & $A_d$ & $s_d$\\
\midrule
15&81&332&3\\
18&153&404&2\\
22&225&480&2\\
24&243&512&3\\
30&243&1112&2\\
\bottomrule
\end{tabular}
\end{center}
\end{proposition}

\begin{proof}
The schoolbook counts are exact: $d^2$ coefficient products and
$d^2-(2d-1)=(d-1)^2$ additions.  For the recursive candidate, write each
operand as a low block of length $a$ plus $X^a$ times a high block of length
$b$.  The products of the two low blocks and of their sums have length $a$;
the high-block product has length $b$, giving $2M_a+M_b$ coefficient
multiplications.  Forming the sums, subtracting the low and high products from
the middle product, and combining the three shifted outputs uses at most
$4d-4$ additional coefficient additions, giving
\eqref{eq:karatsuba-recurrence}.  Recursion therefore realizes the displayed
pair $(M_d,A_d)$.

The unreduced product has only $d-1$ high coefficients.  Eliminating each one
modulo $\phi_d$ uses $s_d$ additions or subtractions and no nontrivial constant
multiplication because every nonleading coefficient is $+1$ or $-1$.  This
proves \eqref{eq:sparse-extension-envelope}.  The companion audit executes the
same recursive plan and compares $48$ random products against schoolbook
convolution before certifying the gate counts.
\end{proof}

\section{Local-separator descent and candidate Level-I break}\label{sec:local-separator-upgrade}

Let $q=127$, $a=m-4$, $C=2^a$, and $C_0=2^{a-1}(a+2)$ as in the parent attack.  The supplied audit optimizes the extension degree over certified sparse irreducible polynomials.

\begin{theorem}[Candidate local-separator bound]\label{thm:headline}
Assume the public-random model, homotopy degree bound, and Boolean-gate conventions of the preceding direct-attack construction and the split-before-lift refinement developed here.  Replace global separation by the local conditions in \cref{sec:local}, accept only simple linear target factors, and use the exact descent test in \cref{sec:frob}.  Then the retry-aware diagnostic costs are those in \cref{tab:headline}.  In particular, Level I costs
\[
 2^{140.383995063}
\]
Boolean gates, leaving $2.616004937$ bits below the category target.
\end{theorem}

\begin{table}[h]
\centering
\caption{Audited stress estimates.  ``Global split'' is the preceding split-before-lift bundle.}
\label{tab:headline}
\begin{tabular}{@{}lrrrrr@{}}
\toprule
Level & Parent & Global split & Local split & Target & Margin\\
\midrule
I   & 151.639 & 141.126 & 140.384 & 143 & $+2.616$\\
III & 203.507 & 191.630 & 191.027 & 207 & $+15.973$\\
V   & 260.389 & 247.637 & 247.125 & 272 & $+24.875$\\
\bottomrule
\end{tabular}
\end{table}

The Level-I optimum in the main ledger uses $K'=\F_{127^8}$ with the Rabin-certified modulus
\[
 X^8+X^4-1.
\]
Degree $9$, using $X^9+X^2+1$, gives total $140.418965$ and is retained as a nearby cross-check.

\section{Characteristic polynomials do not need separation}\label{sec:charpoly}

Let $\Gamma\subset\mathbb A^{1+a}$ be the union of the saturated homotopy-curve components that dominate the $t$-line, with degree at most $D=C_0$.  The Jacobian saturation removes the non-etale locus, so over $K'(t)$ its generic fiber is a finite product of separable field extensions,
\[
 B=K'(t)[X_1,\ldots,X_a]/J,
 \qquad \dim_{K'(t)}B=C.
\]
Write $x_j$ for the image of $X_j$ and $M_{x_j}$ for multiplication by $x_j$ on $B$.

Introduce independent coefficients $\Lambda=(\Lambda_1,\ldots,\Lambda_a)$ and the generic linear element
\[
 u_\Lambda=\sum_j\Lambda_jx_j.
\]
Its characteristic polynomial is
\begin{equation}\label{eq:generic-characteristic}
 \chi_\Lambda(t,U)
 =\det\!\left(UI-\sum_j\Lambda_jM_{x_j}\right).
\end{equation}
Schost proves two facts that are exactly suited to the present use: specialization of $\Lambda$ gives the characteristic polynomial of the specialized linear element, and the generic characteristic polynomial can be written
\[
 \chi_\Lambda=\frac{\Xi_\Lambda}{\Theta},
\]
where the $t$-degrees of $\Xi_\Lambda$ and $\Theta$ are at most $\deg\Gamma\le D$ and $\Theta$ is independent of $\Lambda$~\cite[Lemmas~2--3 and Proposition~2]{Schost2003}.  This is the characteristic-polynomial, or generic $u$-resultant, form of the parametric degree theorem.

Schost phrases one corollary using a primitive element with coefficients in the ground field, together with a field-size hypothesis.  The generic characteristic polynomial itself does not require such an element.  After extending to an algebraic closure of $K'(t)$, the finite etale algebra has $C$ distinct geometric points.  For every distinct pair, the difference
\[
 \Lambda\cdot(x_s-x_r)
\]
is a nonzero polynomial in the independent indeterminates $\Lambda$.  Hence $u_\Lambda$ separates the geometric points over the rational-function field $K'(t)(\Lambda)$ and is primitive there, irrespective of the size of $K'$.  Apply Schost's degree argument after this transcendental scalar extension.  The determinant in \eqref{eq:generic-characteristic} is formed from the original multiplication matrices, so its coefficients already lie in $K'(t)[\Lambda,U]$; scalar extension cannot lower a reduced $t$-degree and thereby hide a larger degree over $K'$.  Thus no $K'$-rational global separator, and in particular no $|K'|\gg C^2$ assumption, is used.

\begin{lemma}[Nonseparating eliminant]\label{lem:charpoly}
For every $\lambda\in K'^a$, including a colliding one,
\[
 q_\lambda(t,U)
 :=\prod_{s=1}^{C}\bigl(U-\lambda(x_s(t))\bigr)
 =\chi_\Lambda(t,U)\big|_{\Lambda=\lambda}.
\]
Every coefficient of $q_\lambda$ has a common-denominator representation with numerator and denominator $t$-degrees at most $D$.
\end{lemma}

\begin{proof}
The first identity is the characteristic-polynomial product formula in the reduced generic fiber, with repeated factors allowed after specialization.  The second assertion follows by substituting $\Lambda=\lambda$ in $\Xi_\Lambda/\Theta$.  Since $\Theta$ is independent of $\Lambda$, a separator collision cannot create a new parameter denominator or increase its degree.
\end{proof}

The coordinate interpolation numerators are already encoded in the same generic characteristic polynomial.

\begin{lemma}[Characteristic-polynomial polar]\label{lem:polar}
Define
\[
 v_j(t,U)=\sum_s x_{s,j}(t)
       \prod_{r\ne s}\bigl(U-\lambda(x_r(t))\bigr).
\]
Then
\begin{equation}\label{eq:polar}
 v_j(t,U)
 =-\left.\frac{\partial\chi_\Lambda(t,U)}{\partial\Lambda_j}
   \right|_{\Lambda=\lambda}.
\end{equation}
Consequently each $v_j$, and every fixed affine-linear combination of the $v_j$ and $q_U$, has numerator and denominator $t$-degrees at most $D$, without assuming that $\lambda$ separates the branches.
\end{lemma}

\begin{proof}
Over a splitting field, differentiate
\[
 \chi_\Lambda(t,U)=\prod_s\left(U-\sum_i\Lambda_i x_{s,i}(t)\right)
\]
with respect to $\Lambda_j$ and then specialize $\Lambda=\lambda$.  This gives \eqref{eq:polar}.  Differentiating $\Xi_\Lambda$ in the linear-form coefficients does not increase its $t$-degree, while the denominator $\Theta$ is unchanged.
\end{proof}

This directly covers the sole linear sketch used later:
\[
 b_1(P)=\lambda^{q^{d-1}}\!\cdot P.
\]
Its interpolation numerator is the corresponding fixed $K'$-linear combination of the $v_j$.

For the rational oil filter
\[
 \psi(X)=\frac{H_0(X)}{c+\lambda(X)},
\]
the preceding split-before-lift note proves that the saturated graph curve has degree at most $2D$.  Apply the same construction to its generic linear element
\[
 u_{\Lambda,\mu}=\Lambda\cdot X+\mu Y,
\]
where $Y=\psi(X)$ is the graph coordinate.  Specializing $(\Lambda,\mu)=(\lambda,0)$ gives $q_\lambda$, and
\[
 -\left.\frac{\partial\chi_{\Lambda,\mu}}{\partial\mu}
 \right|_{(\lambda,0)}
 =v_\psi.
\]
The graph degree bound therefore gives numerator and denominator $t$-degrees at most $2D$ for the filter numerator.  Consequently the existing precision
\[
 P=4C_0+1
\]
remains valid.

\begin{remark}
Separation is needed only when one wants the linear element to be a primitive coordinate for every point.  The attack instead keeps the full characteristic polynomial and later accepts only a simple target factor.  Characteristic polynomials and their coefficient polars remain defined at colliding linear forms.
\end{remark}

\section{Only local separation is needed}\label{sec:local}

At target specialization, normalize $q$ and each scalar numerator by the common $t$-valuation exactly as in Safey El Din--Schost Lemma~12~\cite{SafeySchost2018}.  Call $\lambda$ \emph{valuation-generic} if every branch escaping to infinity with minimum coordinate valuation $\mu_s<0$ satisfies
\[
 \ord \lambda(x_s)=\mu_s.
\]
This condition, not pairwise branch separation, is what the specialization proof uses to remove the infinite factors from the specialized eliminant.

\begin{lemma}[Simple-factor recovery]\label{lem:simple}
Assume $\lambda$ is valuation-generic.  Let $\bar q(U)$ and $\bar v_b(U)$ be the normalized target specializations.  If $\tau\in K'$ is a simple root of $\bar q$, then there is exactly one finite represented target point $P$ with $\lambda(P)=\tau$, and
\[
 b(P)=\frac{\bar v_b(\tau)}{\bar q'(\tau)}.
\]
No separation condition is needed away from that one fiber value.
\end{lemma}

\begin{proof}
Let $S_\infty$ and $S_0$ index the escaping and finite branches.  With
\[
 e=-\sum_{s\in S_\infty}\mu_s,
\]
valuation genericity gives
\[
 (t^e q)(0,U)=\gamma\prod_{s\in S_0}
   \bigl(U-\lambda(P_s)\bigr)
\]
for a nonzero scalar $\gamma$.  In the normalized numerator, the summand indexed by a finite branch $s$ specializes to
\[
 \gamma b(P_s)\prod_{r\in S_0,\ r\ne s}
   \bigl(U-\lambda(P_r)\bigr).
\]
A summand indexed by an escaping branch need not vanish term by term; rather, every such surviving summand is divisible by the full finite product $\prod_{r\in S_0}(U-\lambda(P_r))$.  It therefore vanishes modulo $\bar q$, and in particular at every finite root of $\bar q$.  Thus $\bar v_b$ modulo $\bar q$ is the ordinary interpolation numerator on the finite limits.  A simple root $\tau$ belongs to exactly one finite branch, and evaluation followed by division by $\bar q'(\tau)$ gives the claimed value.
\end{proof}

Condition on the existence of a canonical eligible nonsingular base-field target point $P_\star$; for example, choose the first under any fixed ordering.  Nonsingularity makes it an isolated multiplicity-one point.  Let $I$ generic branches escape to infinity and let $F$ have finite target limits, counted with branch multiplicity.  Then $I+F=C$.  There are at most $I$ valuation-cancellation hyperplanes and at most $F-1$ collision hyperplanes against $P_\star$.  Each bad equality
\[
 \lambda(w)=0
 \quad\text{or}\quad
 \lambda(P-P_\star)=0
\]
cuts out at most $|K'|^{a-1}$ choices of $\lambda$ in $(K')^a$.  Thus a uniform $\lambda\in(K')^a$ fails valuation genericity or collides at $P_\star$ with probability at most
\begin{equation}\label{eq:lambdafail}
 \frac{I+F-1}{|K'|}=\frac{C-1}{|K'|}.
\end{equation}
The coefficient vector of the equation may lie only over an algebraic closure, so its solution set need not literally be a $K'$-defined hyperplane.  The count is still valid: choose a nonzero coordinate of that vector and fix all other coefficients of $\lambda$; the remaining equation has at most one solution in $K'$ (and may have none).

For the filter denominator, first compute the $C$ known start values $\lambda(\xi_s)$ and sample $c$ uniformly from the complement of the forbidden set $\{-\lambda(\xi_s)\}_s$ in $K'$.  Rejection sampling costs an expected factor at most $(1-C/|K'|)^{-1}$ for a scan of the start values; this scan is negligible beside one length-$P$ filter lift and is included in the miscellaneous envelope.  Conditional on this choice, only the at most $F\le C$ finite target values remain forbidden.  Therefore
\begin{equation}\label{eq:cfail}
 \Pr[c+\lambda(P)=0\text{ for a finite target branch}]
 \le \frac{C}{|K'|-C}.
\end{equation}
The old $C^2/|K'|$ event has disappeared.

\section{Exact one-sketch Frobenius descent}\label{sec:frob}

Set
\[
 b_1(P)=\lambda^{q^{d-1}}\!\cdot P,
 \qquad K'=\F_{q^d}.
\]
The separator value $\tau=\lambda(P)$ is already available from the factor of $\bar q$, so no numerator for $\lambda\cdot P$ is needed.

\begin{lemma}[Simple-root descent is exact]\label{lem:exactfrob}
Let $\tau\in K'$ be a simple linear factor of the target eliminant, and let $P$ be its unique represented point.  Then
\[
 b_1(P)^q=\tau
 \quad\Longleftrightarrow\quad
 P\in\F_q^a.
\]
\end{lemma}

\begin{proof}
The forward identity for a base-field point is immediate.  Conversely,
\[
 b_1(P)^q
 =\lambda\cdot P^q.
\]
If it equals $\tau=\lambda\cdot P$, then $P$ and $P^q$ have the same separator value.  The target equations and homotopy are defined over $\F_q$, so the represented finite target limits are stable under Frobenius.  If $P^q\ne P$, the two distinct represented points produce multiplicity at least two at $U=\tau$, contradicting simplicity.  Hence $P^q=P$.
\end{proof}

This eliminates the probabilistic $C/|K'|$ false-positive term.  It also eliminates the abort cap: every simple factor passing the Frobenius test and the nonzero-square oil filter is a genuine base-field target solution.  The algorithm reconstructs coordinates for the first such factor and stops, so exactly one coordinate re-lift is charged per successful forgery.

\section{The sharpened root bound is explicit}\label{sec:rootbound}

For completeness, this pass expands the probability calculation inherited from the split-before-lift note.  Let $M$ count all off-oil affine roots of the normalized residual system, and let $N$ count those with invertible affine Jacobian.  The parent manuscript proves, for an actual parameter $\mu\ge 1-q^{-1}$,
\[
 \mathbb E[M]=\mu,
 \qquad
 \mathbb E[M^2]=2\mu+\mu^2-(q-1)\mu q^{-r},
 \qquad
 \mathbb E[N]=\mu\pi,
\]
where $\pi=\prod_{j=1}^{r-1}(1-q^{-j})$.  The roots occur in pairs $z,-z$, and each nonsingular normalized pair yields one eligible projective point.  Applying the second Bonferroni inequality to the projective-point events and bounding pairs of nonsingular roots by pairs of all roots gives
\begin{align*}
 \Pr[\text{eligible nonsingular projective root}]
 &\ge \frac{\mathbb E[N]}2
      -\mathbb E\!\left[\binom{M/2}{2}\right]\\
 &=\frac{\mu\pi}{2}-\frac{\mu^2}{8}
   +\frac{(q-1)\mu q^{-r}}8.
\end{align*}
The right-hand side is increasing for $\mu\in[1-q^{-1},1]$, so substituting $\mu=1-q^{-1}$ yields
\[
 p_{\rm root}=0.3690869822\ldots.
\]
As a separate fallback, the ledger also substitutes the parent manuscript's already-proved lower bound $0.326336998767509\ldots$; the Level-I total then becomes $140.531461$, still $2.469$ bits below target.

\section{Finite ledger}

At Level I the local variant has two scalar families, $b_1$ and $\psi$, instead of three.  The terminal product tree therefore uses
\[
 \frac a2(1+2\cdot2)=125
\]
long-product units.  Rational reconstruction contributes $99.885$ units.  Candidate coordinate evaluation is charged once over the complete Las Vegas execution, rather than per trial, and the $64$-unit miscellaneous cushion is retained.  This pass also charges the one-time formation of $u_s=\lambda\cdot x_s$ during that coordinate rerun.

\begin{table}[h]
\centering
\caption{Level-I component ledger for extension degree $d=8$.}
\label{tab:components}
\begin{tabularx}{0.88\textwidth}{@{}Xr@{}}
\toprule
Component & $\log_2$ gates\\
\midrule
Initial base-field branch lifts & 139.358\\
Valuation-safe filter series & 132.120\\
Terminal scalar recombination & 138.791\\
One on-demand coordinate re-lift & 137.860\\
Extension-valued leaf formation & 127.610\\
Coordinate-rerun separator leaves & 124.898\\
Factorization and exact verification envelope & 115.707\\
Preprocessing & 63.365\\
\midrule
\textbf{Total} & \textbf{140.383995}\\
\bottomrule
\end{tabularx}
\end{table}

The one-trial success lower bound is $0.353997$, giving restart factor $2.824886$.  At $d=8$, the local-separator failure bound in \eqref{eq:lambdafail} is $0.016637$, and the conditional target-denominator bound in \eqref{eq:cfail} is $0.016918$.

\begin{table}[h]
\centering
\caption{Level-I sensitivity to the multiplier on $\M(n)\log n$ in rational reconstruction, retaining $d=8$.}
\label{tab:sensitivity}
\begin{tabular}{@{}lrrrrrr@{}}
\toprule
Multiplier & 8 & 10 & 16 & 32 & 64 & 128\\
\midrule
Total & 140.384 & 140.424 & 140.538 & 140.805 & 141.225 & 141.815\\
\bottomrule
\end{tabular}
\end{table}

\begin{table}[ht]
\centering
\caption{Level-I cross-checks.  Each row reoptimizes the extension degree.}
\label{tab:robustness}
\begin{tabular}{@{}lrrrrr@{}}
\toprule
Online-envelope multiplier & $1$ & $2$ & $4$ & $8$ & $16$\\
\midrule
Best extension degree & $8$ & $8$ & $9$ & $9$ & $9$\\
Total & 140.384 & 141.117 & 141.952 & 142.856 & 143.806\\
\bottomrule
\end{tabular}
\end{table}

Thus the candidate remains below $2^{143}$ throughout the rational-reconstruction sensitivity range.  It also remains below target after multiplying the entire claimed online-product envelope by eight, but not by sixteen; the finite margin therefore does not excuse the remaining obligation to provide a fully explicit schedule.  The aggregate Level-I margin corresponds to an overall multiplicative cushion of $2^{2.616}=6.13$ in the inherited gate model; this does not address memory traffic or a change of model.


\section{Reduced-parameter end-to-end instantiation}\label{sec:implementation}

The artifact now contains an executable instantiation of the complete attack rather than only independent checks of its ingredients.  The reduced scheme follows the QR-UOV algebraic construction in the official specification~\cite{QRUOVSpec}: central symmetric matrices over an extension field have a zero oil--oil block, the secret change of variables has the compact form
\[
 S=\begin{pmatrix}I&G\\0&I\end{pmatrix},
\]
and the extension-field matrices are pulled back to base-field quadratic forms through a nonzero linear functional, implemented here by the trace pairing.  Seed expansion, byte serialization, and the exact domain-separated hash interface are deliberately simplified; the public quadratic relation attacked by the algorithm is unchanged.

The default demonstration uses
\[
 q=13,\qquad \ell=3,\qquad V=3,\qquad M=2,
\]
so the expanded UOV dimensions are
\[
 v=9,\qquad m=o=6,\qquad n=15.
\]
The residual projective core has $a=m-4=2$, hence $C=4$, $C_0=8$, and precision $P=4C_0+1=33$.  The attack function receives only the expanded public matrices, message, and salt.  It executes target normalization, the anchored Chevalley--Warning construction, a common totally isotropic subspace for three forms, affine charting, all four coefficientwise homotopy lifts, local-separator recombination, rational reconstruction, valuation-safe specialization, exact Frobenius descent, one coordinate rerun, affine rescaling, and unchanged public verification.

For the canonical seed $20260731$, the program first generated an honestly verifiable signature and then produced the forged signature
\[
 (0,8,2,10,0,0,0,4,10,5,11,8,0,12,1)
 \in\F_{13}^{15}
\]
for the recorded message and salt.  The public verifier accepted it.  The successful run used the trace in \cref{tab:toytrace}.

\begin{table}[ht]
\centering
\caption{Canonical reduced-parameter attack trace.  Wall time is machine-dependent and is not used in the security ledger.}
\label{tab:toytrace}
\begin{tabular}{@{}lr@{}}
\toprule
Quantity & Value\\
\midrule
Salt attempts / attack records & $1/1$\\
Target eliminant degree & $4$\\
Branch lifts & $8$ ($4$ initial, $4$ coordinate rerun)\\
Checked lifted-branch residuals & $4$\\
Rational reconstructions & $31$\\
Coordinate reruns & $1$\\
Target roots examined / simple & $2/1$\\
Frobenius passes / exact verified candidates & $1/1$\\
Exhaustive $\F_{13}$ / $\F_{13^2}$ residual roots & $2/2$\\
Python wall time & $0.415$ seconds\\
\bottomrule
\end{tabular}
\end{table}

The exhaustive residual-root enumeration is deliberately performed only \emph{after} the local-separator algorithm has selected and reconstructed its candidate.  It is not a solver subroutine.  It checks that every extension-field root appears in the specialized eliminant, that every simple eliminant root has a singleton fiber, and that the accepted point is among the base-field roots.

The canonical transcript records the public key, target, isotropic subspace, chart, affine core, lifted-system parameters, separator, specialized characteristic polynomial and scalar polars, accepted root, recovered coordinates, signature, and operation counts.  A separate public replay program consumes that transcript without the secret key, reconstructs the affine core, reruns the complete local-separator solver with the recorded separator, and reproduces the accepted forgery.

A fixed regression suite over the five consecutive seeds $20260000,\ldots,20260004$ produced five publicly valid forgeries and five successful public replays.  At the imposed finite retry limits, the largest salt count was two and the slowest individual run took $10.252$ seconds.  This is a regression test, not a statistical estimate of the asymptotic success probability.  A supplementary run over $q=29$ with the same dimensions also forged and replayed successfully.

This closes the narrow composability objection: the artifact now instantiates one finite end-to-end attack on reduced QR-UOV instances.  It does \emph{not} establish the Level-I gate count.  The prototype uses naive truncated-series arithmetic at $a=2$ and therefore neither implements nor benchmarks the fast dyadic multiplication schedule assumed by the full-parameter ledger.  Full parameters remain infeasible to execute, and the public-random-model, seeded-transfer, characteristic-polynomial degree, and memory-model caveats remain.

\section{Validation and remaining risk}

The bundle includes:
\begin{enumerate}
\item \texttt{QRUOV\_local\_separator\_audit.py}, which computes the optimized degree, finite probabilities, and all component costs;
\item \texttt{verify\_local\_separator\_ledger.py}, a separately written implementation importing neither the primary audit nor its parent modules, matching all displayed totals to below $5\cdot10^{-10}$ bits;
\item the end-to-end reduced attack \texttt{implementation/qruov\_local\_separator\_forgery.py}, its public replay checker, five-instance regression suite, and complete public transcripts;
\item Rabin certificates for every newly used sparse modulus;
\item the original split-before-lift toys validating branch lifting, symmetric recombination, and Frobenius arithmetic, together with an exhaustive $\F_9/\F_3$ toy showing that a non-base point passes the one-sketch test exactly on separator collisions; and
\item \texttt{toy\_infinite\_branch\_specialization.py}, which exhibits the subtle case in which an escaping-branch numerator term survives normalization but is a multiple of the finite eliminant and hence vanishes at every finite simple root.
\end{enumerate}

The principal remaining proof obligations are external checks that the characteristic-polynomial degree theorem in \cref{lem:charpoly,lem:polar} applies to the exact saturated dominant homotopy algebra and to the rational-filter graph with the stated degree bounds.  The determinant identities are exact and no defect is currently known, but this identification with the parent solver's saturated dominant algebra is what the degree-$8$ retuning depends on.  The reduced implementation supplies a literal finite schedule for the toy instance, but the dyadic fast online-product envelope used at Level I still needs a line-by-line full-scale algorithmic derivation; the multiplier table shows exactly how much error that part can absorb.  The earlier caveats about the idealized public-random model, seeded transfer, and astronomical memory traffic are unchanged.

\section{Assessment}

This is a real strengthening rather than a cosmetic ledger improvement.  The preceding candidate crossed Level I by splitting before lifting; the present argument removes the largest remaining auxiliary-field overkill.  It replaces a global primitive-element event with a one-sided local isolation event and turns Frobenius filtering from probabilistic to exact.  Under the inherited model, the cumulative improvement over the strict parent bundle is $11.254751$ bits.

The result is now robust to far larger rational-reconstruction constants than the original $1.87$-bit-margin version.  The reduced implementation materially strengthens the evidence: the claimed transformations and filters have now been composed into an actual public-key forgery rather than tested only in isolation.  This manuscript distinguishes this demonstrated toy-scale attack from the Level-I extrapolation and should await specialist verification that the characteristic-polynomial degree theorem applies to the exact saturated dominant algebra and its rational-filter graph.



\section{Combined cost, memory, and parameter implications}\label{sec:security}

\subsection{Work unit and scope}
The QR-UOV specification expresses classical attack costs as Boolean-gate
exponents and charges multiplication of $b$-bit field elements by
\begin{equation}\label{eq:spec-gmul-new}
 \Gmul(b)=2b^2+b
\end{equation}
Boolean gates~\cite{QRUOVSpec}.  We retain that broad unit.  Every
extension-field operation is overcharged as one multiplication using the
explicit Karatsuba circuit of \eqref{eq:sparse-extension-envelope}.

The principal symbolic-homotopy term uses the corrected reconstruction
precision $2C_0$.  For a square solver core of dimension $a$,
\[
 C=2^a,
 \qquad C_0=2^{a-1}(a+2),
\]
and the flattened lifting schedule is charged by
\begin{equation}\label{eq:unified-flat-cost}
 3(a^3+a^2)\Mmul(8CC_0)
\end{equation}
operations in the chosen working field.  The projective direct row uses
$a=m-4$ and includes the fixed-key seed bound, nonsingular-root factor, random
chart, and finite-field separator.  The directional row uses $a=V$ and
includes direction regularity and its separator.  Polynomial-time
preprocessing, all-equation filtering, factorization, graph completion,
square testing and rescaling, and public verification are included through
explicit lower-stage envelopes; they do not change the displayed decimals.

These are reproducible diagnostics for symbolic algorithms, not measured
circuits.  In particular, setting
\[
 \Mmul(t)=t\log_2t\log_2\log_2t
\]
assigns unit coefficient to an asymptotic multiplication bound.  The exact
reductions, Karatsuba circuits, sparse fields, and success probabilities do not
depend on that convention; the final conversion to a practical implementation
does.

\subsection{The two attacks}
\begin{table}[t]
\centering
\small
\setlength{\tabcolsep}{5pt}
\caption{Classical gate diagnostics, including the local-separator split-before-lift refinement and the audited Karatsuba field circuit.  A positive margin means that the attack row lies below the
category target.}
\label{tab:combined-costs}
\begin{tabular}{@{}lrrrrrr@{}}
\toprule
& \multicolumn{2}{c}{Level I} & \multicolumn{2}{c}{Level III} & \multicolumn{2}{c}{Level V}\\
Attack & Cost & Margin & Cost & Margin & Cost & Margin\\
\midrule
Original projective forgery
 &151.639&$-8.639$&203.507&$3.493$&260.389&$11.611$\\
Local-separator projective forgery
 &140.384&$2.616$&191.027&$15.973$&247.125&$24.875$\\
Directional equivalent-key recovery
 &154.826&$-11.826$&206.177&$0.823$&260.857&$11.143$\\
\midrule
Best attack
 &\textbf{140.384}&$\mathbf{2.616}$
 &\textbf{191.027}&$\mathbf{15.973}$
 &\textbf{247.125}&$\mathbf{24.875}$\\
\midrule
Category target &143&&207&&272&\\
\bottomrule
\end{tabular}
\end{table}

The original direct rows use projective solver dimensions $50,74,101$.
The local-separator refinement keeps those algebraic cores but lowers the
working-field degrees and removes global branch separation, the second
Frobenius numerator, probabilistic false candidates, and repeated coordinate
recovery.  Its optimized field degrees are $8$, $12$, and $16$ in the supplied
ledger.  The directional rows remain an independent equivalent-key route.
The reduced implementation executes the local-separator route at toy scale;
the full-parameter symbolic lift still dominates both time and memory.

For context, the QR-UOV team's January 2026 response reports classical
Hashimoto estimates of $158,217,285$ bits for the same three public
dimensions~\cite{QRUOVResponse2026}.  The new direct rows are lower by
approximately $6.36,13.49,24.61$ bits.  This is not a claim that the
implementations are identical: the team uses a semi-regular direct-MQ model,
whereas our row uses the finite-characteristic nonsingular-root solver and its
explicit success probability.

\subsection{What follows at each level}
\paragraph{Level I.}
The local-separator row is $140.384$, or $2.616$ bits below the $143$-bit
target.  Thus the paper gives a candidate Level-I break in the stated
Boolean-gate model.  This is not a practical attack claim: the symbolic state
is physically unrealizable, and the full-scale online-product schedule and
characteristic-polynomial degree transfer remain explicit obligations.  The
reduced-parameter implementation establishes that the algebraic pipeline is a
single executable attack, not that the Level-I extrapolation is practical.

\paragraph{Level III.}
The local-separator row is $191.027$, leaving $15.973$ bits below target.  The
older projective and independent equivalent-key rows are also below target.
The combined evidence strongly supports changing the Level-III parameters.

\paragraph{Level V.}
The local-separator row is $247.125$, leaving $24.875$ bits below target.  The
older projective and independent equivalent-key rows are likewise below target.
Level V should not be retained unchanged.

\subsection{Why changing only the vinegar count is not a repair}
The official parameter sets sit on the direct side of the projective scissors
crossing:
\begin{center}
\begin{tabular}{@{}lrrrr@{}}
\toprule
Level & $V$ & $m-4$ & Directional core & Projective direct core\\
\midrule
I&52&50&52&50\\
III&76&74&76&74\\
V&102&101&102&101\\
\bottomrule
\end{tabular}
\end{center}
Increasing $V$ raises the directional core but leaves the projective direct
core unchanged.  Decreasing $V$ eventually makes directional recovery no
larger than $m-4$.  Consequently a fixed-$m$ vinegar retuning cannot move the
combined algebraic bottleneck.

The next table gives attack-specific thresholds under the same padded ledger.
It is not a complete replacement proposal: every other attack, failure
probability, key size, and implementation cost must be recomputed.  The point
is that both columns must move.  Raising only $V$ leaves the direct attack
unchanged; raising only $m$ leaves the directional attack unchanged.

\begin{table}[t]
\centering
\small
\caption{Smallest quotient-compatible $(V,m)$ pairs found by the companion
audit that put both displayed attack rows at the category target, or at the
target plus a 16-bit attack-specific cushion.}
\label{tab:combined-retuning}
\begin{tabular}{@{}lrrrrrr@{}}
\toprule
& \multicolumn{2}{c}{Current} & \multicolumn{2}{c}{Reach target} & \multicolumn{2}{c}{Reach target $+16$}\\
Level & $V$ & $m$ & $V$ & $m$ & $V$ & $m$\\
\midrule
I&52&54&52&54&57&60\\
III&76&78&78&81&87&90\\
V&102&105&108&111&117&120\\
\bottomrule
\end{tabular}
\end{table}

The table enforces $m\equiv0\pmod3$ and the direct-side feasibility condition
$V\ge m-3$.  Level I already exceeds its bare target in this model, but a
16-bit cushion requires moving both dimensions.  At Levels III and V, even the
target-only diagnostic requires increases in both $V$ and $m$; a 16-bit
cushion requires substantially larger changes.  A structural redesign may be
preferable to paying both costs.

\subsection{Memory remains the central practical limitation}\label{subsec:memory-new}
The desired outputs are small: one direct attack needs a single vector in $\F_q^n$, and the equivalent-key attack ultimately needs a $V\times O$ graph.  The direct symbolic representations are not small.  One flattened coefficient
array at terminal precision requires approximately
\[
 2^{49.41}\ \mathrm{EiB},
 \qquad2^{98.51}\ \mathrm{EiB},
 \qquad2^{153.40}\ \mathrm{EiB}
\]
for the projective direct cores.  The corresponding directional arrays require
about $2^{53.73}$, $2^{102.68}$, and $2^{155.41}$ EiB.  A full-parameter
execution using these representations is physically unrealistic.

This limitation prevents a claim of a practical full-parameter implementation.  No full-parameter forgery or equivalent-key recovery is reported, and the paper draws no conclusion from unimplemented attack variants.  The exact reductions and the fixed-output scissors theorem are algebraic results; time, storage, and correctness are reported as separate resources.

\section{Executable audits}\label{sec:audits-new}
The bundle contains the following executable checks and an end-to-end reduced-parameter attack.

\paragraph{Combined numerical audit.}
\texttt{QRUOV\_scissors\_audit.py} recomputes the Chevalley--Warning slacks,
the projective $m-4$ cores, fixed-key seed--oil intersection bound,
second-moment nonsingular-root bound, chart and separator probabilities,
Karatsuba field circuits, direct and directional cost rows, memory envelopes,
and attack-specific retuning thresholds.  It also applies Rabin's criterion to
every sparse auxiliary-field polynomial used in the official and threshold
rows.  The script emits both human-readable and machine-readable output.

\paragraph{Directional algebra audit.}
\texttt{QRUOV\_directional\_audit.py} checks the corrected $2C_0$ precision,
Vandermonde product starts, characteristic-$127$ lifting identities, flattened
bivariate multiplication, sparse-field certificates, the recursive Karatsuba
plans, and the complete graph-recovery pipeline on $200$ random toy instances.  Of those instances, $199$ had a full-rank planted direction and all $199$ yielded the exact graph; the remaining direction was correctly rejected for retry.

\paragraph{Exact degree-three direct-forgery toy.}
\texttt{QRUOV\_direct\_qr\_toy.py} generates independent extension-field UOV
forms over $\F_{3^3}$, hides their oil space, pulls each form back by a trace to
a scalar base-field quadric, and performs full target normalization, random
seed-subspace search, Chevalley--Warning extension, projective residual search,
affine rescaling, and public verification.  In the frozen $200$-key run, $199$ keys yielded a verified forgery within twelve trials.  Every returned vector passed two independent evaluators; $185$ final seed subspaces were disjoint from the hidden oil space, and every selected seed point was off oil.  The experiment checks the attack logic; it is not a full-parameter benchmark.


\paragraph{End-to-end local-separator forgery.}
\texttt{implementation/qruov\_local\_separator\_forgery.py} instantiates the
complete reduced attack described in \cref{sec:implementation}.  Its public
replay checker reconstructs the residual system and reproduces the forgery
without secret material.  The frozen regression suite succeeds on five of five
instances, and the supplementary $q=29$ instance also succeeds.  Exhaustive
root enumeration is used only after acceptance as a validation oracle, never as
the attack solver.

\section{Conclusion}\label{sec:conclusion}
QR-UOV admits a projective direct-forgery reduction with square cores of
dimensions $50,74,101$ and a complementary directional attack that recovers an
equivalent signing key.  The local-separator refinement changes the concrete
conclusion: global branch separation is unnecessary, characteristic-polynomial
polars suffice for scalar recovery, and one exact Frobenius sketch identifies a
base-field point on a simple target factor.  Under the stated symbolic
Boolean-gate model, the resulting rows lie below all three NIST category
targets, including Level I by $2.616$ bits.

The reduced-parameter implementation now composes the complete attack and
produces publicly accepted forgeries from public inputs alone.  Consequently,
the result is no longer merely a collection of independently tested
micro-lemmas.  What remains conditional is the full-parameter complexity
projection: the exact characteristic-polynomial degree transfer to the
saturated dominant algebra, the fast online-product schedule, seeded-instance
transfer, and memory traffic all require further scrutiny.  The paper therefore
claims a finite, executable structural attack and a model-based candidate
full-parameter break, not a practical full-parameter computation.

The parameter lesson is unchanged but stronger.  Increasing only the vinegar
count cannot repair the fixed-output scissors, while the new direct row already
crosses Level I.  Any QR-UOV revision should reconsider both dimensions and the
underlying compressed structure rather than apply a vinegar-only retuning.

\section*{Appendix overview}
The main body presents the direct forgery, directional equivalent-key recovery, characteristic-$127$ solver, and combined parameter consequences.  The first appendices retain an independent smooth-projective route and two solver cross-checks.  Later appendices cover seeded expansion, optional identification with the planted oil space, alternative low-space directions, and executable pseudocode.  These details are separated so that readers can understand the attack and its parameter implications before entering the algebraic-geometry audit.

\section{Independent projective smooth-core route}\label{sec:projective}

This section supplies the geometric input to \cref{thm:attack}.  Outside three explicit bad key loci, we show that some slice and output recombination gives a smooth $2^V$-point projective core; we then bound the density of those loci.  The next section chooses a finite working field and turns this existence statement into a verified randomized algorithm.  Secret central coordinates are used only to expose the key distribution.

Let $K=\F_{127^3}$, $Q=|K|$, and
\[
(V,O,m)\in\{(52,18,54),(76,26,78),(102,35,105)\}.
\]
Put $e=m-V\in\{2,2,3\}$ and $N=V+O$.  Work first in secret central coordinates
\[
\cV=K^V\oplus K^O,
\qquad
\cO_0=\{0\}\oplus K^O.
\]
The $m$ central quadrics are
\begin{equation}\label{eq:central-slice}
F_i(v,o)=v^{\mathsf T}A_i v+2v^{\mathsf T}B_i o,
\qquad
A_i\in\Sym_V(K),\ B_i\in K^{V\times O}.
\end{equation}
In the public-random expansion model, the pairs $(A_i,B_i)$ are independent and uniform.  Equivalently, each $F_i$ is uniform in
\[
\cU:=H^0\!\left(\PP(\cV),\mathcal I_{\PP(\cO_0)}(2)\right),
\qquad
D_{\rm form}:=\dim\cU=\frac{V(V+1)}2+VO.
\]
The proof uses central coordinates only to expose the distribution.  The secret change from public to central coordinates lies in $\operatorname{GL}_N(K)$.  Left multiplication by its inverse is a bijection on $L^{N\times(V+1)}$, so a uniform public slice remains uniform in central coordinates.

Smoothness, intersection multiplicity, and derivative-kernel recovery are coordinate invariant.  In particular, the public and central tuples have identically vanishing joint discriminant polynomials simultaneously.  Let
\[
X_F:=V(F_1,\ldots,F_m)\subset\PP(\cV).
\]
It contains the planted oil space $\PP(\cO_0)\simeq\PP^{O-1}$.

\begin{proposition}[Exact marginal law on a fixed transverse slice]\label{prop:planted-slice-law}
Fix, independently of the sampled central forms, a $K$-rational projective $V$-plane $\Lambda$ that meets $\PP(\cO_0)$ transversely in one point $p$.  Under the marginal public-random law of the central pairs $(A_i,B_i)$, the restrictions $F_1|_\Lambda,\ldots,F_m|_\Lambda$ are independent uniform elements of
\[
H^0\!\left(\PP^V,\mathcal I_p(2)\right),
\]
the $K$-vector space of projective quadrics through $p$.
\end{proposition}

\begin{proof}
Choose homogeneous coordinates $[z:t_1:\cdots:t_V]$ on $\Lambda$ with $p=[1:0:\cdots:0]$.  Transversality makes the vinegar projection of $\Lambda/p$ invertible, so after changing the $t$-coordinates its central-coordinate parametrization is
\[
v=t,\qquad o=cz+Lt,
\]
for some nonzero $c\in K^O$ and $L\in K^{O\times V}$.  Hence
\[
F_i|_\Lambda(z,t)
 =2z\,t^{\mathsf T}B_ic
  +t^{\mathsf T}\!\left(A_i+B_iL+L^{\mathsf T}B_i^{\mathsf T}\right)t.
\]
For every prescribed linear coefficient $\ell\in K^V$ and symmetric quadratic coefficient $H\in\Sym_V(K)$, the matrices $B_i$ satisfying $B_ic=\ell$ form an affine space of the same size, and then $A_i=H-B_iL-L^{\mathsf T}B_i^{\mathsf T}$ is forced.  Thus $(A_i,B_i)\mapsto F_i|_\Lambda$ is a surjective linear map with equal-size fibers.  The pairs $(A_i,B_i)$ are marginally independent and uniform by \cref{prop:ideal-factorization}, proving the claim.
\end{proof}

\begin{remark}[Why the statement is marginal]\label{rem:slice-law-marginal}
The model permits the secret graph $G$ to depend arbitrarily on $A$.  Conditioning on $G$ can therefore change the law of $A$, so no conditional-on-$G$ uniformity is asserted.  None is needed: the bad-key events in this section depend only on the marginal central tuple $(A_i,B_i)_i$, while a fresh public slice remains uniform after the fixed secret coordinate change for every individual key.
\end{remark}

\begin{proposition}[Generic isolation by the spare equations]\label{prop:generic-full-isolation}
Fix $p\in\PP^V$ and let $m>V$.  A Zariski-generic $m$-tuple in $H^0(\PP^V,\mathcal I_p(2))^m$ has $p$ as its only common projective zero.  Generically its full Jacobian has rank $V$ at $p$.
\end{proposition}

\begin{proof}
For each $x\ne p$, evaluation at $x$ is a nonzero linear functional on $H^0(\PP^V,\mathcal I_p(2))$.  Requiring all $m$ forms to vanish at $x$ therefore imposes codimension $m$.  The incidence with $x\in\PP^V\setminus\{p\}$ has dimension at most the parameter-space dimension plus $V-m$, strictly smaller than the parameter-space dimension.  By Chevalley's theorem, its image is constructible and its Zariski closure has dimension no larger than the incidence, so that closure is proper.  Hence a generic tuple has no common zero away from $p$.  At $p$, the first derivatives of a uniform quadric through $p$ range over the full cotangent space.  The rank-deficient $m\times V$ derivative matrices form a proper determinantal subvariety.
\end{proof}

The two propositions explain why the overdetermined sliced system generically contains only the planted point.  The main attack nevertheless solves a square core because \cref{prop:generic-full-isolation} is qualitative: it does not supply the finite-field density and solver interface needed for a quantified faster theorem.

Let $L/K$ be the finite field chosen below.  One invocation samples
\[
M\leftarrow L^{N\times(V+1)}.
\]
If $\rank M<V+1$, the trial fails.  Otherwise the attack restricts the public quadrics to the projective $V$-plane $\Lambda_M=\PP(\operatorname{im}M)$:
\[
q_i(t)=F_i(Mt),\qquad t\in\PP^V_L.
\]
It next samples $R\leftarrow L^{V\times m}$.  A rank defect again ends the trial; for full-rank $R$, the square core is
\[
g_a(t)=\sum_{i=1}^m R_{a,i}q_i(t),\qquad 1\le a\le V.
\]
Rank defects are counted as failed draws rather than conditioned away.  This keeps the entries of $(M,R)$ uniform, so the Schwartz--Zippel bound applies directly.  By \cref{lem:theta-ranks}, every rank defect already lies in $\{\Theta_F=0\}$.

In secret affine coordinates, a transverse $\Lambda_M$ is exactly a slice $o=c+Lv$ after translating its unique oil intersection.  The projective parameterization is preferable in the proof because it has no inverses or denominators: every restricted coefficient is simply quadratic in the entries of $M$.

\subsection{Why the full discriminant is enough}

The finite-field Kronecker theorem used below assumes that every affine prefix is a reduced complete intersection.  A smooth full projective core already implies these prefix conditions; no product of prefix discriminants is needed.

\begin{lemma}[A smooth full core gives reduced regular prefixes]\label{lem:full-implies-prefix}
Let $g_1,\ldots,g_V$ be homogeneous quadrics in $\PP^V$ over a field of odd characteristic.  If their common zero scheme is a smooth zero-dimensional complete intersection, then the ordered quadrics form a regular sequence and every homogeneous prefix ideal $(g_1,\ldots,g_s)$ is geometrically radical.  After any projective coordinate change and dehomogenization in an affine chart containing a final root, the affine sequence remains regular and every prefix ideal remains geometrically radical.  These are precisely the solver's hypotheses with localization polynomial equal to $1$.
\end{lemma}

\begin{proof}
Work over an algebraic closure and let $S$ be the homogeneous coordinate ring.  The final ideal has height $V$ and is generated by $V$ elements.  Since $S$ is Cohen--Macaulay, $g_1,\ldots,g_V$ is a regular sequence.  Every prefix $I_s=(g_1,\ldots,g_s)$ is therefore an unmixed complete intersection, and its coordinate ring is Cohen--Macaulay.

Suppose that some prefix were not radical.  A Cohen--Macaulay ring satisfies Serre's condition $(S_1)$, so a nonreduced prefix cannot be generically reduced: some irreducible component occurs in its complete-intersection cycle with multiplicity at least two.  Each remaining quadric is a nonzerodivisor modulo the preceding prefix and hence contains no component.  Until the final cut, every surviving component has positive projective dimension, so every remaining positive-degree hypersurface meets it.  The associativity formula for intersection multiplicities preserves the component's multiplicity factor through these proper cuts.  The final zero-cycle would therefore contain a point with multiplicity at least two, contradicting smoothness of the final scheme.  Thus every homogeneous prefix is geometrically radical.

Localization at a chart equation and dehomogenization preserve regular sequences and reducedness.  The same conclusions hold in every affine chart containing a final root; roots at infinity are simply absent from the affine zero scheme.
\end{proof}

Let $\Delta_{V,V}$ be the reduced discriminant for $V$ quadrics in $\PP^V$.  It vanishes exactly when the common zero scheme is not a smooth complete intersection of codimension $V$.  Because $\operatorname{char}K=127\ne2$, Benoist's arbitrary-characteristic formula~\cite{Benoist} gives the partial degree
\begin{equation}\label{eq:dpartial}
d_V=(V+1)2^{V-1}
\end{equation}
in the coefficients of each equation.  Thus it suffices to make the single polynomial
\begin{equation}\label{eq:theta}
\Theta_F(M,R):=
\Delta_{V,V}\!\left(\sum_iR_{1i}F_i(Mt),\ldots,
\sum_iR_{Vi}F_i(Mt)
\right)
\end{equation}
nonzero.  Each coefficient of a restricted equation has degree at most two in $M$ and one in the corresponding row of $R$.  Since the discriminant has degree $d_V$ in the coefficients of each of the $V$ equations,
\begin{equation}\label{eq:theta-degree}
\deg_{M,R}\Theta_F
\le (2V+V)d_V
=3V(V+1)2^{V-1}.
\end{equation}

\begin{lemma}[Nonvanishing certifies full parameter rank]\label{lem:theta-ranks}
If $\Theta_F(M,R)\ne0$, then $\rank M=V+1$ and $\rank R=V$.
\end{lemma}

\begin{proof}
If $\rank M<V+1$, then $\PP(\ker M)$ is nonempty and every pull-back $F_i(Mt)$, together with all of its first derivatives, vanishes there; the core is singular.  If $\rank R<V$, the output equations are linearly dependent and cannot cut a smooth codimension-$V$ complete intersection.  In either case the full discriminant vanishes.
\end{proof}

\subsection{Geometric nonvanishing of the joint slice-and-output discriminant}

For $x\in X_F$, let
\[
\rho_F(x):=\rank\bigl(dF_1(x),\ldots,dF_m(x)\bigr),
\]
where the differentials are taken on $T_x\PP(\cV)$.

\begin{definition}[Three key-level bad loci]\label{def:bad-loci}
Define:
\begin{align*}
\cB_{\oil}&:=\{\,F:\rho_F(p)<V\text{ for every }p\in\PP(\cO_0)(\overline K)\},\\
\cB_{\off}&:=\{\,F:\exists x\in X_F(\overline K)\setminus\PP(\cO_0)\text{ with }\rho_F(x)<V\},\\
\cB_{\spanbad}&:=\{\,F:\dim\Span(F_1,\ldots,F_m)\le V\}.
\end{align*}
\end{definition}

The next three lemmas separate the proof into linear-section geometry, a second-order separability argument, and generic output mixing.

\Needspace{14\baselineskip}
\begin{lemma}[A transverse finite base slice through a regular oil point]\label{lem:linear-section}
Assume $F\notin\cB_{\oil}\cup\cB_{\off}$.  Choose
$p\in\PP(\cO_0)(\overline K)$ with $\rho_F(p)=V$.  Then a Zariski-open dense set of projective $V$-planes $\Lambda\subset\PP(\cV_{\overline K})$ through $p$ has all of the following properties:
\begin{enumerate}
\item $\Lambda\cap\PP(\cO_0)=\{p\}$;
\item $X_F\cap\Lambda$ is finite;
\item at every $x\in X_F\cap\Lambda$, the restricted differential $T_x\Lambda\to\overline K^m$ has rank $V$.
\end{enumerate}
\end{lemma}

\begin{proof}
\proofstep{Dimension of the base locus.}
First, $\dim X_F\le O-1$.  Otherwise choose a closed off-oil point $x$ on a component of dimension at least $O$.  The Zariski tangent space is the common differential kernel and has dimension at least the local dimension, so
\[
\rho_F(x)=(N-1)-\dim T_xX_F\le V-1,
\]
contrary to $F\notin\cB_{\off}$.  Since $\PP(\cO_0)\subset X_F$, the maximal dimension is exactly $O-1$.

\proofstep{Behavior at the planted point.}
At $p$, every differential annihilates $T_p\PP(\cO_0)$.  Both this tangent space and $\ker dF(p)$ have dimension $O-1$, so they are equal.  Therefore a general $V$-plane through $p$ is transverse at $p$ and intersects the oil space only in $p$.  We henceforth restrict to this nonempty open subset of the Grassmannian; no other oil point can occur in the slice.

\proofstep{Off-oil points not lying on a base line.}
Let $\mathcal G_p$ be the Grassmannian of projective $V$-planes through $p$, of dimension $V(O-1)$.  We now consider only $x\in X_F\setminus\PP(\cO_0)$.  Put $K_x=\ker dF(x)$; the off-oil hypothesis gives $\dim K_x\le O-1$.  For fixed $x$, planes through $p$ and $x$ form a family of dimension $(V-1)(O-1)$.  If the tangent direction of the line $\overline{px}$ is not in $K_x$, the additional condition $T_x\Lambda\cap K_x\ne0$ is a proper Schubert condition and lowers this dimension by at least one.  Allowing $x$ to vary in dimension at most $O-1$ gives total dimension at most $V(O-1)-1$.

\proofstep{Lines through the planted point.}
If the tangent direction of $\overline{px}$ lies in $K_x$, then every quadric $F_i$ vanishes identically on that line: both endpoint values and the polar cross term vanish.  Each such direction contributes a full projective line contained in $X_F$, so passing from points of $X_F$ to line directions lowers dimension by at least one.  The direction variety therefore has dimension at most $\dim X_F-1\le O-2$.  Planes containing one such line again form a family of dimension at most
\[
(O-2)+(V-1)(O-1)=V(O-1)-1.
\]
\proofstep{Conclusion.}
The off-oil rank-defect incidence is locally closed and has dimension at most $V(O-1)-1$ by the two cases above.  Its closure in the projective product $X_F\times\mathcal G_p$ has the same dimension, and the projection of that closure is therefore a proper closed subset of $\mathcal G_p$.  Outside it, every intersection point has restricted derivative rank $V$.  At any such point the Zariski tangent space of $X_F\cap\Lambda$ is $T_x\Lambda\cap\ker dF(x)=0$.  Since local dimension never exceeds tangent-space dimension, no point can lie on a positive-dimensional component.  Hence the intersection is finite.
\end{proof}

\Needspace{11\baselineskip}
\begin{lemma}[A second-order spare equation forces separability]\label{lem:second-order}
Let $f_1,\ldots,f_m$ be quadrics on $\PP^V$ with a common point $p$.  Suppose their differential matrix at $p$ has rank $V$.  If there is a nonzero
\[
h\in\Span(f_1,\ldots,f_m)
\]
whose first jet at $p$ is zero, then the rational map
\[
\phi_f:\PP^V\dashrightarrow\PP^{m-1},
\qquad
x\longmapsto[f_1(x):\cdots:f_m(x)]
\]
has differential rank $V$ on a nonempty open subset.
\end{lemma}

\begin{proof}
Choose projective coordinates with $p=[1:0:\cdots:0]$ and affine coordinates $y=(y_1,\ldots,y_V)$.  After an equation-basis change, the first $V$ forms have expansions
\[
f_j(y)=y_j+q_j(y),
\]
where each $q_j$ is quadratic.  Since $h$ is a quadric with zero value and zero differential at $p$, it is a nonzero homogeneous quadratic $q(y)$ on this chart.

Consider the $(V+1)\times(V+1)$ first-jet minor whose rows are $f_1,\ldots,f_V,h$ and whose columns are value followed by the $V$ first derivatives.  Up to a nonzero sign, its lowest-degree homogeneous term is
\[
q(y)-\nabla q(y)\cdot y=-q(y),
\]
using Euler's identity $\nabla q(y)\cdot y=2q(y)$.  It is therefore nonzero.  Hence the full first-jet rank is $V+1$ generically, which is equivalent to differential rank $V$ for $\phi_f$ away from the base locus.
\end{proof}

\begin{lemma}[Finite regular base locus plus separability]\label{lem:finite-base-mixing}
Let $f_1,\ldots,f_m$ be quadrics on $\PP^V$ over an algebraically closed field.  Assume:
\begin{enumerate}
\item their common base scheme $B=V(f_1,\ldots,f_m)$ is finite and nonempty;
\item the $m\times V$ projective differential matrix has rank $V$ at every point of $B$;
\item $\phi_f$ has differential rank $V$ on a nonempty open subset.
\end{enumerate}
Then a nonempty Zariski-open set of matrices $R\in k^{V\times m}$ makes
\[
\left(\sum_iR_{1i}f_i,\ldots,\sum_iR_{Vi}f_i\right)
\]
a smooth zero-dimensional complete intersection of degree $2^V$.
\end{lemma}

\begin{proof}
Let $\cR=k^{V\times m}$.  For $x\notin B$, write
\[
a(x)=(f_1(x),\ldots,f_m(x))\ne0
\]
and let $J(x)$ be the differential matrix.  A matrix $R$ has $x$ as a root exactly when $Ra(x)=0$, a codimension-$V$ linear condition on $\cR$.

On the open set where the first-jet matrix $[a(x)\mid J(x)]$ has rank $V+1$, the map
\[
a(x)^\perp\longrightarrow T_x^*\PP^V,
\qquad
r\longmapsto rJ(x)
\]
is surjective.  Consequently the additional condition $\det(RJ(x))=0$ has codimension one inside $Ra(x)=0$.  After allowing $x$ to vary in a $V$-dimensional space, this singular-root incidence still has codimension at least one in $\cR$.

The first-jet rank-defect locus away from $B$ has dimension at most $V-1$ by assumption~3.  There the root condition $Ra(x)=0$ alone has codimension $V$, so its total incidence also has dimension at most $\dim\cR-1$.

Finally, at each $x\in B$, the root condition is automatic, but the rank-$V$ derivative assumption makes $\det(RJ(x))=0$ a genuine hypersurface in $\cR$.  Since $B$ is finite, these basepoint hypersurfaces do not fill $\cR$.

Globally, the singular-root incidence is the closed subset of $\PP^V\times\cR$ cut out by $Ra(x)=0$ and the maximal minors of $RJ(x)$.  Its projection is closed because $\PP^V$ is proper, and the preceding stratified dimension bounds show that this projection is proper.  Choose $R$ outside it.  The resulting projective zero scheme is nonempty because it contains $B$.  It cannot have positive dimension: at every root the Jacobian has rank $V$, so the Zariski tangent space is zero-dimensional.  Therefore the zero scheme is smooth and zero-dimensional.  Since it is cut out by $V$ quadrics in $\PP^V$, it is a complete intersection of degree $2^V$.
\end{proof}

\begin{theorem}[Existence of a smooth projective core]\label{thm:joint-nonzero}
If
\[
F\notin\cB_{\oil}\cup\cB_{\off}\cup\cB_{\spanbad},
\]
then the joint polynomial $\Theta_F(M,R)$ in \eqref{eq:theta} is nonzero.
\end{theorem}

\begin{proof}
Choose a regular geometric oil point $p\in\PP(\cO_0)(\overline K)$.  By \cref{lem:linear-section}, a dense open subset of the irreducible Grassmannian of $V$-planes through $p$ has finite base intersection and full restricted derivative rank at every basepoint.

After scalar extension, let
\[
T_{F,\overline K}:\overline K^m\longrightarrow\cU_{\overline K},\qquad
(\lambda_i)_i\longmapsto\sum_i\lambda_iF_i.
\]
The coefficient kernel of the derivative map at $p$ has dimension $m-V=e$ over $\overline K$.  Scalar extension preserves the rank of the coefficient map, so
\[
\dim_{\overline K}\ker T_{F,\overline K}
=m-\dim_K\Span_K(F_1,\ldots,F_m).
\]
This relation space lies inside the derivative kernel.  Since $F\notin\cB_{\spanbad}$, its dimension is strictly smaller than $e$.  We may therefore choose a derivative-kernel vector outside $\ker T_{F,\overline K}$; it produces a nonzero global quadric $H\in\Span_{\overline K}(F_i)$ with zero first jet at $p$.

The condition $H|_\Lambda\not\equiv0$ is open and nonempty: if $H$ vanished on every $V$-plane through $p$, it would vanish on their union, which is all of $\PP(\cV)$.  Intersecting this open set with the one supplied by \cref{lem:linear-section} gives a plane $\Lambda$ satisfying both properties.  On this plane, \cref{lem:second-order} makes the restricted linear system generically separable, and \cref{lem:finite-base-mixing} supplies an output matrix $R$ whose core is a smooth complete intersection of degree $2^V$.  Hence $\Theta_F$ is not the zero polynomial.
\end{proof}

\subsection{Finite-field density of the three bad loci}

\subsubsection{First-jet surjectivity away from oil}

\begin{lemma}[The central-form space spans every off-oil first jet]\label{lem:jet-surj}
For every $x\in\PP(\cV)\setminus\PP(\cO_0)$, the first-jet evaluation map
\[
\cU\longrightarrow K\oplus T_x^*\PP(\cV)
\]
is surjective after scalar extension to $\overline K$.
\end{lemma}

\begin{proof}
Choose a vinegar coordinate $v_j$ that is nonzero at $x$ and normalize to the chart $v_j=1$.  The quadrics
\[
v_j^2,\quad v_jv_k\ (k\ne j),\quad v_jo_a\ (1\le a\le O)
\]
belong to $\cU$ and restrict to the affine functions $1$, the remaining vinegar coordinates, and the oil coordinates.  Their first jets form a basis of the value-and-cotangent space.
\end{proof}

\Needspace{8\baselineskip}
\subsubsection{The off-oil incidence}

The codimension count has a simple structure.  At a fixed off-oil point, the values and first derivatives of each central form can be prescribed independently by \cref{lem:jet-surj}.  Making the point a common zero costs $m$ conditions, and forcing derivative rank at most $V-1$ adds the usual determinantal codimension.  Projecting away the point can recover at most its $N-1$ coordinates.  The next lemma makes this count uniform and controls the projection degree.

Let $\mathcal A=\cU^m$ be the affine key space.  For an affine variety, write $\cdeg$ for the sum of the degrees of its irreducible components.  For each vinegar chart $C_j\simeq\mathbb A^{N-1}$, let $Z_j\subset\mathcal A\times C_j$ be the incidence of pairs $(F,x)$ such that all $F_i(x)$ vanish and the $m\times(N-1)$ derivative matrix has rank at most $V-1$.

\begin{lemma}[Codimension and cumulative degree of the off-oil incidence]\label{lem:off-incidence}
Each $Z_j$ is pure of codimension
\[
m+c_{\det},\qquad c_{\det}=(m-V+1)O=(e+1)O,
\]
and satisfies
\[
\cdeg(Z_j)\le 3^m(2V)^{c_{\det}}.
\]
If $Y_j$ is the Zariski closure of its projection to key space, then
\[
\codim_{\mathcal A}Y_j\ge c_{\off}:=m+c_{\det}-(N-1)=e(O+1)+1,
\qquad
\cdeg(Y_j)\le\cdeg(Z_j).
\]
\end{lemma}

\begin{proof}
\proofstep{Structure and codimension.}
On the chart $v_j=1$, the monomials $v_j^2$, $v_jv_k$ for $k\ne j$, and $v_jo_a$ restrict to $1$, the remaining vinegar coordinates, and the oil coordinates.  They give a polynomial splitting of value-and-first-jet evaluation.  After a triangular polynomial change of coefficient coordinates, $Z_j$ is the product of an affine space, the equations setting $m$ value coordinates to zero, and the determinantal variety of $m\times(N-1)$ matrices of rank at most $V-1$.  Hence $Z_j$ is irreducible and pure of codimension
\[
m+(m-V+1)((N-1)-(V-1))=m+(m-V+1)O.
\]

\proofstep{Cumulative degree.}
In the original coefficient-and-point coordinates, every value equation has total degree at most three.  Every derivative entry has degree at most two, so each $V\times V$ minor has degree at most $2V$.  At a generic rank-$(V-1)$ point, the determinantal variety is smooth of codimension $c_{\det}$, and the differentials of its maximal minors span the conormal space.  Choose $c_{\det}$ generic linear combinations of those minors whose differentials form a basis there.  Together with the $m$ value equations, they cut a complete intersection of codimension $m+c_{\det}$ that contains the irreducible variety $Z_j$ as a component.  Cumulative B\'ezout therefore gives
\[
\cdeg(Z_j)\le 3^m(2V)^{c_{\det}}
\]
by the cumulative-degree bounds of Heintz and Lachaud--Rolland~\cite{Heintz,LachaudRolland}.

\proofstep{Projection to key space.}
Projection can recover at most the $N-1$ chart coordinates, which gives
\[
\codim_{\mathcal A}Y_j\ge m+c_{\det}-(N-1)=e(O+1)+1.
\]
The coordinate projection $\mathcal A\times C_j\to\mathcal A$ is linear.  The standard degree-under-linear-projection inequality gives
\[
\deg \overline{\pi(Z_j)}\le\deg Z_j
\]
for an irreducible affine variety~\cite{Heintz,LachaudRolland}.  Since both $Z_j$ and its image closure $Y_j$ are irreducible, this is exactly $\cdeg(Y_j)\le\cdeg(Z_j)$.
\end{proof}

\begin{proposition}[Explicit off-oil density]\label{prop:off-density}
Under the ideal-expansion distribution,
\begin{equation}\label{eq:etaoff}
\Pr[F\in\cB_{\off}]
\le
\eta_{\off}:=
V\,3^m(2V)^{(e+1)O}
Q^{-\{e(O+1)+1\}}.
\end{equation}
\end{proposition}

\begin{proof}
Every off-oil point lies in one of the $V$ vinegar charts, so $\cB_{\off}\subseteq\bigcup_jY_j$.  The affine finite-field point bound
\[
|Y_j(\F_Q)|\le\cdeg(Y_j)Q^{\dim Y_j}
\]
from Lachaud--Rolland~\cite{LachaudRolland}, followed by the union bound and \cref{lem:off-incidence}, gives the formula after division by $|\mathcal A(\F_Q)|$.
\end{proof}

\subsubsection{The oil and low-span loci}

Let
\[
\tau_{m,V}=1-\prod_{j=0}^{V-1}(1-Q^{j-m})
\]
be the rank-deficiency probability of a uniform $m\times V$ matrix.  If every geometric oil direction is rank deficient, then in particular all $O$ standard oil-coordinate directions are rank deficient.  Their derivative matrices are independent in the public-random expansion model.  Hence
\begin{equation}\label{eq:etaoil}
\Pr[F\in\cB_{\oil}]\le\eta_{\oil}:=\tau_{m,V}^{O}.
\end{equation}

Let $e=m-V$ and let $T_F:K^m\to\cU$ be the coefficient map from the proof of \cref{thm:joint-nonzero}.  The event $F\in\cB_{\spanbad}$ is exactly $\dim\ker T_F\ge e$.  For each fixed $e$-space $W\subset K^m$, the condition $W\subseteq\ker T_F$ imposes $eD_{\rm form}$ independent linear conditions on the uniform map $T_F$.  A union bound over the Grassmannian therefore gives
\begin{equation}\label{eq:etaspan}
\Pr[F\in\cB_{\spanbad}]
\le
\eta_{\rm span}:=
\genfrac{[}{]}{0pt}{}{m}{e}_Q Q^{-eD_{\rm form}}.
\end{equation}
Thus the geometric proof requires only $\dim\Span(F_1,\ldots,F_m)>V$; full linear independence is unnecessary.

Combining \cref{thm:joint-nonzero} with
\eqref{eq:etaoff}, \eqref{eq:etaoil}, and \eqref{eq:etaspan} gives the key theorem.

\Needspace{15\baselineskip}
\begin{theorem}[Almost-all-key joint regularizability]\label{thm:key-density}
In the public-random expansion model,
\begin{equation}\label{eq:keybad-total}
\Pr[\Theta_F\equiv0]
\le\eta_{\off}+\eta_{\oil}+\eta_{\rm span}.
\end{equation}
For the Round-2 parameter sets:
\begin{center}
\begin{tabular}{@{}lrrrr@{}}
\toprule
Level & $c_{\off}$ & $\log_2\eta_{\off}$ & $\log_2\eta_{\oil}$ & $\log_2\Pr[\Theta_F\equiv0]$\\
\midrule
I   & 39  & $-364.564$  & $-1132.167$ & $<-364.563$\\
III & 55  & $-457.920$  & $-1635.352$ & $<-457.919$\\
V   & 109 & $-1038.067$ & $-2935.248$ & $<-1038.066$\\
\bottomrule
\end{tabular}
\end{center}
The low-span term is below $2^{-94850}$ already at Level I and is invisible at the displayed precision.
\end{theorem}

\begin{proof}
By \cref{thm:joint-nonzero}, the event $\Theta_F\equiv0$ is contained in $\cB_{\oil}\cup\cB_{\off}\cup\cB_{\spanbad}$.  Apply \eqref{eq:etaoff}, \eqref{eq:etaoil}, \eqref{eq:etaspan}, and the union bound.  The displayed values are direct evaluations of those formulas, rounded outward at the shown precision.
\end{proof}

\section{Finite-field solver interfaces and verified restart}\label{sec:solver-attack}

We first specialize a finite-field theorem for locally closed sets to the graph-direction system.  This gives a recovery route whose correctness needs only the local Jacobian condition of \cref{sec:local-geometry}.  We then return to the faster projective route and its near-quadratic solver.  In both cases, every accepted key is checked against the unchanged public map.

\subsection{Locally closed finite-field solving}\label{subsec:local-solver}

The following proposition records the exact external interface used below.  It is a specialization of Theorem~12 of Gim\'enez et al.~\cite{Gimenez2026}.

\begin{proposition}[Locally closed finite-field solver]\label{prop:gimenez-theorem}
Let $F_1,\ldots,F_r,H\in\F_{\mathfrak q}[X_1,\ldots,X_n]$ have degree at most $d$.  Suppose that, in the localization $\F_{\mathfrak q}[X]_H$,
\begin{enumerate}
\item $F_1,\ldots,F_r$ form a regular sequence; and
\item $(F_1,\ldots,F_s)$ is radical for every $1\le s\le r$.
\end{enumerate}
Let $V_s$ be the Zariski closure of
$\{F_1=\cdots=F_s=0,\ H\ne0\}$, put $\delta_s=\deg V_s$, and let $\delta=\max_s\delta_s$.  If the input has division-free straight-line-program length $L$, $0<\epsilon<1/4$, and
\begin{equation}\label{eq:gimenez-field-condition}
 \mathfrak q>2\epsilon^{-1}n^2rd\delta^3,
\end{equation}
then a probabilistic algorithm computes a Kronecker representation of a lifting fiber of $V_r$ with probability at least $1-2\epsilon$ and bit complexity
\begin{equation}\label{eq:gimenez-cost}
 \widetilde O\!\left(
   \bigl(n^4+r(L+n^2+r^4)\bigr)d\delta^2\log_2\mathfrak q
 \right).
\end{equation}
When $r=n$, the lifting fiber is the entire zero-dimensional locally closed set.  The solver is Monte Carlo: its paper does not provide an efficient certificate that a returned representation is complete.
\end{proposition}

\begin{proof}
The two displayed algebraic conditions are hypotheses (H1)--(H2) of the cited theorem.  Equation \eqref{eq:gimenez-field-condition}, success probability $1-2\epsilon$, and \eqref{eq:gimenez-cost} are its finite-field statement.  The source explicitly notes that its algorithms do not appear to be Las Vegas because output correctness cannot in general be verified at comparable cost.  For $r=n$, no free projection coordinates remain, so a lifting fiber is the zero-dimensional set itself.  The source treats $\delta>d$ as without loss of generality; if $\delta\le d$, replacing it by the safe upper bound $d+1$ only enlarges the field and cost bounds.
\end{proof}

\begin{corollary}[Graph-local specialization]\label{cor:gimenez-local}
Assume $\rank\mathcal B(c)=V$ and $R\mathcal B(c)$ is invertible.  Apply \cref{prop:gimenez-theorem} to the $V$ affine quadrics $f=Rq_c$ and the localization polynomial $h=\det J_f$.  Then
\[
 n=r=V,\qquad d=V,\qquad \delta\le D:=2^V.
\]
The input has straight-line-program length
\begin{equation}\label{eq:local-slp}
 L_{\rm loc}=O(L_{\rm core}+V^4),
 \qquad
 L_{\rm core}:=M_V+2V(M_V+V+1),
 \quad M_V:=\frac{V(V+1)}2.
\end{equation}
Over any extension $L=\F_{Q^s}$ satisfying
\begin{equation}\label{eq:local-field-condition}
 Q^s>2\epsilon^{-1}V^4 8^V,
\end{equation}
the solver contains the planted root $Gc$ with probability at least $1-2\epsilon$ and has bit complexity
\begin{equation}\label{eq:local-solver-cost}
 \widetilde O\!\left(
 \bigl(V^4+V(L_{\rm loc}+V^2+V^4)\bigr)V4^V\log_2(Q^s)
 \right).
\end{equation}
\end{corollary}

\begin{proof}
\Cref{cor:local-hypotheses} gives the localized regularity and radicality hypotheses and places $Gc$ in the open set.  B\'ezout bounds every intermediate degree by $2^s\le2^V$.  The maximum input degree is $V$ because $h$ has degree $V$.  The dense quadratic core has the shared-monomial circuit in \eqref{eq:local-slp}; Baur--Strassen differentiation followed by a division-free Samuelson--Berkowitz determinant circuit adds $O(V^4)$ operations.  Substitution into \eqref{eq:gimenez-field-condition}--\eqref{eq:gimenez-cost} proves the result.
\end{proof}

\begin{corollary}[Smooth-projective Gim'enez cross-check]\label{cor:gimenez-projective}
On a smooth projective core, choose a hyperplane at infinity and dehomogenize in the complementary affine chart.  Conditional on the hyperplane avoiding all $D=2^V$ projective roots, apply \cref{prop:gimenez-theorem} with $H=1$.  Then
\[
 n=r=V,\qquad d=2,\qquad\delta\le D,
\]
and, over $\F_{Q^s}$ with
\begin{equation}\label{eq:projective-gimenez-field}
 Q^s>4\epsilon^{-1}V^3 8^V,
\end{equation}
the direct displayed complexity is
\begin{equation}\label{eq:projective-gimenez-cost}
 \widetilde O\!\left(
  2\bigl(V^4+V(L_{\rm core}+V^2+V^4)\bigr)4^V\log_2(Q^s)
 \right).
\end{equation}
The conditional solver success probability is at least $1-2\epsilon$.  If the hyperplane is sampled uniformly over $\F_{Q^s}$, the unconditional success probability of the chart-plus-solver step is at least
\begin{equation}\label{eq:projective-gimenez-success}
 1-2\epsilon-\frac{D}{Q^s}.
\end{equation}
\end{corollary}

\begin{proof}
Smoothness of the full projective zero scheme implies regularity and geometric radicality of every prefix by \cref{lem:full-implies-prefix}.  A hyperplane avoiding the finite zero set exists, and dehomogenization in its complementary affine chart preserves those properties.  Substitute $d=2$ and $\delta\le2^V$ into \cref{prop:gimenez-theorem}.  A uniformly random hyperplane contains any fixed projective point with probability at most $1/Q^s$; a union bound over the $D$ roots gives the additional $D/Q^s$ term in \eqref{eq:projective-gimenez-success}.
\end{proof}

\paragraph{Why Monte Carlo output is sufficient here.}
A successful call to \cref{prop:gimenez-theorem} contains the planted nonsingular root.  A failed call may omit it or return unusable data.  The attack parses every returned representation, retains only actual roots of all original equations, completes and verifies the graph, checks the inverse pull-back and public-map factorization, and finally checks an invertible oil system.  If no candidate passes, it retries with fresh coins.  Thus the external solver's lack of a completeness certificate can delay termination but cannot create a false accepted key.

\subsection{Fast projective route: field size, failure probability, and cost}\label{subsec:field}

Put
\[
 D:=2^V,\qquad L_s:=\F_{Q^s},\qquad q_s:=|L_s|.
\]
The finite-field theorem of van der Hoeven and Lecerf has an exponent parameter and a field-size requirement.  They must be chosen together.  We use $\varepsilon>0$ for the desired final exponent and set the theorem's internal parameter to $\eta=\varepsilon/3$.  For $V$ quadrics, its field-size threshold is
\begin{equation}\label{eq:vdhl-field-bound}
 B_V(\eta):=
 \max\!\left\{
 \max\!\left\{2,(1+\eta)(8/\eta+1)V2^{\eta/8}\right\}(2^V-1),
 100\,8^V
 \right\}.
\end{equation}
Define the smallest admissible extension degree and a one-degree safety choice by
\begin{equation}\label{eq:solver-extension-degree}
 s_\varepsilon^0:=\min\{s\ge1:Q^s>B_V(\varepsilon/3)\},
 \qquad
 s_\varepsilon^+:=s_\varepsilon^0+1.
\end{equation}
The first choice minimizes field size.  The second divides all finite-invocation failure bounds by the large factor $Q=2{,}048{,}383$.

For any $s\ge s_\varepsilon^0$, sample
\[
 M\leftarrow L_s^{(V+O)\times(V+1)},
 \qquad
 R\leftarrow L_s^{V\times m}
\]
and form the restricted square core.  For every fixed key with $\Theta_F\not\equiv0$, Schwartz--Zippel gives
\begin{equation}\label{eq:fast-beta}
 \beta_{V,s}:=\Pr[\Theta_F(M,R)=0]
 \le \frac{E_V^{\rm joint}}{Q^s},
 \qquad
 E_V^{\rm joint}:=3V(V+1)2^{V-1}.
\end{equation}
The proof of the finite solver theorem bounds its three internal bad choices by
\begin{equation}\label{eq:fast-alpha}
 \alpha_{V,s}:=\frac{5D^2+26D^3}{Q^s}.
\end{equation}
These are bounds for one finite invocation.  A denominator failure, degree mismatch, or failed final check returns \textsf{fail}; the algorithm never runs an unproved internal restart loop on an irregular outer draw.

For the primary choice $\varepsilon=0.3$, the field threshold is governed by $100\,8^V$.  The minimum and safety choices are:
\begin{center}
\small
\begin{tabular}{@{}lrrrr@{}}
\toprule
Level & $s_{0.3}^0$ & success at $s_{0.3}^0$ & $s_{0.3}^+$ & success at $s_{0.3}^+$ \\
\midrule
I   &8  &$\ge0.992337625602$ &9  &$\ge0.999999996259$\\
III &12 &$\ge0.999997944684$ &13 &$\ge0.999999999998$\\
V   &15 &$\ge0.927725586161$ &16 &$\ge0.999999964716$\\
\bottomrule
\end{tabular}
\end{center}
Here success means that the outer core is smooth and the solver's internal choices are good, conditional on $\Theta_F\not\equiv0$; it is at least
$(1-\beta_{V,s})(1-\alpha_{V,s})$.  At the safety fields, the base-two logarithms of $(\alpha_{V,s},\beta_{V,s})$ are
\[
(-27.994,-124.681),\qquad
(-39.858,-183.459),\qquad
(-24.756,-219.513).
\]
Thus the minimum fields are adequate for expected-time verified restart, while the safety fields give a clean one-invocation theorem.  The dependence on $\varepsilon$ in \eqref{eq:solver-extension-degree} is essential: no fixed list of extension degrees works for every arbitrarily small fixed $\varepsilon>0$.

\subsection{Published near-quadratic Kronecker solver}\label{subsec:solver}

\begin{proposition}[Fast smooth-core specialization]\label{prop:finite-field-kronecker}
Let $g_1,\ldots,g_V$ be homogeneous quadrics over a finite field $\F_{p^k}$.  Assume that they form a regular sequence and every projective prefix is geometrically radical.  Put $D_i=2^i$ and $D=D_V$.  Fix a rational $\varepsilon>0$, set $\eta=\varepsilon/3$, and assume
\[
 p^k>B_V(\eta).
\]
One finite invocation of Theorem~5.4 of van der Hoeven and Lecerf uses
\begin{equation}\label{eq:fast-kronecker}
 \widetilde O\!\left(D\,2^{(1+\varepsilon)(V-1)}k\log_2p\right)
\end{equation}
bit operations and fails because of its random coordinate, intersection, or evaluation choices with probability at most
\begin{equation}\label{eq:vdhl-internal-failure}
 \frac{5D^2+26D^3}{p^k}.
\end{equation}
Outside those bad choices it returns a complete Kronecker parametrization.  The checks in \cref{prop:resolution-checks}, including the requirement that the separator have degree $D$, certify completeness for this task.  Repeating checked finite invocations on an input satisfying the hypotheses gives the Las Vegas expected-time form of Corollary~5.5.

For $K=\F_{127^3}$ and any $s\ge s_\varepsilon^0$, the explicit finite-invocation cost over $L_s$ is obtained from \eqref{eq:fast-kronecker} with $k=3s$ and $p=127$.  Since $s=O(V)=O(\log D)$ for the choices above, this field-degree factor is absorbed by soft-$O$ notation.  Including field construction, casting, public slicing, output recombination, solving, and verification, the full expected projective route has scale
\begin{equation}\label{eq:fast-solver-scale}
 T_{\varepsilon,V}
 :=\widetilde O\!\left(
 2^V2^{(1+\varepsilon)(V-1)}3\log_2 127
 \right).
\end{equation}
\end{proposition}

\begin{proof}
The first statement is Theorem~5.4 specialized to the degree sequence $(2,\ldots,2)$, with internal parameter $\eta=\varepsilon/3$.  Its complete field condition is \eqref{eq:vdhl-field-bound}; its complexity exponent $1+3\eta$ becomes $1+\varepsilon$.  The three bad-choice bounds in the proof sum to \eqref{eq:vdhl-internal-failure}.

A smooth zero-dimensional projective core satisfies the required regularity and prefix-radicality assumptions by \cref{lem:full-implies-prefix}.  Corollary~5.5 of the source constructs a suitable extension, casts the input, verifies the returned parametrization, and obtains the same expected soft-$O$ scale.  Sampling $M$ and $R$, restricting the dense quadrics, and forming the output combinations costs only polynomially many operations in the working field and is lower order than the degree-dependent term.
\end{proof}

\begin{corollary}[One checked projective invocation]\label{cor:fast-single-invocation}
Fix $\varepsilon>0$, choose any $s\ge s_\varepsilon^0$, and fix a key with $\Theta_F\not\equiv0$.  Run one outer draw over $L_s$, one finite invocation from \cref{prop:finite-field-kronecker}, and all checks from \cref{prop:resolution-checks}; return \textsf{fail} if any check rejects.  The invocation has soft-$O$ cost $T_{\varepsilon,V}$ and succeeds with probability at least
\begin{equation}\label{eq:fast-single-invocation-success}
 (1-\beta_{V,s})(1-\alpha_{V,s}).
\end{equation}
For $\varepsilon=0.3$, the minimum-field and safety-field probabilities are those in the preceding table.
\end{corollary}

\begin{proof}
The outer discriminant fails with probability at most $\beta_{V,s}$.  Conditional on a smooth core, \cref{prop:finite-field-kronecker} fails internally with probability at most $\alpha_{V,s}$.  The two groups of coins are fresh.  Exact checking changes malformed or incomplete output into \textsf{fail}, never into a false accepted result.
\end{proof}

\subsection{A classical arithmetic cross-check}\label{subsec:slp-solver}

Let
\begin{equation}\label{eq:slp-length}
 M_2:=\binom{V+2}{2},
 \qquad
 L_{\rm quad}:=M_2+V(2M_2-1).
\end{equation}
The first term computes every degree-two monomial in the $V+1$ homogeneous variables once.  Each of the $V$ dense quadrics then uses $M_2$ scalar multiplications and $M_2-1$ additions.  Thus $L_{\rm quad}=\Theta(V^3)$ is an explicit division-free circuit for the full square core.

Giusti--Lecerf--Salvy give the older arithmetic bound
\begin{equation}\label{eq:classical-slp-scale}
 L_{\rm quad}\,\widetilde O(D^2)
\end{equation}
field operations under the same regularity conditions~\cite{GiustiLecerfSalvy2001,vdHL}.  We retain \eqref{eq:classical-slp-scale} only as an epsilon-free growth-rate cross-check.  We do \emph{not} transfer the explicit finite-field failure constants of Theorem~5.4 to this older bound: the literature cited here does not state that combined theorem.  Accordingly, the SLP line below has no proved one-invocation failure probability and is not used for the universal-forgery corollary.

\Needspace{16\baselineskip}
\begin{proposition}[Task-specific resolution checks]\label{prop:resolution-checks}
Let $f=(f_1,\ldots,f_V)$ be affine quadrics over a finite field $L$, and put $D=2^V$.  Suppose a proposed geometric resolution consists of a monic polynomial $m(T)$, coordinate polynomials $v_1(T),\ldots,v_V(T)$, and a linear form $u$.  Reject it unless every $v_j$ has degree below $\deg m$ and
\begin{enumerate}
\item $1\le\deg m\le D$ and $m$ is squarefree;
\item $u(v_1(T),\ldots,v_V(T))\equiv T\pmod m$;
\item $f_i(v_1(T),\ldots,v_V(T))\equiv0\pmod m$ for every $i$;
\item $\gcd\!\left(m(T),\det J_f(v_1(T),\ldots,v_V(T))\right)=1$.
\end{enumerate}
If the solver returns standard numerators $w_j$ with denominator $m'$, first verify $\gcd(m,m')=1$ and set $v_j=w_j(m')^{-1}\bmod m$.  Passing these checks certifies that the represented points are distinct simple solutions of the square core.  When this proposition is used with a smooth zero-dimensional complete intersection of $V$ quadrics, requiring in addition $\deg m=D$ certifies completeness by B\'ezout.  The checks cost $\widetilde O(\operatorname{poly}(V)D\log|L|)$ Boolean gates.
\end{proposition}

\begin{proof}
Squarefreeness gives distinct roots $\tau$ of $m$ over $\overline L$.  The primitive-form identity makes the represented points $P_\tau=(v_1(\tau),\ldots,v_V(\tau))$ distinct.  The equation checks put them in $V(f)$, and the Jacobian gcd makes each one simple.  Evaluate the quadrics and Jacobian in $L[T]/(m)$; compute the determinant there with a division-free algorithm because this quotient need not be a field.  Polynomial arithmetic, gcds, and the $V\times V$ determinant have the stated cost.
\end{proof}

\paragraph{Why the extra checks remain useful.}
The published Las Vegas solver already verifies its complete parametrization.  The task-specific checks make the cryptanalytic interface self-contained and protect against malformed implementations, truncated outputs, and conversion errors.  On every branch, final acceptance additionally requires all unrecombined equations, exact public-map factorization, and an invertible oil system.  No malformed solver transcript can create a false accepted key.

\subsection{Root selection and oil-space recovery}\label{subsec:root-recovery}

\begin{lemma}[Recovery from the projective oil point]\label{lem:descent}
Let $L/K$ be a finite extension.  Suppose $\Theta_F(M,R)\ne0$, and let $[t^\star]\in\PP^V(L)$ be the unique point for which $x^\star=Mt^\star$ lies in $\PP(\cO)_L$.  Then
\[
\rank(P_1x^\star\ \cdots\ P_mx^\star)=V,
\qquad
\ker(P_1x^\star\ \cdots\ P_mx^\star)^{\mathsf T}=\cO_L.
\]
Row reduction recovers the public graph of the oil space.  For the planted point that graph is defined over $K$ and, after the specification's inverse pull-back, yields the planted expanded UOV decomposition.
\end{lemma}

\begin{proof}
The full-rank matrix $M$ defines a projective $V$-plane over $L$.  Its vector-space intersection with the $O$-dimensional oil space is nonzero by dimension count.  If that intersection had dimension at least two, all restricted quadrics would vanish on a positive-dimensional projective linear space, contradicting smooth zero-dimensionality.  Thus it is one $L$-rational point.

Smoothness of the square core gives derivative rank at least $V$ there.  Total isotropy puts the $O$-dimensional oil space inside the ambient common derivative kernel, so the rank is at most $(V+O)-O=V$.  Equality holds, and \cref{lem:kernel} identifies the kernel with $\cO_L$.  The remaining statements follow from $K$-descent and \cref{prop:verified-key}.
\end{proof}

\begin{lemma}[Separator descent]\label{lem:separator-descent}
Let a reduced zero-dimensional scheme over $L=\F_{\mathfrak q}$ have a Kronecker representation whose primitive linear form $u$ separates its geometric points.  A geometric point $P$ is $L$-rational if and only if $u(P)\in L$.
\end{lemma}

\begin{proof}
The forward implication is immediate.  Conversely, if $u(P)\in L$, then
\[
u(P^{(\mathfrak q)})=u(P)^{\mathfrak q}=u(P).
\]
Because $u$ is defined over $L$ and separates the geometric points, Frobenius must fix $P$.
\end{proof}

\Needspace{15\baselineskip}
\begin{proposition}[Filter all original equations before factoring]\label{prop:separator-prefilter}
Let a consistency-checked resolution over $L_s$ have squarefree separator polynomial $m_u(T)$ and coordinate polynomials $v(T)$.  Let $\widetilde q_1,\ldots,\widetilde q_m$ be the dehomogenized original restricted quadrics in the same affine chart, and define
\begin{equation}\label{eq:all-equation-factor}
a(T):=\gcd\bigl(m_u(T),\widetilde q_1(v(T)),\ldots,\widetilde q_m(v(T))\bigr).
\end{equation}
Then the roots of $a$ are exactly the separator values of represented points satisfying all $m$ original equations.  Moreover,
\begin{equation}\label{eq:rational-factor}
a_{L_s}(T):=\gcd\bigl(a(T),T^{q_s}-T\bigr)
\end{equation}
has exactly the separator values of those common points that are $L_s$-rational.  Computing $a$ and $a_{L_s}$ costs $\softO\!\left(D\,\operatorname{poly}(V,m,\log|L_s|)\right)$ Boolean gates.  Only $a_{L_s}$, whose degree may be much smaller than $D$, must be factored.
\end{proposition}

\begin{proof}
Work over an algebraic closure of $L_s$.  Because $m_u$ is squarefree, it is the product of the distinct factors $T-\tau$ indexed by the represented points.  For each root $\tau$,
\[
\widetilde q_i(v(T))\bmod m_u\;\big|_{T=\tau}=\widetilde q_i(v(\tau)).
\]
Thus $T-\tau$ divides \eqref{eq:all-equation-factor} exactly when every original restricted equation vanishes at the represented point.  The second claim follows from \cref{lem:separator-descent} and the identity $z^{q_s}=z$ for $z\in L_s$.  Dense quadratic evaluation in $L_s[T]/(m_u)$, modular exponentiation, and polynomial gcd computation at degree at most $D$ give the stated cost.  Since $\log|L_s|=O(V)$ for the selected fields, it is lower order than the degree-squared solver bound.
\end{proof}

\paragraph{Check the resolution, prefilter, and extract rational roots.}
Convert the solver output to canonical coordinate polynomials and run the consistency checks of \cref{prop:resolution-checks}.  On failure, discard the invocation.  On a successful Las Vegas solver call, the returned separator polynomial $m_u(T)\in L_s[T]$ represents the full affine zero scheme and contains the planted oil point.  The additional checks reject malformed or incomplete implementation outputs.

Apply \cref{prop:separator-prefilter}: evaluate all $m$ original restricted equations in $L_s[T]/(m_u)$, compute $a(T)$, and then compute $a_{L_s}(T)$.  Factor only $a_{L_s}$ over $L_s$.  A Las Vegas finite-field factorization algorithm terminates with probability one and has expected subquadratic dependence on the input degree, uniformly for inputs of degree at most $D$; its expected contribution is therefore below the solver's near-quadratic-in-$D$ scale~\cite{KaltofenShoup}.  Reconstruct the corresponding affine points, undo the coordinate changes, and recheck all $m$ original equations as a final exact verification.

\paragraph{Certify the decomposition.}
For each survivor $x$, compute the common derivative kernel.  Require an $O$-dimensional graph defined over $K$, total isotropy, and every identity \eqref{eq:graphid}.  Apply the specification's inverse pull-back and verify the base-field factorization against the original public map.  An output is called an \emph{exact candidate} only after all these checks pass.

\paragraph{Certify signing or restart.}
Deduplicate exact candidates by their $K$-rational graph.  For each remaining candidate, first test the $V$ basis assignments from \cref{prop:signing-search}.  If they all fail, scan the remaining nonzero combinations on the designated pair lines.  Return the first candidate with an invertible oil matrix, together with that matrix's inverse.  If no candidate passes, discard the trial and sample a fresh slice.  On a correct solver branch the candidate list contains the planted decomposition, so outside the event in \eqref{eq:eta-sign} this scan succeeds deterministically.  At most
\[
 D\bigl(V+\lfloor V/2\rfloor(q-1)\bigr)
\]
oil matrices are tested per trial, below the solver's near-quadratic-in-$D$ scale.  In the public-random model, the basis-first part succeeds after fewer than $1.008$ tests on average.

Let $\mathsf A_{\varepsilon,V,s}$ denote the event that the full trial accepts a publicly verified signing decomposition and inverse oil matrix.  On a key outside $\cE_{\rm sign}$, a draw with $\Theta_F(M,R)\ne0$ yields the planted exact candidate whenever the finite solver invocation succeeds, after which the deterministic signing scan succeeds.  The outer and internal coins are fresh, so
\begin{equation}\label{eq:acceptance-lower-bound}
\Pr[\mathsf A_{\varepsilon,V,s}]
\ge(1-\beta_{V,s})(1-\alpha_{V,s}).
\end{equation}
This one-sided implication is the only interface needed between the solver guarantee and the verified attack.  On irregular inputs or bad internal choices, a failed check returns \textsf{fail}; it never creates a false accepted key.

\begin{lemma}[Verified-restart accounting]\label{lem:restart-accounting}
Suppose each fresh trial, conditional on any history of earlier failures, accepts with probability at least $p>0$.  Then the number $N$ of trials satisfies
\[
\Pr[N>t]\le(1-p)^t,
\qquad
\mathbb E[N]\le p^{-1}.
\]
If the current trial uses fresh coins and has expected work at most $T$, uniformly over prior histories, then the expected total work is at most $T/p$.
\end{lemma}

\begin{proof}
The tail bound follows inductively from the conditional success lower bound, and summing the tail probabilities gives $\mathbb E[N]\le1/p$.  For the work bound, the event that trial $j$ is reached depends only on earlier trials.  Fresh current-trial coins therefore give
\[
\mathbb E\!\left[\mathbf 1_{\{N\ge j\}}W_j\right]
\le T\Pr[N\ge j].
\]
Summing over $j$ and applying the trial bound proves the claim.
\end{proof}

\Needspace{15\baselineskip}
\subsection{Main recovery theorem}\label{subsec:main-theorem}
Use $T_{\varepsilon,V}$ from \eqref{eq:fast-solver-scale}.  For $0<\epsilon<1/4$ and an extension degree $s$ satisfying \eqref{eq:local-field-condition}, write
\[
 T_{\rm loc}(\epsilon,s):=
 \widetilde O\!\left(
 \bigl(V^4+V(L_{\rm loc}+V^2+V^4)\bigr)V4^V\log_2(Q^s)
 \right).
\]

\begin{theorem}[Key-by-key recovery from one graph direction]\label{thm:graph-attack}
Let an expanded \QRUOV{} key over $K$ have a valid planted decomposition.  Suppose that at least one maximal minor of $\mathcal B(c)$ is a nonzero polynomial in $c$.  Fix $0<\epsilon<1/4$ and choose $s$ so that \eqref{eq:local-field-condition} holds.  Put
\[
 p_{\rm dir}:=\left(1-\frac VQ\right)
 \left(1-\frac{V(m-V)}Q\right)(1-2\epsilon).
\]
Then the following verified attacks are correct for this fixed key.
\begin{enumerate}
\item A randomized graph-direction trial finds an exactly verified expanded UOV decomposition with probability at least $p_{\rm dir}$.  Repeating fresh trials terminates almost surely, uses at most $p_{\rm dir}^{-1}$ trials on average, and has expected work
\begin{equation}\label{eq:local-expected-decomposition}
 O\!\left(T_{\rm loc}(\epsilon,s)/p_{\rm dir}\right).
\end{equation}
Every accepted output is correct.

\item If the determinant polynomial $\Delta$ of the planted decomposition is nonzero, testing one fresh uniform scalar-block vinegar assignment in each trial returns a signing certificate with expected trial count at most
\begin{equation}\label{eq:local-functional-trials}
 \bigl(p_{\rm dir}(1-m/q)\bigr)^{-1}.
\end{equation}

\item If one planted basis matrix is invertible, the deterministic basis-first signing scan succeeds on every successful recovery trial.  Under the designated-pair condition of \cref{prop:signing-search}, the full basis-and-line scan does so after at most $V+\lfloor V/2\rfloor(q-1)$ oil-matrix tests per candidate.  In either case the expected trial and work bounds are those of item~1.

\item If a coordinate direction $e_j$ satisfies $\rank\mathcal B(e_j)=V$, all outer attack choices can be made deterministic: scan the $O\,[V(m-V)+1]$ direction/condenser pairs of \cref{cor:deterministic-outer-list}, and repeat the finitely many probabilistic solver calls in rounds.  This still terminates almost surely and adds only a polynomial factor.  The shorter lists of \cref{cor:compact-direction-list} give category-scale failure bounds in the public-random model.
\end{enumerate}
The theorem uses only a local nonsingularity condition at the planted root.  It does not require the projective discriminant, the off-oil incidence argument, a key-generation distribution, or uniqueness of the recovered decomposition.
\end{theorem}

\begin{proof}
\proofstep{A good trial contains the planted root.}
By \cref{prop:graph-direction-density}, a random $c$ makes $\mathcal B(c)$ full rank with probability at least $1-V/Q$.  Conditional on that event, a random $t$ makes $\mathsf C(t)\mathcal B(c)$ invertible with probability at least $1-V(m-V)/Q$.  The planted root then lies in the localized smooth set by \cref{cor:local-hypotheses}.  By \cref{cor:gimenez-local}, the solver returns a representation containing it with probability at least $1-2\epsilon$.

\proofstep{One planted root gives the whole hidden graph.}
For every returned root $x$, solve the public systems \eqref{eq:direction-completion-system}.  At $x=Gc$, \cref{prop:graph-direction-completion} returns every column of $G$.  The attack then checks all original direction equations, rank conditions, graph identities, inverse pull-back identities, and public-map factorizations.  Any root that does not yield a valid graph is discarded.

\proofstep{Restart and signing.}
A fresh trial therefore reaches the planted expanded decomposition with probability at least $p_{\rm dir}$.  Apply \cref{lem:restart-accounting}.  The randomized and deterministic signing claims follow from \cref{prop:signing-search}.  Filtering and completing at most $D$ roots costs $\widetilde O(D\operatorname{poly}(V,O,m,\log|L|))$, below \eqref{eq:local-solver-cost}.  Exact verification prevents false acceptance even when the internal Monte Carlo solver fails.

For item~4, \cref{cor:deterministic-outer-list} guarantees that a regular coordinate direction has a condenser value with nonsingular planted Jacobian.  The deterministic scan reaches it.  Repeating all finitely many solver calls in rounds gives a fixed positive success probability per round and hence almost-sure termination.
\end{proof}

\begin{corollary}[Public-random density and compact deterministic certificates]\label{cor:graph-ideal-recovery}
In the public-random expansion model, the probability that all $O$ coordinate directions are rank deficient is exactly $\tau_{m,V}(Q)^O$, whose base-two logarithms are
\[
 -1132.166907,\qquad-1635.352198,\qquad-2935.247544.
\]
Hence, except with probability at most
\begin{equation}\label{eq:eta-local-functional}
 \eta_{\rm local}:=\tau_{m,V}(Q)^O+\eta_{\rm sign},
\end{equation}
the full deterministic direction list and deterministic signing fallback recover a universal signing certificate.

A much shorter fixed attack already has category-scale coverage.  Use the first $3,4,4$ coordinate directions and their condenser lists, together with the first $22,31,40$ basis signing assignments, at Levels I, III, and V.  These attacks solve only $315,612,1228$ square systems and test at most $22,31,40$ oil matrices per exact candidate.  Their failure probabilities are at most
\[
 \tau_{m,V}(Q)^{r_{\rm dir}}+\sigma_m^{b_{\rm sign}},
\]
with base-two logarithms below
\begin{equation}\label{eq:compact-combined-logs}
 -153.502,\qquad-216.298,\qquad-279.095.
\end{equation}
No independence between the recovery and signing failures is needed for this union bound.
\end{corollary}

\begin{proof}
The all-direction statement is \eqref{eq:all-coordinate-directions-bad}, followed by \cref{prop:signing-search} and the union bound.  For the compact statement, combine \cref{cor:compact-direction-list} with the fixed-prefix signing bound \eqref{eq:compact-sign-logs}.  On the complement, the deterministic outer list contains a square system with the planted point nonsingular, and the fixed signing list contains an invertible planted oil matrix.  Round-robin solver repetition and exact verification then give the claimed universal certificate.
\end{proof}

\begin{theorem}[Fast projective key-by-key recovery]\label{thm:attack}
Let an expanded \QRUOV{} key over $K$ have a valid planted expanded decomposition, and assume that the public joint slice-and-output discriminant $\Theta_F(M,R)$ is not the zero polynomial.  Fix a rational $\varepsilon>0$, choose any $s\ge s_\varepsilon^0$, and put
\[
 p_{\varepsilon,V,s}:=(1-\beta_{V,s})(1-\alpha_{V,s}).
\]
Then:
\begin{enumerate}
\item A decomposition-only verified-restart algorithm terminates with probability one and returns an exactly verified expanded UOV decomposition.  Every accepted output is correct, the expected number of outer trials is at most $p_{\varepsilon,V,s}^{-1}$, and the expected work is
\begin{equation}\label{eq:expected-decomposition}
 O\!\left(\frac{T_{\varepsilon,V}}{p_{\varepsilon,V,s}}\right).
\end{equation}
For $\varepsilon=0.3$, the minimum fields are $s=8,12,15$.  Their expected-trial bounds are below
\[
1.00773,\qquad1.000003,\qquad1.07791.
\]
The safety fields $s=9,13,16$ give the near-one single-invocation probabilities in \cref{cor:fast-single-invocation}.

\item Suppose the planted determinant polynomial $\Delta$ from \cref{prop:signing-search} is nonzero.  If each outer trial tests one fresh uniform scalar-block vinegar assignment for every exact candidate, verified restart returns an equivalent signing key and inverse oil matrix.  The expected number of outer trials is at most
\[
 \frac{1}{p_{\varepsilon,V,s}(1-m/q)},
\]
and the expected work is
\begin{equation}\label{eq:expected-attack-randomized}
 O\!\left(\frac{T_{\varepsilon,V}}{p_{\varepsilon,V,s}(1-m/q)}\right).
\end{equation}
At most $D$ oil matrices are tested per trial.

\item If one planted basis matrix $M_j$ is invertible, the basis-first deterministic scan returns an equivalent signing key on every successful core invocation and tests at most $DV$ oil matrices per outer trial.  Under the designated-pair condition of \cref{prop:signing-search}, the full deterministic fallback tests at most
\[
 D\bigl(V+\lfloor V/2\rfloor(q-1)\bigr)
\]
oil matrices per trial.  In either case the expected trial count is at most $p_{\varepsilon,V,s}^{-1}$ and the expected work remains $O(T_{\varepsilon,V}/p_{\varepsilon,V,s})$.
\end{enumerate}
These are statements about each fixed expanded key satisfying the displayed conditions.  They do not depend on how the key was generated.  Recovering the original seed or serialized private key is unnecessary.
\end{theorem}

\begin{proof}
\proofstep{A regular draw yields a complete core resolution.}
The discriminant hypothesis and \eqref{eq:fast-beta} imply that a fresh $(M,R)$ gives a smooth full core with probability at least $1-\beta_{V,s}$.  \Cref{lem:full-implies-prefix} supplies the solver hypotheses.  Conditional on a smooth core, the finite invocation in \cref{prop:finite-field-kronecker} returns and verifies the complete geometric resolution with probability at least $1-\alpha_{V,s}$.  Every invocation has a finite operation schedule and returns \textsf{fail} if a denominator, degree, or verification check rejects.  An irregular outer draw therefore cannot trap the attack inside an unjustified restart loop.

\proofstep{The planted decomposition appears among the exact candidates.}
On a regular draw, the quotient-algebra prefilter of \cref{prop:separator-prefilter}, rational-root extraction, and \cref{lem:descent} retain the planted oil point and recover the planted expanded decomposition.  Exact all-equation checks, inverse pull-back, and public-map factorization reject every invalid candidate.

\proofstep{Signing certification.}
If $\Delta\not\equiv0$, \cref{prop:signing-search} gives success probability at least $1-m/q$ for a fresh uniform assignment on the planted candidate.  If a basis matrix or designated pair witness exists, the corresponding deterministic scan succeeds on every regular draw.

\proofstep{Restart and auxiliary work.}
Apply \cref{lem:restart-accounting}.  Finite-field factorization has subquadratic expected dependence on degree, and all filtering, derivative-kernel, descent, pull-back, verification, and signing work costs $\widetilde O(D\operatorname{poly}(V,O,m,q,\log|L_s|))$.  For every fixed $\varepsilon>0$ this is lower order than $T_{\varepsilon,V}$.  Exact verification prevents false acceptance on every branch.
\end{proof}

\begin{corollary}[Public-random density for the fast projective route]\label{cor:ideal-recovery}
In the public-random expansion model, the probability that either $\Theta_F\equiv0$ or the planted determinant polynomial $\Delta$ is identically zero is below
\[
2^{-364.563},\qquad2^{-457.919},\qquad2^{-1038.066}
\]
at Levels I, III, and V.  Every other key satisfies the randomized equivalent-key conclusion of \cref{thm:attack}.  In fact, outside the same bounded union of bad events, one designated pair span contains an invertible matrix, so the deterministic refinement applies.
\end{corollary}

\begin{proof}
The event $\Theta_F\equiv0$ is bounded by \cref{thm:key-density}.  If $\Delta$ is the zero polynomial, every designated pair span is singular, so this event is contained in $\cE_{\rm sign}$ and is bounded by \cref{prop:signing-search}.  The union has probability at most
\[
\eta_{\off}+\eta_{\oil}+\eta_{\rm span}+\eta_{\rm sign},
\]
which gives the displayed numerical bounds.  On the complement of this particular union, the designated-pair condition itself holds, so item~3 of \cref{thm:attack} applies.
\end{proof}

\begin{corollary}[One-invocation key-only universal forger]\label{cor:single-invocation-forger}
In the public-random expansion model, set $\varepsilon=0.3$, use the safety field $s=s_{0.3}^+$, run the checked projective invocation of \cref{cor:fast-single-invocation}, and then run the basis-first scan and projective-line fallback of \cref{prop:signing-search}.  The adversary uses no signing oracle, never outputs an invalid signature, and can forge every chosen message whenever it succeeds.  Its success probability over the key and attack coins is at least
\begin{equation}\label{eq:one-pass-forger-success}
 1-\delta_{\rm key}-\alpha_{V,s_{0.3}^+}-\beta_{V,s_{0.3}^+},
\end{equation}
where $\delta_{\rm key}$ is the key-level exceptional probability in \cref{cor:ideal-recovery}.  In particular, the success probabilities at Levels I, III, and V are respectively at least
\[
 0.999999996259,\qquad0.999999999998,\qquad0.999999964716.
\]
\end{corollary}

\begin{proof}
Outside the key-level exceptional set, the projective discriminant and a deterministic pair-line witness both exist.  A successful checked solver pass returns the complete core, so the planted point survives exact filtering and yields the planted expanded decomposition.  The deterministic signing scan then returns a certificate that inverts every public target.  Apply \cref{cor:universal-forgery} and union-bound the key, outer-discriminant, and internal-solver failures.  Exact verification ensures that failure affects only whether a forgery is returned, never whether a returned forgery is valid.
\end{proof}




\section{Official seeded expansion: ideal primitives and fixed functions}\label{app:seeded-transfer}

The key-by-key theorem needs no expansion model.  This appendix asks how often official seeded key generation satisfies the sufficient algebraic conditions for deterministic equivalent-key recovery.  The official expansion uses public counter $2i-1$ for the symmetric block $A_i$, public counter $2i$ for the rectangular block $E_i$, and secret counter $1$ for the graph $G$~\cite[Secs.~4.7--4.10]{QRUOVSpec}.  For every recommended set, the secret buffer length $\tau_2$ is at most the public $A_i$ buffer length $\tau_1$.  Hence, if the independently sampled public and secret seeds happen to coincide, the secret counter-$1$ raw stream is a prefix of the public $A_1$ stream.  The resulting correlation between $G$ and $A_1$ is expressly allowed by \cref{def:ideal}; it is not a bad event.

For Levels I, III, and V, let $\lambda\in\{128,192,256\}$ be the seed length.  The specification uses $m$ public buffers of length $\tau_1$ bytes and $m$ public buffers of length $\tau_2$ bytes.  Each public rejection sampler exhausts its prescribed buffer with probability at most $2^{-\lambda}$.  Put
\begin{equation}\label{eq:public-rejection}
\varepsilon_{\rm pub}:=2m2^{-\lambda}.
\end{equation}
No secret-buffer exhaustion term is needed: \cref{def:ideal} permits an arbitrary conditional law for $G$.

\subsection{SHAKE in the ideal-XOF model}

\begin{proposition}[Exact public-random transfer in the ideal-XOF model]\label{prop:rom-transfer}
Model SHAKE as a random function from each domain-separated input to one consistent infinite output string.  The total-variation distance between the official expanded-key distribution and the public-random model of \cref{def:ideal} is at most
\[
\delta_{\rm XOF}:=\varepsilon_{\rm pub}=2m2^{-\lambda}.
\]
This statement includes keys with equal public and secret seeds.
\end{proposition}

\begin{proof}
The $2m$ public seed/counter inputs are pairwise distinct, so their ideal-XOF strings are independent and uniform.  If the seeds differ, the secret string is independent of all public strings.  If they coincide, the secret counter-$1$ string is the required prefix of the public $A_1$ string; after rejection sampling, this induces an arbitrary conditional distribution of $G$ given $A_1$, exactly as permitted by \cref{def:ideal}.  In either case, the even-counter $E_i$ strings are independent uniform strings conditioned on $(A,G)$.  Abort only if one of the $2m$ public buffers contains too few accepted values.  The union bound \eqref{eq:public-rejection} pays for that event.  Conditioned on no public exhaustion, the accepted $A_i$ values have their product-uniform marginal and the accepted $E_i$ values are product-uniform conditioned on $(A,G)$, so the distribution is exactly \cref{def:ideal}.
\end{proof}

\subsection{AES in the ideal-cipher model}

Let
\[
T:=m\left(\left\lceil\frac{\tau_1}{16}\right\rceil
          +\left\lceil\frac{\tau_2}{16}\right\rceil\right)
\]
be the number of $128$-bit public cipher blocks, and define
\begin{equation}\label{eq:ideal-cipher-distance}
\delta_{\rm perm}(T)
:=1-\prod_{j=0}^{T-1}\left(1-\frac{j}{2^{128}}\right),
\qquad
\delta_{\rm IC}:=\delta_{\rm perm}(T)+\varepsilon_{\rm pub}.
\end{equation}

\begin{proposition}[Exact public-random transfer in the ideal-cipher model]\label{prop:icm-transfer}
Model AES as an ideal cipher with an independent random permutation for every key.  The total-variation distance between the official expanded-key distribution and the public-random model of \cref{def:ideal} is at most $\delta_{\rm IC}$.  This statement also includes equal public and secret AES keys.
\end{proposition}

\begin{proof}
For every fixed public key, the specification's $T$ public counter blocks are distinct.  Their ideal-cipher outputs are therefore a uniform ordered sample without replacement from $2^{128}$ blocks.  The exact total-variation distance from $T$ independent uniform blocks is the collision probability of the independent sample, namely $\delta_{\rm perm}(T)$.

Couple the without-replacement public block vector to an independent-block vector at this distance.  If the public and secret keys differ, generate $G$ from the independent secret-key permutation in both experiments.  If the keys coincide, generate $G$ from the prescribed prefix of the coupled public $A_1$ block vector in each experiment.  In both cases this is a common randomized postprocessing kernel, so total variation cannot increase.  Under the independent-block vector, $A$ has the product-uniform marginal and $E$ is product-uniform conditioned on $(A,G)$.  Truncation and rejection sampling are further postprocessing.  Finally, public-buffer exhaustion contributes at most \eqref{eq:public-rejection}.
\end{proof}

\begin{table}[t]
\centering
\small
\setlength{\tabcolsep}{4.3pt}
\caption{Official public-buffer parameters and exact ideal-primitive transfer bounds.  The last column gives base-two logarithms.  No seed-collision term is present.}
\label{tab:seeded-model-bounds}
\begin{tabular}{@{}c r r r r r r@{}}
\toprule
Level & $\lambda$ & $m$ & $\tau_1$ & $\tau_2$ & $T$ & $\log_2\delta_{\rm XOF}\ \big/\ \log_2\delta_{\rm IC}$\\
\midrule
I   &128&54 &4267 &2916 &24300  &$-121.245\ /\ -99.863$\\
III &192&78 &9020 &6123 &73866  &$-184.715\ /\ -96.655$\\
V   &256&105&16144&11018&178290 &$-248.286\ /\ -94.112$\\
\bottomrule
\end{tabular}
\end{table}


\subsection{Independent-graph transfer for direct forgery}

The direct residual-root theorem uses \cref{def:direct-ideal}, so equal public
and secret seeds or AES keys must now be paid for rather than absorbed as an
allowed graph--coefficient correlation.  Define
\begin{equation}\label{eq:direct-transfer-distances}
 \delta_{\rm XOF}^{\rm ind}:=\delta_{\rm XOF}+2^{-\lambda},
 \qquad
 \delta_{\rm IC}^{\rm ind}:=\delta_{\rm IC}+2^{-\lambda}.
\end{equation}

\begin{proposition}[Independent-graph transfer in the ideal-primitive models]\label{prop:direct-seeded-transfer}
The total-variation distance between official expanded-key generation and
\cref{def:direct-ideal} is at most $\delta_{\rm XOF}^{\rm ind}$ when SHAKE is
modeled as an ideal XOF, and at most $\delta_{\rm IC}^{\rm ind}$ when AES is
modeled as an ideal cipher.
\end{proposition}

\begin{proof}
Condition first on the independently sampled public and secret seeds (or AES
keys) being distinct.  This event fails with probability exactly $2^{-\lambda}$.
On the conditioning event, the ideal primitive supplies an independent secret
stream for $G$ and independent public streams for $A$ and $E$.  The public
buffer and, in the ideal-cipher model, without-replacement couplings are
exactly those of \cref{prop:rom-transfer,prop:icm-transfer}.  Their losses are
$\delta_{\rm XOF}$ and $\delta_{\rm IC}$, respectively.  A union bound with the
seed-collision event gives \eqref{eq:direct-transfer-distances}.
\end{proof}

For Levels I, III, and V, the corresponding base-two logarithms are
\[
 \log_2\delta_{\rm XOF}^{\rm ind}
   =-121.232,-184.705,-248.279,
 \qquad
 \log_2\delta_{\rm IC}^{\rm ind}
   =-99.863,-96.655,-94.112,
\]
rounded upward in the last displayed digit.

\begin{corollary}[Official seeded direct-forgery invocation in ideal-primitive models]\label{cor:direct-seeded-ideal}
Let $p_{\rm dir}$ denote the right-hand side of the direct-forgery success bound
in \eqref{eq:direct-total-success}.  Under official seeded key generation, the
success probability of the checked direct-forgery invocation is at least
$p_{\rm dir}-\delta_{\rm XOF}^{\rm ind}$ in the ideal-XOF model and at least
$p_{\rm dir}-\delta_{\rm IC}^{\rm ind}$ in the ideal-cipher model.  Every
returned signature remains valid against the unchanged verifier.
\end{corollary}

\begin{proof}
Total variation changes the probability of the ideal-model success event by at
most the corresponding distance in \cref{prop:direct-seeded-transfer}.  Public
verification is distribution-free and rejects every invalid candidate.
\end{proof}

\subsection{From key-distribution distance to verified recovery}

Let
\[
 \eta_{\rm fast}:=\eta_{\off}+\eta_{\oil}+\eta_{\rm span}+\eta_{\rm sign},
 \qquad
 \eta_{\rm local}:=\tau_{m,V}(Q)^O+\eta_{\rm sign}
\]
be the public-random exceptional-key probabilities for the fast projective and graph-direction equivalent-key theorems.  Total variation is applied only to these key-level events, not to a bounded attack execution.

\begin{corollary}[Official seeded recovery in ideal-primitive models]\label{cor:seeded-ideal-primitives}
Under the official seeded key generation:
\begin{enumerate}
\item with SHAKE modeled as an ideal XOF, the probability that a key fails the graph-direction equivalent-key hypotheses is at most $\eta_{\rm local}+\delta_{\rm XOF}$, while the corresponding fast-projective bound is $\eta_{\rm fast}+\delta_{\rm XOF}$;
\item with AES modeled as an ideal cipher, the corresponding bounds are $\eta_{\rm local}+\delta_{\rm IC}$ and $\eta_{\rm fast}+\delta_{\rm IC}$.
\end{enumerate}
Every other key for the relevant route satisfies the corresponding key-by-key recovery theorem: accepted outputs are always correct, and verified restart terminates almost surely with the expected-work bounds of \cref{thm:graph-attack,thm:attack}.  At Levels I, III, and V, both ideal-key exceptions are negligible compared with the transfer terms, so the base-two logarithms round respectively to
\[
(-121.245,-184.715,-248.286)
\quad\text{and}\quad
(-99.863,-96.655,-94.112)
\]
for SHAKE and AES.
\end{corollary}

\begin{proof}
By \cref{prop:rom-transfer,prop:icm-transfer}, the probability of every key event changes by at most the corresponding total-variation distance.  Apply this separately to the graph-direction event of \cref{cor:graph-ideal-recovery} and the fast-projective event of \cref{cor:ideal-recovery}.  On either complement, the corresponding theorem is key by key and independent of how the expanded key was sampled.  Exact verification is distribution-free.  The ideal exceptional probabilities are far below the transfer terms at the displayed precision.
\end{proof}

\subsection{Why this does not prove a fixed-primitive theorem}

\begin{proposition}[Support barrier for deterministic seeded expansion]\label{prop:seed-support-barrier}
Let $E:\{0,1\}^{\lambda}\to\Omega$ be any deterministic public expansion map, let $S$ be uniform on $\{0,1\}^{\lambda}$, and let $U$ be uniform on the finite set $\Omega$.  Then
\[
\Delta_{\rm TV}(E(S),U)
\ge 1-\frac{2^{\lambda}}{|\Omega|}.
\]
In particular, if $\Omega=\{0,1\}^B$, the distance is at least $1-2^{\lambda-B}$.
\end{proposition}

\begin{proof}
The distribution $E(S)$ is supported on at most $2^\lambda$ points.  Let $R$ be its support.  Then $\Pr[E(S)\in R]=1$, while $\Pr[U\in R]=|R|/|\Omega|\le2^\lambda/|\Omega|$.  The variational characterization of total variation gives the claim.
\end{proof}

Thus the fixed public SHAKE/AES stream vector cannot be information-theoretically close to independent uniform streams.  Nor is there a generic public-seed PRG reduction: the public seed is revealed, expansion is deterministic, and a distinguisher can recompute the official stream and compare it with the transcript.  The ideal-XOF and ideal-cipher key-density conclusions above are exact statements inside those models; they do not instantiate themselves for fixed primitives.

The key-by-key recovery theorem nevertheless applies to every fixed-primitive key satisfying the key-by-key conditions.  A high-probability standard-model statement must therefore analyze the relevant algebraic bad events under the actual fixed expansion map, or introduce an explicitly event-specific hardness assumption.  The support barrier does not rule out such work; it rules out a generic statistical or transcript-pseudorandomness transfer that replaces the revealed deterministic expansion by independent uniform streams.

\section{Optional identification with the planted oil space}\label{app:uniqueness}

Neither recovery theorem requires uniqueness: it keeps every verified candidate and searches all of them for a signing witness.  This appendix proves the stronger generic identification statement that, for almost every key, no second $K$-rational common isotropic $O$-space exists.  Fix the planted oil space $\cO_0$.  Let $\cO'\ne\cO_0$ be another $O$-space and put
\[
s=O-\dim(\cO'\cap\cO_0),\qquad 1\le s\le O.
\]

\begin{lemma}[Relative Grassmannian stratum]\label{lem:relative-grassmann}
The locus of such $\cO'$ with fixed $s$ has dimension $s(V+O-s)$.
\end{lemma}

\begin{proof}
Choose the $(O-s)$-dimensional intersection inside $\cO_0$, the $s$-dimensional image of $\cO'$ in $K^{V+O}/\cO_0\simeq K^V$, and the graph map back to the $s$-dimensional quotient of $\cO_0$.  The dimensions are $(O-s)s$, $s(V-s)$, and $s^2$, respectively.
\end{proof}

\begin{lemma}[Independent extra isotropy conditions]\label{lem:extra-isotropy}
For fixed $\cO'$ in this stratum, requiring one central form that already vanishes on $\cO_0\times\cO_0$ also to vanish on $\cO'\times\cO'$ imposes exactly
\[
c_s=\binom{O+1}{2}-\binom{O-s+1}{2}=\frac{s(2O-s+1)}2
\]
additional independent linear conditions.
\end{lemma}

\begin{proof}
Restriction to $\cO'$ already vanishes on $\Sym^2(\cO'\cap\cO_0)$.  Conversely, after choosing complements, every symmetric form on $\cO'$ that vanishes on this subspace extends to an ambient symmetric form vanishing on $\cO_0\times\cO_0$.  The restriction map therefore has the displayed rank.
\end{proof}

\begin{theorem}[Uniqueness codimension and rational first moment]\label{thm:uniqueness}
For $m$ independent central forms, the key locus admitting an alternative common isotropic $O$-space of type $s$ has codimension at least
\[
\delta_s=m\frac{s(2O-s+1)}2-s(V+O-s).
\]
The minimum over $s$ equals $903,1927,3539$ for Levels~I, III, and V.  With Gaussian binomial coefficients denoted by $\genfrac{[}{]}{0pt}{}{a}{b}_Q$,
\[
\Pr[\text{a second $K$-rational oil space}]
\le
\sum_{s=1}^{O}
\genfrac{[}{]}{0pt}{}{O}{s}_Q
\genfrac{[}{]}{0pt}{}{V}{s}_Q
Q^{s^2-mc_s}.
\]
Writing the displayed first-moment bound as $\eta_{\rm uniq}$, we have
\[
\log_2\eta_{\rm uniq}<-18932.346,\quad -40401.586,\quad -74198.865
\]
at Levels I, III, and V, respectively.
\end{theorem}

\begin{proof}
Over the stratum of dimension $s(V+O-s)$, each of the $m$ form fibers loses $c_s$ dimensions by \cref{lem:extra-isotropy}.  The incidence variety therefore has codimension at least $mc_s-s(V+O-s)$ after projection to key space.  The function $\delta_s$ is concave in $s$, so its minimum on $1\le s\le O$ occurs at an endpoint; direct evaluation gives the displayed values and shows that $s=1$ minimizes in all three parameter sets.

For the finite-field statement, the number of $K$-rational $O$-spaces in the type-$s$ stratum is
\[
\genfrac{[}{]}{0pt}{}{O}{s}_Q
\genfrac{[}{]}{0pt}{}{V}{s}_Q Q^{s^2}.
\]
For each fixed $\cO'$, the $m$ independent forms satisfy the $c_s$ extra conditions with probability $Q^{-mc_s}$.  The union bound gives the formula; exact evaluation gives the stated logarithms.
\end{proof}

\begin{corollary}[Generic identification with the planted decomposition]\label{cor:planted-identification}
In the public-random model, outside the exceptional event of \cref{cor:graph-ideal-recovery} and the additional event of probability $\eta_{\rm uniq}$ in \cref{thm:uniqueness}, every exact graph-direction candidate has the planted oil space.  The same conclusion holds for the fast projective route outside the exceptional event of \cref{cor:ideal-recovery} and the uniqueness event.  In either case, the returned decomposition is the planted expanded decomposition up to the verified coordinate representation.  The secret seed is still not recovered or needed.
\end{corollary}

\begin{proof}
Every exact candidate defines a common isotropic $O$-space.  Uniqueness forces that space to be $\cO_0$, and its graph determines the expanded decomposition used by \cref{prop:verified-key}.
\end{proof}

\section{Alternative elimination and low-space directions}\label{app:alternatives}

The main text uses complete univariate resolutions because they are the only route for which we have both a suitable theorem and full failure accounting.  This appendix records two exact ways to extract a root more compactly once a quotient algebra exists, and then explains why several natural attempts do not yet improve the proved exponent.

\subsection{Extracting an isolated root from an existing quotient algebra}\label{subsec:implicit-extraction}

Let $A=L[y_1,\ldots,y_V]/I$ be the reduced affine coordinate algebra of a square core, of degree $d\le D$.  Write $Z$ for its geometric root set, and let $q_1,\ldots,q_m$ denote the original restricted equations in $A$.

\begin{proposition}[Norm-gradient root formula]\label{prop:implicit-root}
Suppose $p\in L^V$ is the unique point of $Z$ satisfying all $m$ original equations.  Choose $\rho\leftarrow L^m$, put $h=\sum_i\rho_iq_i\in A$, and define
\begin{equation}\label{eq:norm-jet}
\mathcal N_h(z_0,\ldots,z_V)
:=\operatorname{Norm}_{A/L}\!\left(h+z_0+\sum_{j=1}^Vz_jy_j\right).
\end{equation}
With probability at least $1-(d-1)/|L|$, $h$ vanishes at $p$ and at no other geometric core point.  On that event,
\begin{equation}\label{eq:norm-gradient-root}
\partial_{z_0}\mathcal N_h(0)\ne0,
\qquad
p_j=
\frac{\partial_{z_j}\mathcal N_h(0)}
     {\partial_{z_0}\mathcal N_h(0)}
\quad(1\le j\le V).
\end{equation}
Equivalently, if $M_h$ and $M_{y_j}$ are the multiplication operators in $A$, then $\ker M_h$ is one-dimensional and every nonzero $w\in\ker M_h$ satisfies $M_{y_j}w=p_jw$.
\end{proposition}

\begin{proof}
After scalar extension to an algebraic closure, reducedness gives
\[
A\otimes_L\overline L\cong\prod_{x\in Z}\overline L.
\]
At every $x\ne p$, the vector $(q_1(x),\ldots,q_m(x))$ is nonzero.  The condition $h(x)=0$ cuts out a proper $L$-linear subspace of $L^m$, so a uniform $\rho$ satisfies it with probability at most $1/|L|$.  A union bound over at most $d-1$ points proves the selector bound.  On the good event, the norm factors as
\[
\mathcal N_h(z)=
\prod_{x\in Z}\left(h(x)+z_0+\sum_{j=1}^Vz_jx_j\right).
\]
Differentiating at the origin gives
\[
\partial_{z_0}\mathcal N_h(0)=\prod_{x\ne p}h(x),
\qquad
\partial_{z_j}\mathcal N_h(0)=p_j\prod_{x\ne p}h(x),
\]
which proves \eqref{eq:norm-gradient-root}.  In the same product decomposition, $M_h$ is diagonal with a unique zero eigenvalue and the coordinate operators act on its kernel by the scalars $p_j$.
\end{proof}

\begin{remark}[Several surviving roots]\label{rem:norm-multiple}
For an arbitrary selector $h$, let $S=\{x\in Z:h(x)=0\}$.  The lowest nonzero homogeneous part of \eqref{eq:norm-jet} at the origin is
\[
\left(\prod_{x\notin S}h(x)\right)
\prod_{x\in S}\left(z_0+\sum_{j=1}^Vz_jx_j\right).
\]
Thus the first nonzero norm jet is the Chow product of all selected roots.  The linear formula \eqref{eq:norm-gradient-root} is exactly the case $|S|=1$.
\end{remark}

\begin{proposition}[Resultant-gradient root formula]\label{prop:resultant-gradient}
Assume $\operatorname{char}L\ne2$.  Let $F_0,\ldots,F_V\in L[t_0,\ldots,t_V]$ be homogeneous quadrics with a unique geometric common point $p=[p_0:\cdots:p_V]$.  Assume that their value Jacobian on the tangent space $T_p\PP^V$ has rank $V$.  Let $\mathcal R$ be their multivariate resultant, written in the monomial coefficients $c_{i,\alpha}$ of $F_i$, $|\alpha|=2$.  Then the coefficient tuple is a smooth point of the resultant hypersurface.  If $\lambda=(\lambda_0,\ldots,\lambda_V)$ spans the left kernel of the tangent Jacobian, there is a scalar $\kappa\ne0$ such that
\begin{equation}\label{eq:resultant-gradient}
\frac{\partial\mathcal R}{\partial c_{i,\alpha}}
=\kappa\lambda_i p^\alpha
\qquad(0\le i\le V,\ |\alpha|=2).
\end{equation}
In particular, for any $i,k$ with $\lambda_i p_k\ne0$,
\begin{equation}\label{eq:resultant-ratio}
\frac{p_j}{p_k}
=
\frac{\partial\mathcal R/\partial c_{i,e_j+e_k}}
     {\partial\mathcal R/\partial c_{i,2e_k}}
\qquad(0\le j\le V).
\end{equation}
The coefficient convention in \eqref{eq:resultant-ratio} treats $t_jt_k$ as one monomial, with no symmetric-matrix factor of two.
\end{proposition}

\begin{proof}
Consider the incidence variety of coefficient tuples and projective points satisfying $F_i(p)=0$ for every $i$.  Its first-order equations at the given pair are
\[
\delta F_i(p)+dF_i(p)[\delta p]=0\qquad(0\le i\le V).
\]
The tangent Jacobian has rank $V$, so a tangent displacement $\delta p$ exists exactly when
\[
\sum_{i=0}^V\lambda_i\,\delta F_i(p)=0.
\]
Uniqueness and transversality make the incidence projection locally one-to-one onto a smooth branch of the resultant hypersurface.  Its tangent hyperplane is therefore the kernel of the displayed nonzero functional.  The differential of the irreducible resultant is proportional to that functional, giving \eqref{eq:resultant-gradient}; the ratio formula follows by evaluating the two quadratic monomials at $p$.
\end{proof}

\begin{remark}[Compression achieved, evaluator still missing]\label{rem:implicit-memory}
The norm gradient outputs one root using only $V+1$ field elements, and the resultant gradient uses one nonzero coefficient block.  If quotient multiplication is available as a black box, a one-dimensional kernel can be found with $O(d)$ stored field elements and $O(d)$ multiplication-operator products by standard black-box linear algebra~\cite{Wiedemann}.  Alternatively, Baur--Strassen converts any arithmetic circuit for the norm or resultant into a circuit for all first derivatives with constant-factor overhead~\cite{BaurStrassen}.  These statements compress the \emph{root-extraction interface}; they do not construct the degree-$d$ quotient algebra, multiplication operators, norm, or resultant in $o(d)$ space.  With current constructions, the time remains near-quadratic in $d$ (with the fixed $\varepsilon$ loss of \cref{prop:finite-field-kronecker}) and at least one degree-$d$ state vector remains.  Moreover, \cref{prop:generic-full-isolation} gives no production-parameter density bound for uniqueness of the full-system root.  The formulas are therefore exact algorithmic targets, not a practical attack.
\end{remark}

\subsection{Why the spare equations do not yet lower the finite exponent}\label{subsec:exponent-audit}

The full restricted system has $e=m-V\in\{2,2,3\}$ equations beyond the square core, and generically those equations isolate the planted point.  We tested whether they also lower the solving exponent.  The following deliberately optimistic semi-regular screen uses the Hilbert series
\[
\frac{(1-t^2)^{V+e}}{(1-t)^{V+1}}
\]
for $V+e$ homogeneous quadrics in $V+1$ projective variables.  Because the actual system contains a planted point, this is a diagnostic rather than a theorem.  At the last positive coefficient, the monomial counts are:
\begin{center}
\small
\begin{tabular}{@{}c r r r r r@{}}
\toprule
Level & $V$ & degree & $\log_2 N$ & optimistic $\log_2N^2$ & current $\log_2D^2$\\
\midrule
I   &52  &27 &$69.787$  &$139.574$ &$104$\\
III &76  &39 &$102.585$ &$205.170$ &$152$\\
V   &102 &48 &$131.814$ &$263.629$ &$204$\\
\bottomrule
\end{tabular}
\end{center}
Even charging only $N^2$ sparse black-box work is worse than the baseline $D^2$ geometric scale and also worse than the primary $\varepsilon=0.3$ fast-line indicators after field and circuit conversion; dense linear algebra is worse still.  This does not rule out a coefficient-specific method, but it rules out the most direct overdetermined-Macaulay improvement under standard linear algebra.

\begin{proposition}[Underslicing loses more in hit probability than it saves in degree]\label{prop:underslice}
Let $0\le d<V$, and choose a uniform projective $d$-plane over $K=\F_Q$.  Its probability of meeting the planted projective oil space is
\begin{equation}\label{eq:underslice-probability}
\pi_d
=1-\prod_{i=0}^{d}
\frac{1-Q^{i-V}}{1-Q^{i-(V+O)}}.
\end{equation}
If the same degree-squared square-core solver is run independently on every sampled plane, its expected leading degree factor is $4^d/\pi_d$.  For $d=V-r$ and the \QRUOV{} base field $Q=127^3$, removing $r=1,2,3$ slice dimensions increases the base-two exponent relative to the $d=V$ slice by
\[
18.966053,\qquad37.932107,\qquad56.898161
\]
bits, respectively, before any other work is charged.  Choosing the slices over the larger solver field only worsens this tradeoff.
\end{proposition}

\begin{proof}
Write $k=d+1$.  A vector $k$-subspace disjoint from the fixed $O$-dimensional oil space is the graph of a linear map from a $k$-subspace of the $V$-dimensional quotient to the oil space.  Hence there are
\[
Q^{Ok}\genfrac{[}{]}{0pt}{}{V}{k}_Q
\]
disjoint subspaces, out of $\genfrac{[}{]}{0pt}{}{V+O}{k}_Q$ total.  Taking the ratio and expanding the Gaussian binomial coefficients gives \eqref{eq:underslice-probability}.  Independent repetition needs expected $1/\pi_d$ planes.  Substitution into $4^d/\pi_d$ gives the displayed values, which are recomputed by the companion numerical audit.  Asymptotically each removed dimension costs $\log_2Q-2=18.966054\ldots$ bits: it saves two degree-squared bits but loses roughly one factor of $Q$ in intersection probability.
\end{proof}

The same calculation covers the equivalent hybrid strategy of guessing one base-field coordinate of the oil point before solving: every guessed coordinate costs roughly $\log_2Q$ bits and saves only two bits in the degree-squared solver.  Thus the minimal guaranteed slice is also the best random underslice under the current solver model.

The other natural routes also stop short of an exponent improvement.  Explicit resultant matrices remain exponentially larger than the $D$-dimensional quotient representation in the dense quadratic regime~\cite{DAndreaDickenstein}.  Chow-form algorithms are polynomial in geometric degree but do not provide a subquadratic dependence specialized to these parameters~\cite{JeronimoEtAl}.  Faster modular composition lowers an operation used after a univariate algebra exists, not the geometric-intersection stage that constructs that algebra~\cite{NeigerEtAl}.

We also checked two recent refinements of Macaulay/Gr\"obner solving.  Degree smoothing replaces a full Macaulay matrix by a sufficient submatrix lying between two adjacent integer degrees~\cite{DegreeSmoothing}.  Our optimistic screen already charges the complete monomial space at the \emph{lower} adjacent degree and is still $35.574$, $53.170$, and $59.629$ bits above the $D^2$ geometric scale; smoothing within that interval cannot reverse the comparison.  The recent F4-tracer analysis for generic square degree-$\delta$ systems has leading arithmetic scale $\widetilde O(\delta^{\omega V+1}e^{cV})$ and relies on trace compatibility and genericity conjectures~\cite{F4Tracer}.  For quadrics, even its factor $2^{\omega V}$ exceeds $D^2=2^{2V}$ for every currently proved $\omega>2$.  Thus these improvements are valuable in their intended regimes but do not lower the exponent here.  The norm and resultant gradients identify a smaller \emph{output} target; no proved route found in this audit computes that target below the published near-quadratic-in-$D$ line used here.

\Needspace{10\baselineskip}
\section{Reproducibility and executable attack description}\label{app:algorithms}

\subsection{Companion executable audits}

The bundle contains four standard-library Python programs.
\begin{itemize}
\item \path{QRUOV_scissors_audit.py} recomputes the target-orthogonal dimension
inequalities, projective $m-4$ cores, fixed-key seed--oil intersection bounds,
second moments, eligible-root probabilities, random-chart and separator
factors, both finite-size cost rows, memory envelopes, Karatsuba plans,
sparse-field certificates, and attack-specific retuning diagnostics.  It emits
text and JSON and freezes every displayed number.
\item \path{QRUOV_directional_audit.py} checks the corrected $2C_0$
reconstruction precision, Vandermonde product starts, characteristic-$127$
lifting identities, flattened multiplication, sparse-field certificates,
recursive Karatsuba convolution, and complete graph recovery on $200$ toy keys.
\item \path{QRUOV_direct_qr_toy.py} generates exact degree-three extension-field
UOV instances, applies the trace pull-back, and executes full target
normalization, seed search, Chevalley--Warning extension, projective residual
search, affine rescaling, and public verification end to end.
\item \path{QRUOV_cw_helper.py} supplies the finite-field and
Chevalley--Warning helper routines used by the direct experiment.
\end{itemize}
These programs audit formulas and attack logic.  They do not instantiate the
astronomically large full-parameter symbolic state.

\subsection{Executable attack description}

\paragraph{Projective target-orthogonal direct-forgery route.}
\begin{enumerate}
\item Derive the official nonzero verification target $y$ from the message and
salt.  Choose a public invertible output transformation sending $y$ to
$(0,\ldots,0,1)$, and apply the same transformation to the public quadrics.
\item Sample a uniform seven-dimensional seed subspace $S_0\subset\F_q^n$.
Enumerate $\mathbb P(S_0)$ until finding a nonzero common zero of the first
three transformed quadrics; Chevalley--Warning guarantees one.
\item Starting from its span, repeatedly intersect the three polar-orthogonal
spaces and solve the induced three quadrics on a seven-dimensional quotient
slice.  Extend until obtaining a common totally isotropic subspace $U$ of
dimension $m-3$.
\item Choose a random nonzero chart form on $U$.  Restrict the next $m-4$
zero-target quadrics to the chart and run the checked finite-characteristic
nonsingular-root solver on the resulting $m-4$ variable square system.
\item At every represented rational projective root, evaluate the final
transformed quadratic.  If its value is a nonzero square, recover either affine
scale making that value one, reverse the transformations, and re-evaluate all
original public equations.
\item Return the first vector passing the unchanged QR-UOV verifier.  If the
chart, solver invocation, or residual search fails, change the salt and attack
coins and repeat.  Exact verification gives one-sided error.
\end{enumerate}

\paragraph{Primary finite-characteristic direction route.}
\begin{enumerate}
\item Expand the public key to matrices $P_i$ over $K=\F_{127^3}$.  Choose a nonzero public direction $c\in K^O$ and form $q_{i,c}(x)=P_i(-x,c)$.
\item Choose a random $V\times m$ output-recombination matrix $R$ and form the square system $f=Rq_c$.  Work in $K'=\F_{127^d}$ using the sparse degree-$18$, $24$, or $30$ model and the audited recursive Karatsuba circuit for the relevant parameter level.
\item Run the finite-characteristic nonsingular-root algorithm of \cref{thm:finite-char-nonsingular}, using the Vandermonde product start and a full random separator.  A bad invocation may omit roots or fail; it cannot create an accepted false key.
\item Re-evaluate all $m$ original directional equations at every represented $K$-rational root.  At each survivor, solve the $O$ public linear systems of \cref{prop:graph-direction-completion}.  Keep only graphs passing every graph identity, total-isotropy check, inverse pull-back, and exact public-map factorization.
\item Run the basis-first signing scan and, if needed, the pair-line fallback.  Return the first exact candidate with an invertible oil matrix.  If no candidate passes, sample fresh solver coins and then a fresh outer direction.
\end{enumerate}
Coordinate directions and Wronskian condensers provide a deterministic outer fallback; the primary finite-size ledger uses the simpler random full-rank recombination analyzed in \cref{thm:verified-directional-recovery}.

\paragraph{Fast projective route.}
\begin{enumerate}
\item Set $\varepsilon=0.3$.  Use the minimum fields $s=8,12,15$ for expected-restart analysis or the safety fields $s=9,13,16$ for the one-invocation theorem.  Sample a full-rank slice matrix $M$ and an output recombination $R$ over $L_s$; restrict the public quadrics and form the square core.
\item Run one finite checked invocation of the van der Hoeven--Lecerf solver.  Return \textsf{fail} on an exceptional division, degree mismatch, or consistency failure.  Do not hide an unbounded internal restart inside one outer draw.
\item Canonicalize and verify the resolution.  Evaluate all original restricted equations in its quotient algebra, prefilter with \cref{prop:separator-prefilter}, extract $L_s$-rational roots, and recheck all equations at each reconstructed point.
\item At every survivor, compute the common derivative kernel.  Require dimension $O$, descent to a graph over $K$, total isotropy, all graph identities, the official inverse pull-back, and exact factorization of the original public map.
\item Run the basis-first signing scan and, if needed, the pair-line fallback.  Return the first verified signing certificate; otherwise begin a fresh outer trial.
\end{enumerate}

Both wrappers have one-sided error: a solver failure may cause a retry, but exact public verification prevents an invalid equivalent key or invalid signature from being returned.  The output is a verified expanded signing decomposition and an inverse oil matrix, not the original seed or an official serialized private key.


\begin{multicols}{2}
\raggedcolumns
\begin{thebibliography}{99}
\fontsize{8.5pt}{9.6pt}\selectfont
\setlength{\itemsep}{0pt}
\setlength{\parsep}{0pt}

\bibitem{QRUOVSpec}
H.~Furue, Y.~Ikematsu, F.~Hoshino, T.~Takagi, H.~Kosuge, K.~Yamakoshi, R.~Akiyama, S.~Nakamura, S.~Orihara, and K.~Kinjo.
\newblock \emph{QR-UOV Specification Document (Round 2 submission)}.
\newblock NIST Additional Digital Signatures Project, 5 February 2025.

\bibitem{KPG1999}
A.~Kipnis, J.~Patarin, and L.~Goubin.
\newblock Unbalanced Oil and Vinegar signature schemes.
\newblock In \emph{EUROCRYPT 1999}, pp.~206--222, 1999.

\bibitem{NISTRound3}
National Institute of Standards and Technology.
\newblock Nine candidates advance to the third round of the Additional Digital Signatures for the PQC standardization process.
\newblock NIST CSRC, 14 May 2026.

\bibitem{NISTIR8610}
G.~Alagic et al.
\newblock Status report on the second round of the Additional Digital Signature Schemes for the NIST Post-Quantum Cryptography Standardization Process.
\newblock NIST IR~8610, May 2026.

\bibitem{Benoist}
O.~Benoist.
\newblock Degr\'es d'homog\'en\'eit\'e de l'ensemble des intersections compl\`etes singuli\`eres.
\newblock \emph{Annales de l'Institut Fourier}, 62(3):1189--1214, 2012.

\bibitem{vdHL}
J.~van der Hoeven and G.~Lecerf.
\newblock On the complexity exponent of polynomial system solving.
\newblock \emph{Foundations of Computational Mathematics}, 21(1):1--57, 2021.

\bibitem{GiustiLecerfSalvy2001}
M.~Giusti, G.~Lecerf, and B.~Salvy.
\newblock A Gr\"obner free alternative for polynomial system solving.
\newblock \emph{Journal of Complexity}, 17(1):154--211, 2001.
\newblock doi:10.1006/jcom.2000.0571.

\bibitem{Gimenez2026}
N.~Gim\'enez, J.~Heintz, G.~Matera, L.~M.~Pardo, M.~P\'erez, and M.~Privitelli.
\newblock A Kronecker algorithm for locally closed sets over a perfect field.
\newblock arXiv:2512.14888v2, 5 June 2026.

\bibitem{Beullens}
W.~Beullens.
\newblock Improved cryptanalysis of UOV and Rainbow.
\newblock In \emph{EUROCRYPT 2021, Part I}, pp.~348--373, 2021.

\bibitem{FurueIkematsu}
H.~Furue and Y.~Ikematsu.
\newblock A new security analysis against MAYO and QR-UOV using rectangular MinRank attack.
\newblock In \emph{IWSEC 2023}, pp.~101--116, 2023.

\bibitem{Suzuki}
T.~Suzuki, H.~Furue, T.~Ito, S.~Nakamura, and S.~Uchiyama.
\newblock An extended rectangular MinRank attack against UOV and its variants.
\newblock In \emph{IWSEC 2025}, LNCS~16208, pp.~111--130, Springer, 2026.

\bibitem{FurueIkematsu2026}
H.~Furue and Y.~Ikematsu.
\newblock Key recovery attacks on UOV using $p^\ell$-truncated polynomial rings.
\newblock Cryptology ePrint Archive, Paper 2026/298, 2026.

\bibitem{JustGuess}
A.~May, M.~Ostuzzi, and H.~Ressler.
\newblock Just Guess: Improved (quantum) algorithm for the underdetermined MQ problem.
\newblock In \emph{EUROCRYPT 2026, Part IV}, LNCS~16544, pp.~334--362, 2026.

\bibitem{Pebereau}
P.~P\'ebereau.
\newblock One vector to rule them all: Key recovery from one vector in UOV schemes.
\newblock In \emph{PQCrypto 2024}, pp.~92--108, 2024; ePrint 2023/1131.


\bibitem{BaurStrassen}
W.~Baur and V.~Strassen.
\newblock The complexity of partial derivatives.
\newblock \emph{Theoretical Computer Science}, 22(3):317--330, 1983.
\newblock doi:10.1016/0304-3975(83)90110-X.

\bibitem{DAndreaDickenstein}
C.~D'Andrea and A.~Dickenstein.
\newblock Explicit formulas for the multivariate resultant.
\newblock \emph{Journal of Pure and Applied Algebra}, 164(1--2):59--86, 2001.
\newblock doi:10.1016/S0022-4049(00)00145-6.

\bibitem{JeronimoEtAl}
G.~Jeronimo, T.~Krick, J.~Sabia, and M.~Sombra.
\newblock The computational complexity of the Chow form.
\newblock \emph{Foundations of Computational Mathematics}, 4(1):41--117, 2004.
\newblock doi:10.1007/s10208-002-0078-2.

\bibitem{NeigerEtAl}
V.~Neiger, B.~Salvy, E.~Schost, and G.~Villard.
\newblock Faster modular composition using two relation matrices.
\newblock In \emph{ISSAC 2026}, pp.~315--324, 2026; arXiv:2601.17422.

\bibitem{DegreeSmoothing}
S.~Jaques, L.~Ran, S.~Samardjiska, and M.~Seitner.
\newblock Smoothing the degree of regularity for polynomial systems.
\newblock Cryptology ePrint Archive, Paper 2026/408, 2026.

\bibitem{F4Tracer}
R.~Kouba, V.~Neiger, and M.~Safey El Din.
\newblock A complexity analysis of the F4 Gr\"obner basis algorithm with tracer data.
\newblock arXiv:2603.16378, 2026.

\bibitem{KaltofenShoup}
E.~Kaltofen and V.~Shoup.
\newblock Subquadratic-time factoring of polynomials over finite fields.
\newblock \emph{Mathematics of Computation}, 67(223):1179--1197, 1998.

\bibitem{Wiedemann}
D.~Wiedemann.
\newblock Solving sparse linear equations over finite fields.
\newblock \emph{IEEE Transactions on Information Theory}, 32(1):54--62, 1986.

\bibitem{Heintz}
J.~Heintz.
\newblock Definability and fast quantifier elimination in algebraically closed fields.
\newblock \emph{Theoretical Computer Science}, 24:239--277, 1983.

\bibitem{LachaudRolland}
G.~Lachaud and R.~Rolland.
\newblock On the number of points of algebraic sets over finite fields.
\newblock \emph{Journal of Pure and Applied Algebra}, 219(11):5117--5136, 2015.


\bibitem{Rabin1980}
M.~O.~Rabin.
\newblock Probabilistic algorithms in finite fields.
\newblock \emph{SIAM Journal on Computing}, 9(2):273--280, 1980.
\newblock doi:10.1137/0209024.

\bibitem{SafeySchost2018}
M.~Safey El Din and E.~Schost.
\newblock Bit complexity for multi-homogeneous polynomial system solving: Application to polynomial minimization.
\newblock arXiv:1605.07433v2, 2017; manuscript revised 6 November 2018.

\bibitem{Schost2003}
E.~Schost.
\newblock Computing parametric geometric resolutions.
\newblock \emph{Applicable Algebra in Engineering, Communication and Computing}, 13(5), 2002.
\newblock doi:10.1007/s00200-002-0109-x.

\bibitem{QRUOVResponse2026}
S.~Orihara, on behalf of the QR-UOV team.
\newblock Round 2 (Additional Signatures) official comment: QR-UOV.
\newblock NIST PQC Forum, 20 January 2026.

\bibitem{Hashimoto2023}
Y.~Hashimoto.
\newblock An improvement of algorithms to solve under-defined systems of multivariate quadratic equations.
\newblock \emph{JSIAM Letters}, 15:53--56, 2023.
\newblock doi:10.14495/jsiaml.15.53.

\bibitem{ThomaeWolf2012}
E.~Thomae and C.~A.~Wolf.
\newblock Solving underdetermined systems of multivariate quadratic equations revisited.
\newblock In \emph{PKC 2012}, LNCS~7293, pp.~156--171, 2012.
\newblock doi:10.1007/978-3-642-30057-8\_10.

\bibitem{FurueUnderdetermined2021}
H.~Furue, S.~Nakamura, and T.~Takagi.
\newblock Improving Thomae--Wolf algorithm for solving underdetermined multivariate quadratic polynomial problem.
\newblock In \emph{PQCrypto 2021}, LNCS~12841, pp.~65--78, 2021.
\newblock doi:10.1007/978-3-030-81293-5\_4.

\bibitem{Chevalley1935}
C.~Chevalley.
\newblock D\'emonstration d'une hypoth\`ese de M.~Artin.
\newblock \emph{Abhandlungen aus dem Mathematischen Seminar der Universit\"at Hamburg}, 11:73--75, 1935.

\bibitem{Warning1935}
E.~Warning.
\newblock Bemerkung zur vorstehenden Arbeit von Herrn Chevalley.
\newblock \emph{Abhandlungen aus dem Mathematischen Seminar der Universit\"at Hamburg}, 11:76--83, 1935.

\bibitem{JinEtAl2026}
Y.~Jin, Y.~Pan, X.~He, B.~Gong, and J.~Ding.
\newblock Security analysis on UOV families with odd characteristics: Using symmetric algebra.
\newblock Cryptology ePrint Archive, Paper 2025/1137; in \emph{PKC 2026, Part I}, LNCS~16551, pp.~428--454, 2026.

\end{thebibliography}
\end{multicols}

\end{document}
